\documentclass[11pt]{article}
\usepackage[T1]{fontenc}
\usepackage{lmodern}
\usepackage{amsmath,amssymb,amsthm,mathtools,mathrsfs}
\usepackage[margin=1in]{geometry}
\usepackage{booktabs,array,longtable,enumitem}
\usepackage{microtype}
\usepackage[colorlinks=true,linkcolor=blue,citecolor=blue,urlcolor=blue]{hyperref}
\hypersetup{pdftitle={Uniform Laguerre signs and hypercomplex jets of the theta heat flow},
 pdfauthor={Drafted for Edward Baker},pdfsubject={Four-jet positivity, quaternionic and octonionic geometry, and the remaining inner region}}
\allowdisplaybreaks
\setlength{\emergencystretch}{2em}
\numberwithin{equation}{section}
\newtheorem{theorem}{Theorem}[section]
\newtheorem{proposition}[theorem]{Proposition}
\newtheorem{lemma}[theorem]{Lemma}
\newtheorem{corollary}[theorem]{Corollary}
\newtheorem{classicalinput}[theorem]{Classical input}
\theoremstyle{definition}
\newtheorem{definition}[theorem]{Definition}
\theoremstyle{remark}
\newtheorem{remark}[theorem]{Remark}
\newcommand{\R}{\mathbb R}
\newcommand{\C}{\mathbb C}
\newcommand{\HH}{\mathbb H}
\newcommand{\OO}{\mathbb O}
\newcommand{\ii}{\mathbf i}
\newcommand{\jj}{\mathbf j}
\newcommand{\kk}{\mathbf k}
\newcommand{\dd}{\,d}
\newcommand{\Lag}{\mathcal L}
\newcommand{\RH}{\mathrm{RH}}
\newcommand{\Ephys}{E}
\DeclareMathOperator{\Rea}{Re}
\DeclareMathOperator{\Ima}{Im}
\title{Uniform Laguerre signs and hypercomplex jets\\of the theta heat flow}
\author{Drafted for Edward Baker}
\date{15 October 2026}
\begin{document}
\maketitle
\begin{abstract}
For the theta heat family $H$, the two forms
$H_x^2-HH_{xx}$ and $H_{xx}^2-H_xH_{xxx}$ control the evolution of its
first jet at genuine stationary points. We prove both forms strictly
positive, uniformly for $0<t\le1/2$ and sufficiently large $|x|$, when
$t\log(|x|/(4\pi))\ge61/24$, including the closed boundary.
Five exponent pairs and Polymath's effective Riemann--Siegel approximation
give a common nonzero amplitude for all four derivatives; shrinking
complex discs pay the derivative errors. Exact staircase cutoffs satisfy
energy-relative and complete Gaussian-tail lower bounds. A conditional
middle-block estimate would extend the signs to the closed strip
$5/2\le t\log(|x|/(4\pi))\le8/3$; that hypothesis remains unproved.
Complete-source certificates establish both stationary signs and a
physical first-jet gap through $|x|\le112$. Quaternionic and octonionic
fields preserve the signed observations. Delay obstructions, a
phase-controlled outer comparison, and all-earlier-time finite
comparisons distinguish valid correlations from unsupported signs.
Collision expansions, analytic rigidity of logarithmic losses, and
conditional cutoff transfers isolate a signed estimate near small
physical energy as a remaining route to collision exclusion. A weaker
weighted-remainder hypothesis also suffices for eventual signs. The
inner cusp and the finite transition to the asymptotic region remain open.
\end{abstract}
\tableofcontents

\section{The theta heat flow and the main result}\label{sec:introduction}

The four real derivatives
\begin{equation}
 h_j(t,x)=\partial_x^jH(t,x),\qquad j=0,1,2,3,
 \label{eq:four-jets}
\end{equation}
form the smallest consecutive jet containing the collision detector
$(H,H_x)$ and both quadratic expressions
\begin{equation}
 \Lag_0=h_1^2-h_0h_2,\qquad
 \Lag_1=h_2^2-h_1h_3.
 \label{eq:laguerre-definitions}
\end{equation}
Here ``four jets'' means these four consecutive derivative components.
Their common scalar origin is essential: they are not four independent
fields. The first two give the quaternionic observation
$F=h_0+\jj h_1$ and its energy $E=h_0^2+h_1^2$. Since $H_t=-H_{xx}$,
\begin{equation}
 \frac12E_t=-h_0h_2-h_1h_3
 =\begin{cases}\Lag_0,&h_1=0,\\ \Lag_1,&h_2=0.\end{cases}
 \label{eq:stationary-energy-intro}
\end{equation}
The stationary sets in this identity are those of the original scalar
function and its first derivative. Their signs will remain the same
objects throughout the analytic and hypercomplex constructions.

\subsection{The fixed source and normalization}
The family under consideration is
\begin{align}
 \Phi(u)&=\sum_{n\ge1}\phi_n(u),\qquad
 \phi_n(u)=(2\pi^2n^4e^{9u}-3\pi n^2e^{5u})e^{-\pi n^2e^{4u}},
 \label{eq:theta-source}\\
 H(t,z)&=\int_0^\infty e^{tu^2}\Phi(u)\cos(zu)\dd u,
 \qquad H(0,z)=\frac18\xi\left(\frac12+\frac{iz}{2}\right).
 \label{eq:theta-family}
\end{align}
We use the convention
$\xi(s)=s(s-1)\pi^{-s/2}\Gamma(s/2)\zeta(s)/2$, agreeing with
\cite{polymath}. In particular the critical line corresponds to real
$z$, and the heat equation has the backward sign
\begin{equation}
 \partial_tH=-\partial_z^2H.
 \label{eq:backward-heat}
\end{equation}

For completeness, the normalization follows directly from the theta
Mellin identity. Put $A(u)=e^u\sum_{n\ge1}e^{-\pi n^2e^{4u}}$.
Poisson summation makes $P(u)=e^u+2A(u)$ even, and differentiation gives
$\Phi=(P''-P)/16=(A''-A)/8$. Splitting the theta Mellin integral at one
and substituting $y=e^{4u}$ gives
\[
 \xi(1/2+iz/2)=\frac12-(1+z^2)
                    \int_0^\infty A(u)\cos(zu)\dd u.
\]
Evenness of $P$ gives $A'(0)=-1/2$. Two integrations by parts now prove
\eqref{eq:theta-family}; all positive-infinity boundary terms vanish.
This also proves that the complete source $\Phi$ is even. For $u\ge0$
its terms are positive, and their successive majorant ratios yield
\begin{equation}
 0<\Phi(u)<4\pi^2e^{9u-\pi e^{4u}}.
 \label{eq:source-majorant}
\end{equation}
Indeed the ratio of successive terms of $n^4e^{-\pi n^2e^{4u}}$ is at
most $16e^{-3\pi}<1/2$. The bound justifies differentiation of every
order on compact sets of complex $(t,z)$. Thus $H$ is jointly entire,
even in $z$, and real on real $(t,z)$; \eqref{eq:backward-heat} follows
under the integral. For each fixed real time it is a real entire
function of order at most one, as is also immediate from the same
source bound by splitting the integral at $u=O(\log(2+|z|))$.

\subsection{Uniform positivity outside the inner region}
The leading result concerns the dimensionless parameter $t\log|x|$,
which remains nontrivial as time tends to zero while height increases.
\begin{theorem}[Uniform four-jet asymptotic]\label{thm:outer}
There is a finite, presently unevaluated $X_*>4\pi$ with the
following property. Uniformly for
\begin{equation}
 x\ge X_*,\qquad 0<t\le\tfrac12,\qquad
 \theta:=tL\ge61/24,\qquad L:=\log\frac{x}{4\pi},
 \label{eq:outer-domain}
\end{equation}
there is an explicitly defined nonzero complex amplitude $Q=Q_t(x)$
such that, with $\omega=L/4$ and $\phi=\arg Q$,
\begin{equation}
 \frac{1}{2|Q|}\left(h_0,\frac{h_1}{\omega},
                   \frac{h_2}{\omega^2},\frac{h_3}{\omega^3}\right)
 = (\cos\phi,-\sin\phi,-\cos\phi,\sin\phi)
       +O(L^{-1}).
 \label{eq:main-four-jet}
\end{equation}
Consequently
\begin{align}
 h_0^2+(h_1/\omega)^2&=4|Q|^2(1+O(L^{-1})),
 \label{eq:main-balanced}\\
 \Lag_0&=4|Q|^2\omega^2(1+O(L^{-1})),\qquad
 \Lag_1=4|Q|^2\omega^4(1+O(L^{-1})).
 \label{eq:main-signs}
\end{align}
After increasing $X_*$, both Laguerre forms are strictly
positive everywhere in \eqref{eq:outer-domain}, and $F(t,x)\ne0$ there.
Evenness gives the reflected statements at negative heights.
\end{theorem}

The theorem combines the effective approximation in
\cite[Theorem 1.3]{polymath} with five exponent-pair estimates, uniformly
in all four derivatives. Strict rational margins absorb the positive
arithmetic losses and the complex-disc crossing, so the lower boundary
$61/24$ is included exactly. The starting height remains unevaluated;
no optimality claim is made. Lemma~\ref{lem:middle-block-transfer} gives
a conditional continuation to $5/2$, while
Proposition~\ref{pg:weighted-transfer} records a weaker sufficient
weighted-remainder hypothesis. Neither hypothesis is proved here.
A different, more conservative region has a fully explicit height.
\begin{proposition}[An explicit outer region]\label{prop:explicit-outer}
Set $X_{\mathrm{exp}}=4\pi\exp(10^9)$. For $x\ge X_{\mathrm{exp}}$,
$0<t\le1/2$, and $tL\ge5$, every coordinate error in
\eqref{eq:main-four-jet} is at most $108000000/L\le0.108<1/8$.
Consequently
\begin{equation}
 \frac{h_0^2+(h_1/\omega)^2}{4|Q|^2}\ge\frac12,
 \qquad
 \frac{\Lag_j}{4|Q|^2\omega^{2j+2}}\ge\frac{31}{64},\quad j=0,1.
 \label{eq:explicit-signs-intro}
\end{equation}
Reflection gives the negative-height statements.
\end{proposition}
The complete explicit estimates are in Appendix~\ref{app:explicit-height}.
This enormous cutoff does not provide a practical bridge from height 112,
and it is not an evaluated cutoff for the sharper $61/24$ theorem.

The geometric interpretation is that a common nonzero amplitude controls
all four jets, including at phases where individual leading components
vanish. In quaternionic and octonionic coordinates, the same quadratic
forms become signed geometric observations. Those constructions do not
supply positivity by themselves; they identify precisely which part of
the theta analysis is needed. A second, complementary result is the
validated bounded cover
\begin{equation}
 |F(t,x)|>10^{-17}\quad(0\le t\le1/2,\ |x|\le112),
 \label{eq:intro-bounded}
\end{equation}
together with both strict stationary signs there. Its complete-source
proof and certificate structure are given in
Section~\ref{sec:fc-certificate}, including Theorem~\ref{fc:height112}.

Real multiple zeros at positive times are the relevant obstruction:
the classical heat-flow inputs imply
\begin{equation}
 \RH\quad\Longleftrightarrow\quad
 F(t,x)\ne0\quad\hbox{for every }t>0\hbox{ and }x\in\R.
 \label{eq:intro-rh}
\end{equation}
Section~\ref{sec:fc-collision} gives the reduction and the sufficient
stationary bounded-loss criterion. Theorem~\ref{thm:outer} and
\eqref{eq:intro-bounded} leave an inner small-time region and a finite
transition strip. The exact shared-source formula later in the paper
formulates the remaining arithmetic estimate without discarding its
signed cross terms.

% Prepared for Edward Baker; gpt-6-astra, effort ultra; 14 October 2026.
\section{Proof of the uniform outer-region theorem}\label{sec:outer-proof}

\subsection{A holomorphic comparison amplitude}
Use the principal logarithm and set
\begin{align}
 M_0(s)&=\frac{s(s-1)}{16}\pi^{-s/2}\sqrt{2\pi}
 \exp\left\{\left(\frac s2-\frac12\right)\operatorname{Log}\frac s2
                         -\frac s2\right\},\label{eq:M0}\\
 \alpha(s)&=\frac{M_0'(s)}{M_0(s)}
 =\frac{1}{2s}+\frac{1}{s-1}+\frac12\operatorname{Log}\frac{s}{2\pi},
 \qquad M_t(s)=e^{t\alpha(s)^2/4}M_0(s),\label{eq:alpha-Mt}\\
 s_\pm(z)&=\frac{1\pm iz}{2},\qquad
 w_\pm(z)=s_\pm(z)+\frac t2\alpha(s_\pm(z)),\label{eq:wpm}\\
 Q_\pm(z)&=M_t(s_\pm(z))\zeta(w_\pm(z)),\qquad
 \mathcal P_t(z)=Q_+(z)+Q_-(z).\label{eq:analytic-model}
\end{align}
For real $x$ we take $Q=Q_+(x)$; the other branch is its conjugate.
The functions in these definitions are analytic on the small discs
used below. Their arguments stay at imaginary parts approximately
$\pm x/2$, away from all logarithmic cuts, while $\Re w_\pm>1$ keeps
the zeta arguments away from the pole.

Write $z=u+iv$, $L_u=\log(u/(4\pi))$, and
\[
 N=\left\lfloor\sqrt{u/(4\pi)+t/16}\right\rfloor,
 \qquad S_{t,N}(w)=\sum_{n\le N}e^{t\log^2n/4}n^{-w}.
\]
Polymath's estimate applies when $u\ge200$, $0\le v\le1$, and
$0<t\le1/2$. In the present notation it reads
\begin{equation}
 H(t,z)=M_t(s_-(z))S_{t,N}(w_-(z))
       +M_t(s_+(z))S_{t,N}(w_+(z))+\mathcal E(t,z),
 \label{eq:RS-input}
\end{equation}
where $|\mathcal E|\le |M_t(s_-)|(e_A+e_B+e_{C,0})$.
This is the exact two-branch expression of the cited theorem: its first
exponent $s_*$ equals $w_-$ and its second exponent, after incorporating
the factor $n^v$, is $\overline{s_*}+\kappa-v=w_+$. To retain the full
imported error information, define
\[
 \gamma=\frac{M_t(s_+)}{M_t(s_-)},\qquad
 \kappa=\frac t2\left(\alpha\left(\frac{1-v+iu}{2}\right)
                -\alpha\left(\frac{1+v+iu}{2}\right)\right).
\]
The bounds needed from \cite[eqs. (20)--(24)]{polymath} are
\begin{align}
 |\gamma|&\le e^{0.02v}(u/(4\pi))^{-v/2},\qquad
 |\kappa|\le\frac{tv}{2(u-6)},\label{eq:RS-ratios}\\
 e_A+e_B&\le\sum_{n\le N}(1+|\gamma|N^{|\kappa|}n^v)
 e^{t\log^2n/4}n^{-\Re w_-}
 \left[\exp\left\{\frac{(t^2/16)\log^2(u/(4\pi n^2))+0.626}
                             {u-6.66}\right\}-1\right],
 \label{eq:RS-AB}\\
 e_{C,0}&\le (u/(4\pi))^{-(1+v)/4}
 \exp\left\{-\frac{tL_u^2}{16}
 +\frac{1.24(3^v+3^{-v})}{N-0.125}
 +\frac{3|L_u+i\pi/2|+10.44}{u-12}\right\}.
 \label{eq:RS-C}
\end{align}
No derivatives of this finite approximation will be taken. Its cutoff
can jump with $u$; the analytic model \eqref{eq:analytic-model} will
remove that difficulty before differentiation.

\subsection{A maximal estimate for logarithmic-phase blocks}
The gain over absolute summation comes from retaining the oscillation in
$n^{-iT}$ before estimating each arithmetic block.
\begin{lemma}[Logarithmic-phase blocks]\label{lem:logarithmic-block}
For $T\ge1$, $M\ge1$, and every interval $I$ of integers contained in
$(M,2M]$,
\begin{equation}
 \left|\sum_{n\in I}n^{-iT}\right|
 \ll T^{1/6}M^{1/2}+T^{-1/6}M.
 \label{eq:logarithmic-block}
\end{equation}
The implied constant is absolute. The same estimate holds with $T$
replaced by $-T$, and is uniform over all shortened terminal blocks.
\end{lemma}
\begin{proof}
Apply the third derivative test, in the form of
\cite[Theorem 2 and its preceding proof]{robert}, to
$f(y)=-(T/(2\pi))\log y$. On $(M,2M]$ its third derivative has
magnitude comparable to $T/M^3$, with an absolute comparison ratio.
For an interval of length $J\le M$, the test gives
\[
 J(T/M^3)^{1/6}+J^{1/2}(T/M^3)^{-1/6}
 \le T^{1/6}M^{1/2}+T^{-1/6}M.
\]
Empty intervals contribute zero; one-point intervals are absorbed by
increasing the absolute constant. Conjugation handles the opposite phase.
\end{proof}

For a continuously differentiable, possibly complex weight $W$, Abel
summation on a shortened block is the exact identity
\begin{equation}
 \sum_{M<n\le R}W(n)n^{-iT}
 =W(R)S_M(R)-\int_M^R W'(y)S_M(y)\dd y,
 \quad S_M(y)=\sum_{M<n\le y}n^{-iT},
 \quad M\le R\le2M.
 \label{eq:abel-logarithmic-block}
\end{equation}
Thus \eqref{eq:logarithmic-block} is multiplied by
$|W(R)|+\int_M^R|W'(y)|\dd y$. For a nonnegative decreasing weight,
this total-variation factor is exactly $W(M)$. The maximal form of the
lemma is needed because the last block ends at the actual cutoff $N$.

\subsection{Five exponent pairs and a closed outer boundary}
We strengthen the elementary third-derivative estimate above using five
limiting exponent pairs. Write
\begin{equation}
 \begin{aligned}
 P&=(13/84,55/84),& AP&=(13/194,76/97),\\
 A^2P&=(13/414,359/414),& D&=(18/199,593/796),\\
 R&=(4742/38463,35731/51284).
 \end{aligned}
 \label{eq:closed-pairs}
\end{equation}
Here $D$ is the name of the displayed pair, not an assertion that an
arbitrary transformation called $D$ preserves exponent pairs.
Bourgain's theorem supplies $P$ with arbitrarily small positive losses
\cite[Theorem 6]{bourgain14}; Trudgian and Yang supply $D,R$ and the
$A$-process
\cite[Section 1.1, Lemma 1.1, and eq.~(1.10)]{trudgianyang23},
\[
 A(k,l)=\left(\frac{k}{2k+2},\frac{k+l+1}{2k+2}\right).
\]
Their exponent-pair definition is uniform over arbitrary subintervals.
For $f(y)=(T/(2\pi))\log y$, its first derivative is exactly the model
monomial with exponent parameter one, so all derivative-approximation
errors are zero. Conjugation handles the opposite phase. Consequently,
for each fixed $\eta>0$, every pair $(k,l)$ in
\eqref{eq:closed-pairs} gives the complete maximal estimate
\begin{equation}
 \sup_{I\subset(M,2M]}
 \left|\sum_{n\in I}n^{-iT}\right|
 \ll_\eta T^{k+\eta}M^{l-k},
 \qquad M\ge1,\quad T\ge2\pi M.
 \label{eq:closed-maximal-block}
\end{equation}
The supremum is over integer intervals, including shortened terminal
blocks. Bourgain's proper-subinterval completion may cost a logarithm;
choose a smaller initial positive loss to absorb it into the displayed
$\eta$. No zero-loss endpoint theorem is used.

For clarity, the descendants of $P$ also have admissibly small losses.
Under equal positive increments $\epsilon$ in both coordinates, the
increments in the coordinates of $A(k,l)$ are
\[
 \frac{\epsilon}{2(1+k)(1+k+\epsilon)},\qquad
 \frac{\epsilon(1+k-l)}{2(1+k)(1+k+\epsilon)}.
\]
They are nonnegative and at most $\epsilon/2$ for the pairs used here;
a second application remains continuous. If the final coordinate
increments are at most $\eta$, their contribution to the logarithm of
the bound is at most
$\eta(\log T-\log M)+\eta\log M=\eta\log T$.
Thus \eqref{eq:closed-maximal-block} includes all losses and is a full
bound in the indicated range.

Set
\[
 \theta_*:=\frac{61}{24},\qquad \eta=\rho:=\frac1{100000}.
\]
For $(t,x)$ in \eqref{eq:outer-domain}, put
\begin{equation}
 D_x=\{z:|z-x|\le L^{-1}\},\qquad
 L_u=\log(u/(4\pi)),\qquad \theta_u=tL_u,
 \quad z=u+iv.
 \label{eq:shrinking-disc}
\end{equation}
In this proof write $N_{\rm RS}$ for the integer $N$ in
\eqref{eq:RS-input}; it will not be confused with the source cutoff
introduced in Section~\ref{sec:residuals}. Uniformly on $D_x$, the
explicit formula \eqref{eq:alpha-Mt} gives, for either branch
$w=\sigma+i\tau=w_\pm(z)$,
\begin{equation}
 \begin{gathered}
 L_u=L+O((xL)^{-1}),\quad |\tau|\asymp u,\quad
 \log N_{\rm RS}\le L_u/2+O(u^{-1}),\\
 \theta_u\ge\theta_*-\rho,\qquad
 \sigma\ge\tfrac12+\tfrac{\theta_u}{4}-\delta_x,\qquad
 0\le\delta_x=O(L^{-1})\le\rho.
 \end{gathered}
 \label{eq:exponent-bounds}
\end{equation}
These estimates hold after one fixed increase of $X_*$ and are uniform
for $0<t\le1/2$, including growing $\theta_u\le L_u/2$. In particular
$\sigma\ge9/8$ there. The real disc is allowed to cross slightly below
$\theta_*$; the following strict margins pay that displacement.

For $\beta=(\log M)/L_u$ define the unperturbed weighted exponent
\begin{equation}
 F_{k,l}(\beta;\theta)
 =k+(l-k-\tfrac12)\beta-\frac{\theta}{4}\beta(1-\beta).
 \label{eq:closed-weighted-exponent}
\end{equation}
At $\theta=\theta_*$ it is convex, with second derivative $61/48$.
The following endpoint values are exact.
\begin{center}
\small
\begin{tabular}{@{}llll@{}}
\toprule
Pair & Interval for $\beta$ & Left value & Right value\\
\midrule
$A^2P$ & $[1/5,3/10]$ & $-43/13800$ & $-869/662400$\\
$AP$ & $[3/10,3/8]$ & $-459/310400$ & $-145/198656$\\
$D$ & $[3/8,81/200]$ & $-215/407552$ & $-21807/254720000$\\
$R$ & $[81/200,3/7]$ & $-4368359/49232640000$ & $-1/1176$\\
$P$ & $[3/7,1/2]$ & $-1/1176$ & $-11/2688$\\
\bottomrule
\end{tabular}
\end{center}
The intervals cover $[1/5,1/2]$ without gaps. Convexity therefore gives,
with the corresponding pair in each interval,
\begin{equation}
 F_{k,l}(\beta;\theta_*)\le-\mu_*,\qquad
 \mu_*:=\frac{21807}{254720000}.
 \label{eq:closed-margin}
\end{equation}
The positive exponent-pair loss costs at most $\eta$, the real-part
error at most $\delta_x/2$, and the decrease of $\theta_u$ below
$\theta_*$ at most $\rho/16$. For the tiny possible extension
$\beta\le1/2+O((uL_u)^{-1})$, continue to use $P$. Its principal change
beyond $1/2$ is $\theta_u(\beta-1/2)^2/4$, which tends to zero uniformly
even when $\theta_u$ grows; the other perturbations do also. Increase
$X_*$ until this extension costs at most $\rho$. The surviving margin is
\begin{equation}
 \mu_*-\rho(1+\tfrac12+\tfrac1{16}+1)
 =\frac{76399}{1273600000}>\frac1{20000}.
 \label{eq:closed-paid-margin}
\end{equation}
Replacing $T=|\tau|$ by a constant multiple of $u$ changes bounded
factors only. All these losses are paid at the actual lower real
boundary $\theta=61/24$.

\subsection{Uniform replacement of the weighted sums at the closed boundary}
\begin{lemma}[Weighted logarithmic moments]\label{lem:weighted-moments}
Under \eqref{eq:exponent-bounds}, for every fixed nonnegative integer $r$,
\begin{equation}
 \sum_{n\le N_{\rm RS}}(\log n)^r e^{t\log^2n/4}n^{-w}
       -(-1)^r\zeta^{(r)}(w)=O_r(L^{-1}),\qquad z\in D_x.
 \label{eq:sum-replacement}
\end{equation}
For $r=0$ the factor at $n=1$ is interpreted as one. The implied
constant is uniform throughout \eqref{eq:outer-domain}.
\end{lemma}
\begin{proof}
Write $\ell=L_u$ and split at $Y=e^{\ell/5}$. For $n\le Y$, use
$e^q-1\le qe^q$ and
$e^{t\log^2n/4}\le n^{\theta_u/20}$. The small-index discrepancy is
at most
\begin{equation}
 \frac t4\sum_{n\ge2}(\log n)^{r+2}
                    n^{-(1/2+\theta_u/5-\delta_x)}.
 \label{eq:weighted-zeta-difference}
\end{equation}
Put $\theta_u=\theta_*-\rho+s$ with $s\ge0$. Since
\[
 \frac12+\frac{\theta_*}{5}=\frac{121}{120},\qquad
 \frac1{120}-\frac\rho5-\rho>\frac1{240},
\]
the exponent in \eqref{eq:weighted-zeta-difference} is at least
$1+1/240+s/5$. Hence, with the fixed finite constant
$D_r=\sum_{n\ge2}(\log n)^{r+2}n^{-1-1/240}$, that discrepancy is
bounded by
\begin{equation}
 \frac{D_r}{4\ell}(\theta_*-\rho+s)2^{-s/5}=O_r(\ell^{-1}).
 \label{eq:uniform-growing-theta}
\end{equation}
The displayed numerator is bounded on the whole half-line $s\ge0$;
this proves uniformity for growing $\theta_u$, rather than merely for
time parameters in a fixed compact substrip. Since $\sigma\ge9/8$,
the omitted ordinary zeta tail satisfies
\begin{equation}
 \sum_{n>Y}(\log n)^r n^{-\sigma}
       =O_r(\ell^r e^{-\ell/40}).
 \label{eq:ordinary-zeta-tail-three}
\end{equation}

Partition $Y<n\le N_{\rm RS}$ into the exact dyadic intervals
$(M,\min(2M,N_{\rm RS})]$, starting at $M=Y$. The weight
\[
 W_r(y)=(\log y)^r e^{t\log^2y/4}y^{-\sigma}
\]
is nonnegative and, for sufficiently large $\ell$, decreasing, because
\begin{equation}
 y(\log W_r)'=\frac r{\log y}+\frac t2\log y-\sigma
 \le\frac{5r}{\ell}-\frac12+\delta_x+O(u^{-1})<0.
 \label{eq:weighted-block-monotonicity}
\end{equation}
Apply \eqref{eq:abel-logarithmic-block} and
\eqref{eq:closed-maximal-block}, choosing the pair by the left endpoint
$\beta=(\log M)/\ell$. The weighted block is at most a constant times
$W_r(M)T^{k+\eta}M^{l-k}$. A block may cross a boundary in the table:
the estimate is still valid, because \eqref{eq:closed-maximal-block}
already controls every shortened interval at its given scale $M$.
Equations \eqref{eq:closed-margin}--\eqref{eq:closed-paid-margin}
therefore yield
\begin{equation}
 \sum_{Y<n\le N_{\rm RS}}(\log n)^r e^{t\log^2n/4}n^{-w}
       =O_r(\ell^{r+1}e^{-\ell/20000}).
 \label{eq:high-index-decay}
\end{equation}
There are $O(\ell)$ blocks, including the shortened final block. Combining
\eqref{eq:weighted-zeta-difference},
\eqref{eq:ordinary-zeta-tail-three}, and \eqref{eq:high-index-decay}
proves the lemma.
\end{proof}

\subsection{Complete value errors and derivative control}
The positive sums in the imported errors require no further arithmetic
cancellation. For $n\le N_{\rm RS}$,
\[
 e^{t\log^2n/4}n^{-\theta_u/4}
 =\exp\{-\tfrac t4\log n(\ell-\log n)\}\ll1,
 \qquad n^{\delta_x}\ll1.
\]
On the upper half-disc \eqref{eq:RS-ratios} bounds $|\gamma|$ and
$N_{\rm RS}^{|\kappa|}$ uniformly. The factor $n^v$ replaces the first
branch's real part $(1+v)/2$ by $(1-v)/2$, already covered in
\eqref{eq:exponent-bounds}. Thus the complete positive weight sum in
\eqref{eq:RS-AB} is $O(\sum_{n\le N_{\rm RS}}n^{-1/2})=O(u^{1/4})$.
Its per-index error multiplier is $O(L^2/u)$, uniformly for $t\le1/2$.
The correction exponent in \eqref{eq:RS-C} tends uniformly to zero.
Consequently all three imported errors obey
\begin{equation}
 e_A+e_B=O(L^2u^{-3/4}),\qquad
 e_{C,0}\ll e^{-\ell/4-\theta_u\ell/16}=O(u^{-2/5}),
 \label{eq:RS-small-errors}
\end{equation}
where $1/4+(\theta_*-\rho)/16>2/5$.

Set $\Gamma=|M_t(s_+(x))|$. The logarithmic $z$-derivative of either
$M_t(s_\pm(z))$ is $O(L)$, so integration over a segment of length
at most $1/L$ gives
\begin{equation}
 |M_t(s_\pm(z))|\asymp\Gamma\qquad(z\in D_x).
 \label{eq:gamma-comparison}
\end{equation}
Combining \eqref{eq:RS-input}, Lemma~\ref{lem:weighted-moments} at
$r=0$, and \eqref{eq:RS-small-errors}--\eqref{eq:gamma-comparison}
gives
\begin{equation}
 H(t,z)-\mathcal P_t(z)=O(\Gamma/L),\qquad z\in D_x.
 \label{eq:analytic-value-error}
\end{equation}
Initially the imported approximation applies to the upper half-disc;
Schwarz symmetry supplies the lower half.

Absolute convergence of the Dirichlet series for $\zeta$ and
$1/\zeta$ in $\Re w\ge9/8$ gives
\begin{equation}
 \zeta(9/8)^{-1}\le|\zeta(w)|\le\zeta(9/8).
 \label{eq:zeta-nonzero}
\end{equation}
Thus $Q\ne0$ and $|Q|\asymp\Gamma$. The complete difference in
\eqref{eq:analytic-value-error} is holomorphic on a neighborhood of
the closed disc. Cauchy's formula, with time fixed, yields
\begin{equation}
 \partial_x^k(H(t,\cdot)-\mathcal P_t)(x)
       =O_k(|Q|L^{k-1}),\qquad k=0,1,2,3.
 \label{eq:analytic-jet-error}
\end{equation}
At $k=0$ this means $O(|Q|/L)$. The moving Riemann--Siegel floor has
been eliminated before differentiation.

\subsection{The common phase and the two signs}
At the real center the explicit carrier satisfies
\begin{equation}
 \alpha(s_+(x))=L/2+i\pi/4+O(x^{-1}),\qquad
 \alpha^{(m)}(s_+(x))=O_m(x^{-m}),\quad m\ge1.
 \label{eq:alpha-asymptotics}
\end{equation}
Fixed-order logarithmic derivatives of $\zeta$ are uniformly bounded
in $\Re w\ge9/8$. Therefore
\begin{align}
 (\log Q_+)'&=\frac i2\left[\alpha+\frac t2\alpha\alpha'
  +\left(1+\frac t2\alpha'\right)
                        \frac{\zeta'(w_+)}{\zeta(w_+)}\right]
             =i\omega+O(1),\label{eq:phase-derivative}\\
 (\log Q_+)''&=O(1),\qquad
 (\log Q_+)'''=O(1).\label{eq:higher-log-derivatives}
\end{align}
Here all $\alpha$ terms are evaluated at $s_+(x)$ and time is fixed.
The estimates remain uniform as $\theta=tL$ grows because $t\le1/2$.
The exponential derivative formulas give
\[
 Q_+^{(k)}(x)=Q\{(i\omega)^k+O(L^{k-1})\},\qquad k=1,2,3;
\]
the real-axis derivatives of $Q_-$ are their conjugates. Equation
\eqref{eq:analytic-jet-error} proves \eqref{eq:main-four-jet} with
$O(L^{-1})$, throughout the closed domain \eqref{eq:outer-domain}.

Write $P_k=h_k/(2|Q|\omega^k)$ and
$P=(c,-s,-c,s)+e$, where $c^2+s^2=1$. If $|e_k|\le d$, direct
expansion gives
\begin{equation}
 P_0^2+P_1^2\ge1-4d,\qquad
 P_1^2-P_0P_2,\ P_2^2-P_1P_3\ge1-4d-d^2.
 \label{eq:explicit-perturbation}
\end{equation}
For example the first signed form equals
$1-2se_1+e_1^2-ce_2+ce_0-e_0e_2$; retain its nonnegative square
and bound the other terms. Increasing $X_*$ until $\omega\ge1$ and
$d\le1/8$ proves Theorem~\ref{thm:outer} and the source-defined envelopes
\begin{equation}
 \begin{gathered}
 E\ge2|Q|^2\ge\frac{2\Gamma^2}{\zeta(9/8)^2},\\
 \Lag_0\ge\frac{31}{16}|Q|^2\omega^2,\qquad
 \Lag_1\ge\frac{31}{16}|Q|^2\omega^4.
 \end{gathered}
 \label{eq:closed-sign-envelopes}
\end{equation}
The endpoint $\theta=61/24$ is included because the arithmetic margins
were strictly positive before the disc and exponent-pair losses were
paid. This strengthens the earlier $8/3+\varepsilon$ result, and in
particular covers its formerly excluded boundary $8/3$. The height
$X_*$ remains unevaluated. Appendix~\ref{app:explicit-height} retains
the independent explicit height under the stronger condition $tL\ge5$.

\begin{corollary}[Asymptotic collision localization]
\label{cor:threshold-closed}\label{cor:threshold-three}\label{cor:threshold-four}
If $H(T_m,a_m)=H_x(T_m,a_m)=0$, $0<T_m\le1/2$, and
$|a_m|\to\infty$, then every sufficiently large $m$ satisfies
\[
 T_m\log\frac{|a_m|}{4\pi}<\frac{61}{24}.
\]
In particular the limsup of these products is at most $61/24$.
\end{corollary}
\begin{proof}
For $|a_m|\ge X_*$, a product at least $61/24$ contradicts the energy
lower envelope \eqref{eq:closed-sign-envelopes}, using reflection when
$a_m<0$.
\end{proof}
This localizes any hypothetical collision sequence; it does not assert
that collisions exist or impose this threshold below the unevaluated
height cutoff.

\subsection{A conditional continuation to five halves}
The largest blocks alone do not control the remaining interval:
\[
 F_P(1/2;5/2)=-\frac1{672},\qquad
 F_D(2/5;5/2)=\frac9{3980}>0.
\]
At $\beta=2/5$, $D$ is the best of the five displayed pairs. The positive
upper exponent prevents these particular estimates from proving a small
remainder there. It does not prove growth of the actual sum or failure
of a theta sign. Absolute treatment of low indices at $\theta=5/2$
requires a split $Y=e^{\alpha\ell}$ with $\alpha<1/5$, so merely
moving that split beyond the failing middle block is not a repair.

\begin{lemma}[Conditional middle-block transfer]\label{lem:middle-block-transfer}
Suppose there are fixed constants $C,\kappa>0$ and $T_0$ such that
for $T\ge T_0$, $\beta=\log M/\log T\in[7/25,47/100]$, and every
integer interval $I\subset(M,2M]$,
\begin{equation}
 \left|\sum_{n\in I}n^{-iT}\right|
 \le C T^{\beta/2+(5/8)\beta(1-\beta)-\kappa}.
 \tag{MB}\label{eq:middle-block-hypothesis}
\end{equation}
Then, for sufficiently large $x$, uniformly for $0<t\le1/2$ and the
entire closed strip $5/2\le tL\le8/3$, the four-jet asymptotic
\eqref{eq:main-four-jet} holds with $O(L^{-1})$, and
$E,\Lag_0,\Lag_1$ are strictly positive. The two cutoff conclusions of
Corollary~\ref{cor:signed-staircase-cutoffs} also hold there, with
$\zeta(11/10)$ replacing $\zeta(9/8)$ in the relative coefficient.
The hypothesis \eqref{eq:middle-block-hypothesis} is not established here.
\end{lemma}
\begin{proof}
Use the same shrinking disc and put $Y=e^{19\ell/100}$. At the real
lower boundary $\theta=5/2$, the exponent in the absolute small-index
discrepancy is
\[
 \frac12+\frac{(5/2)(1-19/100)}4=\frac{161}{160}>1.
\]
For sufficiently large height the disc perturbations leave it at least
$1+1/320$. Since $\theta$ is bounded on this strip, the discrepancy is
$O_r(L^{-1})$. Also $\Re w\ge11/10$, so the ordinary zeta tail is
$O_r(\ell^r e^{-19\ell/1000})$.

For the high blocks outside the assumed middle range, use $A^2P$ on
$[19/100,2/7]$ and $P$ on $[23/50,1/2]$. Their exact endpoint
values at $\theta=5/2$ are
\[
 \begin{array}{c|cc}
 A^2P&-3293/3312000&-1/4508\\
 P&-41/84000&-1/672.
 \end{array}
\]
Convexity gives strict savings on both full intervals, preserved by
sufficiently small positive exponent-pair losses. The middle range in
\eqref{eq:middle-block-hypothesis} overlaps these intervals by $1/175$
and $1/100$, respectively. Those fixed overlaps absorb the difference
between $\log T=\ell+O(1)$ and the scale $\ell$ used on the complex disc.

For a remaining block, write $b=\log M/\ell$, while the hypothesis uses
$\beta=\log M/\log T$. Multiplication by the actual decreasing weight
and Abel summation give, apart from fixed powers of $\ell$, a logarithm
bounded by
\[
 -\kappa\ell-\frac{\theta_u-5/2}{4}b(1-b)\ell+O(1).
\]
The $O(1)$ includes the logarithmic-coordinate change and the real-part
error, since $\delta_x\ell=O(1)$. On this fixed strip the disc's
possible displacement below $\theta=5/2$ tends to zero. Thus the bound
is at most $-\kappa\ell/2$ after increasing the starting height.
Summing the $O(\ell)$ blocks proves Lemma~\ref{lem:weighted-moments}
on the entire closed strip, conditionally on
\eqref{eq:middle-block-hypothesis}.

All imported errors remain negligible; in particular
$e_{C,0}=O(u^{-2/5})$ since $1/4+(5/2)/16=13/32>2/5$, with room
for the vanishing disc perturbation. The analytic comparison, Cauchy
estimates, and phase-uniform algebra above apply with the zero-free
half-plane $\Re w\ge11/10$. They give
$E\ge2\Gamma^2/\zeta(11/10)^2$ and the same two positive Laguerre
envelopes. The proof of Corollary~\ref{cor:signed-staircase-cutoffs}
then gives both cutoff conclusions.
\end{proof}
Together with Theorem~\ref{thm:outer}, the hypothesis would cover all
$tL\ge5/2$ in the allowed time range at sufficiently large height.
It is a sufficient arithmetic condition, not a necessary condition or
an RH-equivalent reformulation. A direct estimate for the complete
weighted or signed physical observation could bypass it. Neither this
conditional transfer nor the unconditional theorem settles the inner
cusp or evaluates the finite transition to the asymptotic region.

\section{Why positive-time collisions are the obstruction}
\label{sec:fc-collision}

The preceding asymptotic result is useful because a failure of RH forces
a finite collision of real zeros at a positive time.  We record the
precise reduction, including the compactness input, before formulating
the signed estimate still needed outside the proved regions.

We use the normalization
\[
 H(0,z)=\frac18\xi\!\left(\frac12+\frac{iz}{2}\right),
 \qquad H_t=-H_{zz}.
\]
The positive, rapidly decreasing theta source makes $H$ jointly entire
in $(t,z)$, even in $z$, and real on the real $(t,z)$ plane.  Moreover,
$H(t,iy)>0$ for real $t,y$, so no time slice is identically zero and
every zero has finite multiplicity.

The classical de Bruijn--Newman theory gives a finite real threshold
$\Lambda$ such that all zeros of $H(t,\cdot)$ are real exactly when
$t\ge\Lambda$; the endpoint is included
\cite{debruijn,newman}.  In this normalization $\Lambda\le1/2$ and
RH is equivalent to $\Lambda\le0$.  De Bruijn's strip contraction
keeps all zeros in $|\operatorname{Im}z|\le1$ for $t\ge0$.
We also use the uniform confinement theorem of
\cite[Theorem 1.5(i)]{polymath}: there is an absolute constant $C>0$
such that, for $0<t\le1/2$, a nonreal zero satisfies
\begin{equation}
 |\operatorname{Re}z|<e^{C/t}.
 \label{fc:uniform-confinement}
\end{equation}
Uniformity in time, rather than a separate uncontrolled radius for
each slice, is what prevents a first collision from escaping to infinity.

\begin{lemma}[Local splitting at a multiple zero]
\label{fc:hermite}
Let $W$ be real analytic near $(T,a)$, solve $W_t=-W_{xx}$, and have
an exact spatial zero of multiplicity $m\ge2$ there.  Its nearby zeros
have leading locations
\[
 a+\sqrt{2(t-T)}\,\lambda_k+O(|t-T|),\qquad k=1,\ldots,m,
\]
where the $\lambda_k$ are the distinct real roots of the probabilists'
Hermite polynomial $\operatorname{He}_m$.  In particular, $W$ has a
nonreal zero at times immediately before $T$.
\end{lemma}
\begin{proof}
Put $W_q=\partial_x^qW(T,a)$.  The convergent heat Taylor series is
\[
 W(T+\delta,a+w)=
 \sum_{p,q\ge0}\frac{(-\delta)^p w^q}{p!q!}W_{2p+q}.
\]
With $\delta=\eta^2$ and $w=\sqrt2\eta\lambda$, division by $\eta^m$
extends analytically to $\eta=0$, where its value is
$2^{m/2}W_m\operatorname{He}_m(\lambda)/m!$.
The roots of $\operatorname{He}_m$ are real and simple: its Rodrigues
formula gives orthogonality to all lower-degree polynomials for the
positive Gaussian weight, and the polynomial formed from its real
sign-changing factors would contradict that orthogonality unless there
were $m$ such roots. The simple roots therefore continue by the implicit function
theorem.  For $t<T$, take $\eta=i\sqrt{T-t}$; at least one Hermite root
is nonzero, so its leading imaginary displacement cannot be cancelled
by the $O(T-t)$ remainder.
\end{proof}

\begin{proposition}[Finite collision criterion for RH]
\label{fc:rh-collision}
RH holds if and only if
\begin{equation}
 H(t,x)=H_x(t,x)=0
 \quad\hbox{has no solution with }0<t\le1/2,\ x\in\R.
 \label{fc:collision-criterion}
\end{equation}
Equivalently, $E(t,x)=h_0(t,x)^2+h_1(t,x)^2$ is positive throughout
this region.  No uniform positive lower bound for $E$ is required. Equivalently,
there is no such joint zero at any positive time.
\end{proposition}
\begin{proof}
If RH holds, all zeros are real at all nonnegative times.  A multiple
zero at a positive time would contradict Lemma~\ref{fc:hermite}.

Conversely, if RH fails then $0<\Lambda\le1/2$.  Choose
$t_k\uparrow\Lambda$ with $t_k\ge\Lambda/2$, and a nonreal zero
$z_k$ at each $t_k$.  The strip bound and
\eqref{fc:uniform-confinement} put these zeros in the fixed compact set
$|\operatorname{Im}z|\le1$, $|\operatorname{Re}z|\le e^{2C/\Lambda}$.
A subsequence converges to a zero $a$ of $H(\Lambda,\cdot)$, which is
real by the threshold theorem.  If that zero were simple, the local
implicit zero branch would be unique; conjugation symmetry would make
it real for real times on both sides of $\Lambda$.  This contradicts
$z_k\notin\R$.  Thus $H(\Lambda,a)=H_x(\Lambda,a)=0$.
The source positivity gives $H(\Lambda,0)>0$, so $a\ne0$.
\end{proof}

\subsection{A sufficient estimate only on two stationary sets}

Recall
\[
 \mathcal L_j=h_{j+1}^2-h_jh_{j+2},\qquad j=0,1.
\]
On a smooth stationary curve $h_{j+1}(t,\chi(t))=0$, the heat equation
gives
\[
 \frac12\frac{d}{dt}h_j(t,\chi(t))^2=\mathcal L_j(t,\chi(t)).
\]
The following result includes singular stationary sets and all finite
multiplicities.  Its losses may grow with height or as $t\downarrow0$;
only local boundedness through each possible positive-time collision
is required.

\begin{lemma}[A stationary predecessor of an even zero]
\label{fc:predecessor}
Suppose $W_t=-W_{xx}$ and $W$ has exact spatial multiplicity $2r$ at
$(T,a)$, where $r\ge1$.  With $W_q=\partial_x^qW(T,a)$, there is a
real analytic function $\chi(h)$ for small $h\ge0$ such that
\begin{align*}
 W_x(T-h,\chi(h))&=0,&
 \chi(h)&=a-\frac{W_{2r+1}}{rW_{2r}}h+O(h^2),\\
 W(T-h,\chi(h))&=\frac{W_{2r}}{r!}h^r+O(h^{r+1}),&
 W_{xx}(T-h,\chi(h))&=\frac{W_{2r}}{(r-1)!}h^{r-1}+O(h^r).
\end{align*}
\end{lemma}
\begin{proof}
Use the heat Taylor series above with $\delta=-h$ and $w=hy$.
The function $h^{-r}W_x(T-h,a+hy)$ extends analytically to $h=0$,
where it equals
$W_{2r}y/(r-1)!+W_{2r+1}/r!$.
Its $y$ derivative is nonzero.  The implicit function theorem gives
$y(h)$, and substitution in the same series gives the two displayed
asymptotics.
\end{proof}

\begin{theorem}[Stationary bounded-loss criterion]
\label{fc:bounded-loss}
Suppose there are nonnegative functions $K_0,K_1$, locally bounded
on $(0,1/2]\times\R$, for which
\begin{equation}
 \begin{split}
 h_1=0&\quad\Longrightarrow\quad\mathcal L_0\ge-K_0E,\\
 h_2=0&\quad\Longrightarrow\quad\mathcal L_1\ge-K_1E.
 \end{split}
 \label{fc:stationary-loss}
\end{equation}
Then RH holds.  Local boundedness here includes a neighborhood of
every candidate joint zero, using a left neighborhood in time at
$t=1/2$.
\end{theorem}
\begin{proof}
Suppose $H$ has a real multiple zero at $(T,a)$ in this time range,
of exact multiplicity $m\ge2$.
If $m=2r$, apply Lemma~\ref{fc:predecessor} to $W=H$.
On its predecessor branch $h_1=0$, $E=H^2>0$, and
\[
 -\frac{\mathcal L_0}{E}=\frac{H_{xx}}H
 =\frac rh+O(1),\qquad h=T-t\downarrow0.
\]
This contradicts the first inequality with locally bounded $K_0$.
If $m=2r+1$, apply the lemma to $W=H_x$.
Then $h_2=0$, $H_x=c h^r+O(h^{r+1})$, and
$H_{xxx}=rc h^{r-1}+O(h^r)$, with $c\ne0$.
The heat Taylor series also gives $H=O(h^{r+1})$ on this branch,
because its displacement from $a$ is $O(h)$.  Consequently
\[
 E=H_x^2(1+O(h^2)),\qquad
 -\frac{\mathcal L_1}{E}=\frac rh+O(1),
\]
contradicting the second inequality.  All divisions occur on proved
nonzero predecessor branches, not at the supposed collision.
Proposition~\ref{fc:rh-collision} now applies.
\end{proof}

Both levels matter: the heat polynomial
$H(t,x)=x^3-6(t-T)x$ has
$\mathcal L_0=3x^4+36(t-T)^2\ge0$ everywhere, but has a triple zero
at $(T,0)$.  Its second stationary loss diverges on approaching that
collision from earlier times.  Thus positivity of the first observation
alone is insufficient for arbitrary multiplicity.



% Prepared for Edward Baker, 14 October 2026.
% Substantial LLM assistance: OpenAI gpt-6-astra, effort ultra,
% as verified and supplied by the parent session for this turn.
% Standalone section fragment. Required macros: \R,\C,\HH,\OO,\ii,\jj,\kk.
% Required citation keys: sudbery, baez, kauhanenorelma.

\subsection{The finite part of the collision pole and delayed cancellation}
\label{subsec:collision-finite-part}

The preceding collision argument uses only the leading pole.  Its next
coefficient also identifies cancellations that a delayed source estimate
must preserve.  All statements in this subsection are conditional on a
finite collision; they do not assert that the theta family has one.

\begin{proposition}[Finite part on the stationary predecessor]
\label{prop:collision-finite-part}
Suppose $H(T,\cdot)$ has exact multiplicity $m=2r+j$ at $x_*$, where
$T>0$, $r\ge1$, and $j\in\{0,1\}$. Put
\[
 \begin{gathered}
 A=h_m(T,x_*)\ne0,\qquad B=h_{m+1}(T,x_*),\qquad
 D=h_{m+2}(T,x_*),\\
 k=A/r!,\qquad \alpha_* =\frac{D}{(r+1)A}-\frac{B^2}{2rA^2}.
 \end{gathered}
\]
On the predecessor $h_{j+1}(T-s,\chi(s))=0$ of
Lemma~\ref{fc:predecessor}, as $s\downarrow0$,
\begin{align}
 \chi(s)&=x_*-\frac{B}{rA}s+O(s^2),\notag\\
 h_j(T-s,\chi(s))&=k s^r(1+\alpha_*s+O(s^2)),\notag\\
 \frac{\Lag_j}{E}&=-\frac rs-\alpha_*+O(s),
 \qquad E\sim k^2s^{2r},\qquad
 \Lag_j\sim-rk^2s^{2r-1}.
 \label{eq:collision-finite-part}
\end{align}
For odd multiplicity, the more precise comparison is
\begin{equation}
 \frac{H}{H_x}=-\frac{B}{r(r+1)A}s+O(s^2),\qquad
 E=h_1^2\left(1+\frac{B^2s^2}{r^2(r+1)^2A^2}+O(s^3)\right).
 \label{eq:odd-collision-energy}
\end{equation}
\end{proposition}
\begin{proof}
For $W=h_j$, the heat Taylor expansion at the collision is
\[
 W(T-s,x_*+y)=\sum_{p,q\ge0}
 \frac{s^py^q}{p!q!}\,h_{j+2p+q}(T,x_*).
\]
The implicit-function construction in Lemma~\ref{fc:predecessor}
gives $y=-Bs/(rA)+O(s^2)$.  The terms of order $s^{r+1}$ in $W$
are $Ds^{r+1}/(r+1)!$, $Bs^ry/r!$, and
$As^{r-1}y^2/[2(r-1)!]$.  Their sum, divided by $ks^r$,
is $\alpha_*s$, proving the second expansion.  Along the branch,
$dW(T-s,\chi(s))/ds=W_{xx}$ exactly, so
$W_{xx}/W=r/s+\alpha_*+O(s)$.  For $j=0$, $E=W^2$ and the result
follows.  For $j=1$, the same series for $H$ gives
\[
 H(T-s,\chi(s))=-\frac{B}{r(r+1)r!}s^{r+1}+O(s^{r+2}),
\]
which proves \eqref{eq:odd-collision-energy}.  Since
$\Lag_1=-W W_{xx}$ on this branch, the relative $O(s^2)$ energy
correction changes the displayed quotient only by $O(s)$.
\end{proof}

The moving energy derivative must be distinguished from the fixed-position
one.  With $\xi(t)=\chi(T-t)$,
\[
 \frac d{dt}E(t,\xi(t))=
 \begin{cases}2\Lag_0,&h_1=0,\\
 2\Lag_1+2\xi' h_0h_1,&h_2=0.
 \end{cases}
\]
The extra term divided by $E$ is $O(s)$ on the odd predecessor.
The exact branch identity without that term is
$\frac12\frac d{dt}h_1^2=\Lag_1$.
For odd multiplicity, the surviving real zero has displacement
$-Bs/((r+1)A)+O(s^2)$ from $x_*$: substituting $y=sz$ in the heat
series for $H$ and applying the implicit function theorem gives this
coefficient.  Its displacement from the stationary predecessor is
$Bs/[r(r+1)A]+O(s^2)$, consistently with
\eqref{eq:odd-collision-energy}.

Fix $a=\pi/16$, $b=1/4$, and $\mu=a^2+b^2$, and retain
$C_j=-h_{j+2}-2ah_{j+1}-\mu h_j$.  For any declared delay put
\[
 G_j=C_j(t-\delta,x),\qquad P_j=h_jG_j,\qquad
 Q_j=h_j(C_j-G_j/2).
\]
At current stationarity the exact combination is
\begin{equation}
 \frac{P_j/2+Q_j}{E}
 =\frac{\Lag_j}{E}-\mu\frac{h_j^2}{E}
 =-\frac rs-(\alpha_*+\mu)+O(s).
 \label{eq:delay-collision-combination}
\end{equation}
Thus changing the delay changes the two summands but not the observation
whose adverse part must be controlled.

\begin{proposition}[Stronger poles introduced by a positive delay]
\label{prop:positive-delay-cancellation}
Suppose the prescribed delay is analytic near the candidate event,
$\delta_*:=\delta(T,x_*)>0$, and put
$g=C_j(T-\delta_*,x_*)$.  If $g\ne0$, then on the predecessor
\[
 \frac{P_j}{E}=\frac gk s^{-r}(1+O(s)),\qquad
 \frac{Q_j}{E}=-\frac{P_j}{2E}-\frac rs-(\alpha_*+\mu)+O(s).
\]
For $r>1$, the two products therefore have opposite leading poles of
order $s^{-r}$.  Their intervening negative powers also cancel in
\eqref{eq:delay-collision-combination}.  If
$A_{\min}=(-P_j/E)_+$ and $B_{\min}=(-Q_j/E)_+$ are used only as
diagnostic quantities on this nonzero branch, then for $r>1$ and $g\ne0$
\[
 A_{\min}/2+B_{\min}\sim\frac{|g/k|}{2}s^{-r}.
\]
Both split ratios are $O(s^{-1})$ if and only if
$G_j=O(s^{r-1})$.  For $r=1$, writing $\gamma=g/k$ (including
$g=0$),
\[
 P_j/E=\gamma/s+O(1),\qquad
 Q_j/E=-(1+\gamma/2)/s+O(1).
\]
The necessary leading conditions for $P_j\ge0$ and
$Q_j\ge-KE$ with locally bounded $K$ are respectively
$\gamma\ge0$ and $\gamma\le-2$, and are incompatible.
\end{proposition}
\begin{proof}
Analyticity gives $G_j=g+O(s)$, while
$h_j=ks^r(1+O(s))$ and $E=h_j^2(1+O(s^2))$; the even case has
the exact equality $E=h_j^2$.  These facts prove the first expansion,
and the second follows from \eqref{eq:delay-collision-combination}.
Taking the negative part according to the sign of $g/k$ proves the
diagnostic formula when $r>1$.  Finally
$P_j/E=G_j/[ks^r](1+O(s))$ proves the equivalence; the combined
identity gives the same bound for $Q_j/E$.  The case $r=1$ follows
without assuming $g\ne0$.
\end{proof}

For comparison, a diagnostic delay $\delta=\rho s+O(s^2)$ with fixed
$\rho\ge0$ gives, from the leading heat polynomial,
\[
 \frac{P_j}{E}=-\frac{r(1+\rho)^{r-1}}s+O(1),\qquad
 \frac{Q_j}{E}=
 \frac{r\{(1+\rho)^{r-1}/2-1\}}s+O(1).
\]
If only $\delta/s\to\rho$ is known, replace $O(1)$ by $o(s^{-1})$.
These rules refer to an unknown collision time and are stress tests,
not proposed source protocols.  A split can even create an adverse pole
on a favorable branch: for
$H=y^4-12uy^2+12u^2$, $u=t-T>0$, at $y=0$,
\[
 E=144u^4,\quad\Lag_0=288u^3>0,\quad
 \frac{P_0}{E}=-\frac{2\delta+\mu\delta^2}{u^2}
                 +\frac{2+2\mu\delta}{u}-\mu
\]
for a fixed $0<\delta<T$, whereas
$P_0/(2E)+Q_0/E=2/u-\mu$.  This is a polynomial control outside
the theta class.

\subsection{Logarithmic losses and local analytic rigidity}
\label{subsec:logarithmic-loss}

For prescribed $M>0$ and finite $q\ge0$, define
\[
 \mathfrak f_q(e)=e[1+\log_+(M/e)]^q\quad(e>0),\qquad
 \mathfrak f_q(0)=0,\qquad\log_+v=\max\{0,\log v\}.
\]
\begin{proposition}[A logarithmic stationary-loss criterion]
\label{prop:logarithmic-collision-criterion}
Suppose on the two stationary sets that
$\Lag_j\ge-K_j\mathfrak f_q(E)$, with nonnegative $K_j$ locally
bounded through every candidate positive-time collision.  Then no such
collision occurs in a domain containing the required predecessor
neighborhoods.  If the domain covers $(0,1/2]\times\R$, RH follows.
\end{proposition}
\begin{proof}
On the predecessor of a zero of multiplicity $2r+j$,
$E\sim k^2s^{2r}$ and $-\Lag_j/E=r/s+O(1)$.  Hence
\[
 \frac{-\Lag_j}{\mathfrak f_q(E)}
 =\frac{r(1+o(1))}{s[2r\log(1/s)]^q}\longrightarrow\infty,
\]
contradicting the local bound for $K_j$.  Nontriviality of the theta
family ensures finite spatial multiplicity.  All divisions are on the
nonzero predecessor, and Proposition~\ref{fc:rh-collision} supplies
the final RH implication.
\end{proof}

A loss by a fixed smaller energy power does not have the same
all-multiplicity implication: $-\Lag_j/E^p$ has order
$s^{2r(1-p)-1}$, which diverges only if $p>1-1/(2r)$.
At equality its limit is $r|k|^{1/r}$, and below equality it tends
to zero.  Also, a loss coefficient integrable in time along the
predecessor would suffice to exclude the event: any coefficient allowing
it must dominate $r/s+O(1)$, whose integral diverges.  This is a
branchwise requirement; two-dimensional local integrability alone gives
no such control on the stationary curve.
Analyticity is essential for the logarithmic statement when $q>1$.
With $\beta=1/(q-1)$, the smooth function
$y(t)=\exp[-(T-t)^{-\beta}]$, continued by zero for $t\ge T$,
satisfies $y'=-\beta y[\log(1/y)]^q$ and reaches zero.  It is not
analytic at $T$ and is not a theta or backward-heat example.

The logarithmic allowance is useful quantitatively, but on analytic
stationary germs it does not enlarge the qualitative class of locally
bounded losses.
\begin{lemma}[Analytic rigidity of logarithmic losses]
\label{lem:analytic-log-rigidity}
Let $S$ be semianalytic near a finite point $p$, and let $E\ge0$ and
$L$ be real analytic in an ambient neighborhood.  Suppose, on $S$ near
$p$,
\[
 L\ge-K E[1+\log_+(M/E)]^q,
\]
where $K\ge0$ is bounded, $M>0$ and $q\ge0$ are fixed and finite,
and the right side is zero at $E=0$.  On a smaller neighborhood in $S$
there is a finite $C\ge0$ such that $L\ge-CE$.  Conversely this
linear bound implies the logarithmic bound with $K=C$.
\end{lemma}
\begin{proof}
If no such $C$ exists, there are $z_n\in S$ tending to $p$ with
$E(z_n)>0$, $-L(z_n)>0$, and $E(z_n)/[-L(z_n)]\to0$.
The semianalytic set
\[
 Z=\{(z,u):z\in S,\ E(z)>0,\ -L(z)>0,\ u>0,
                     \ E(z)=u[-L(z)]\}
\]
therefore has $(p,0)$ in its closure.  Real analytic curve selection
\cite[(7.2), p.~40]{ruiz} supplies an analytic arc
$(z(\rho),u(\rho))$ into $Z$ for $\rho>0$.  Boundedness of $L$
and $E=u(-L)$ imply $E(p)=0$.  Since $u(0)=0$ and
$E(z(0))=0$, but both functions are positive for $\rho>0$, they have
finite orders
$u(\rho)=a\rho^k(1+o(1))$ and
$E(z(\rho))=b\rho^m(1+o(1))$, with $a,b>0$ and integers
$k,m\ge1$.  The assumed estimate would force
\[
 u(\rho)^{-1}\le K[1+\log_+(M/E(z(\rho)))]^q
                  =O((1+|\log\rho|)^q),
\]
contradicting its order $\rho^{-k}$.  At $E=0$ the assumed estimate
already gives $L\ge0$.  The converse is immediate.
\end{proof}

The theta stationary sets are analytic, including their singular points.
Restrictions by $t\le1/2$ or by a closed $t\log(x/(4\pi))$ strip
are locally semianalytic for $x>4\pi$, so the lemma applies there as
well.  It yields an existence constant on a germ, not a source formula
for that constant, nor uniform control at infinity or as $t\downarrow0$.
In particular it does not authorize defining a proof coefficient from
an observed quotient $-\Lag_j/E$.
The exact record
\path{numerics/COLLISION_BEHAVIOR_CHECKS_20261014.json} contains
1072 named assertions over 60 rational jet cases of multiplicities
2 through 13, including 30 delay comparisons.  Its driver
\path{numerics/check_collision_behavior.py} independently forms the
physical numerators and denominators and checks the finite-part
coefficients.  These finite polynomial checks support the algebra;
they do not prove analytic branch existence, the logarithmic limit, or
a theta-specific sign estimate.  Review 20 records the internal audit.

\section{The four jets as quaternionic and octonionic fields}
\label{sec:hypercomplex}

The four scalar jets enter the collision problem in overlapping pairs:
\(H,H_x\) detect a multiple zero, while the two stationary signs involve
\((H,H_x,H_{xx})\) and \((H_x,H_{xx},H_{xxx})\).
Quaternion multiplication keeps both the norms of these pairs and their
signed determinants. A compatible octonionic extension retains these
observations as oriented associators. We give the constructions explicitly
so that their analytic content and their limitations are visible.

\subsection{A quaternionic heat field with independent external time}

Use Hamilton units \(\ii\jj=\kk\),
\(\ii^2=\jj^2=\kk^2=-1\), and identify
\(\C\) with \(\C_{\ii}=\R+\ii\R\).
Complex coefficients are placed on the right of \(\jj\). Thus
\begin{equation}
 (a+\jj b)(c+\jj d)
 =ac-\overline b d+\jj(bc+\overline a d),
 \qquad \overline{a+\jj b}=\overline a-\jj b,
 \label{eq:hc-quaternion-product}
\end{equation}
for \(a,b,c,d\in\C_{\ii}\), and
\(|a+\jj b|^2=|a|^2+|b|^2\).

Write the independent quaternion coordinates as
\[
 q=x_0+\ii x_1+\jj x_2+\kk x_3=z_0+\jj z_1,
 \qquad z_0=x_0+\ii x_1,\quad z_1=x_2-\ii x_3.
\]
The unnormalized left Fueter operator and its conjugate are
\[
 D_q=\partial_{x_0}+\ii\partial_{x_1}
                  +\jj\partial_{x_2}+\kk\partial_{x_3},
 \qquad
 \overline D_q=\partial_{x_0}-\ii\partial_{x_1}
                  -\jj\partial_{x_2}-\kk\partial_{x_3}.
\]
A differentiable field is \emph{left Fueter regular} if \(D_qF=0\).
The convention using the factor \(1/2\) has the same kernel; all units
here multiply on the left. For the classical theory see \cite{sudbery}.

Let \(H(\tau,z)\) be the jointly entire theta heat family, satisfying
\(H_\tau=-H_{zz}\). With \(t\in\R\) an independent external parameter,
set
\[
 \tau=t+\overline{z_1}=t+x_2+\ii x_3.
\]
\begin{equation}
 G_m(t,q)=\partial_z^mH(\tau,z_0)
                +\jj\partial_z^{m+1}H(\tau,z_0),
 \qquad F=G_0.
 \label{eq:hc-lift}
\end{equation}
The coordinate \(z_1\) supplies a complex heat-time displacement; both
entries of \(G_m\) have the same spectral argument.

\begin{proposition}[Regular lift and jet identities]
\label{prop:hc-fueter}
Each \(G_m\) is globally real analytic and left Fueter regular in \(q\).
For every \(m\ge0\),
\begin{align}
 D_qG_m&=0,&\overline D_qG_m&=2G_{m+1},&
 \partial_tG_m&=-G_{m+2}.\label{eq:hc-jet-operators}
\end{align}
\begin{equation}
 \begin{aligned}
 \partial_{x_0}G_m&=G_{m+1},&\qquad
 \partial_{x_1}G_m&=G_{m+1}\ii,\\
 \partial_{x_2}G_m&=-G_{m+2},&
 \partial_{x_3}G_m&=-G_{m+2}\ii.
 \end{aligned}
 \label{eq:hc-coordinate-jets}
\end{equation}
\end{proposition}
\begin{proof}
For complex-valued \(U,V\), direct use of
\eqref{eq:hc-quaternion-product} gives
\[
 \tfrac12D_q(U+\jj V)
 =U_{\overline z_0}-V_{z_1}
              +\jj\bigl(V_{z_0}+U_{\overline z_1}\bigr).
\]
The derivatives are componentwise Wirtinger derivatives. Taking
\(U=\partial_z^mH(\tau,z_0)\) and
\(V=\partial_z^{m+1}H(\tau,z_0)\), the first two terms vanish by
holomorphy and the final pair cancel by the heat equation. Joint
entireness proves global real analyticity. Differentiating the declared
arguments proves \eqref{eq:hc-coordinate-jets} and the external time
identity. Finally \(D_q+\overline D_q=2\partial_{x_0}\) proves the
remaining assertion.
\end{proof}

On the physical slice \(q=x\in\R\), write
\(h_m=h_m(t,x)=\partial_x^mH(t,x)\). The jets are real and
\begin{equation}
 F=h_0+\jj h_1,\qquad F_x=G_1=h_1+\jj h_2,
 \qquad F_{xx}=G_2=h_2+\jj h_3=-F_t.
 \label{eq:hc-physical-jets}
\end{equation}
Thus the four physical jets occur in three consecutive overlapping
quaternion pairs. The basic observations are
\begin{equation}
 E=|F|^2=h_0^2+h_1^2,\qquad
 \mathcal L_j=h_{j+1}^2-h_jh_{j+2}\quad(j=0,1).
 \label{eq:hc-observations}
\end{equation}
In particular \(F(t,x)=0\) means exactly that \(H(t,\cdot)\) has a multiple
real zero at \(x\). The scalar positive-time collision criterion therefore
translates to
\[
 \mathrm{RH}\quad\Longleftrightarrow\quad
 F(t,x)\ne0\quad(t>0,\ x\in\R).
\]
This uses the scalar collision theorem for the theta family. It makes
no assertion about zeros at nonreal quaternion arguments.

The signed products and the heat equation give, on this slice,
\begin{align}
 \overline{G_j}G_{j+1}
   &=h_jh_{j+1}+h_{j+1}h_{j+2}-\jj\mathcal L_j,
   \label{eq:hc-signed-product}\\
 E_t&=-2(h_0h_2+h_1h_3),&
 (\partial_t+\partial_x^2)E&=2(h_1^2+h_2^2).
 \label{eq:hc-energy-evolution}
\end{align}
Consequently
\begin{equation}
 h_1=0\ \Longrightarrow\ E_t=2\mathcal L_0,
 \qquad
 h_2=0\ \Longrightarrow\ E_t=2\mathcal L_1.
 \label{eq:hc-stationary-energy}
\end{equation}
These formulas preserve the two different stationary conditions. The
positive source term in the backward-heat identity
\eqref{eq:hc-energy-evolution} does not give a forward minimum principle.

\subsection{An eight-dimensional monogenic extension}

For a reader unfamiliar with octonions, it suffices here to regard
\(\OO\) as two copies of \(\HH\), with a specified multiplication.
Write \((a,b)=a+\ell b\), where \(\ell=(0,1)\), and adopt the
Cayley--Dickson convention
\begin{equation}
 (a,b)(c,d)=(ac-d\overline b,\overline a d+cb),
 \qquad \overline{(a,b)}=(\overline a,-b).
 \label{eq:hc-octonion-product}
\end{equation}
Then \(|(a,b)|^2=|a|^2+|b|^2\). This is the convention in
\cite[Section 2.2]{baez}. Octonionic multiplication is nonassociative;
we keep parentheses whenever three octonionic factors occur. All
component products on the right of \eqref{eq:hc-octonion-product} are
associative quaternion products.

Let \(r=y_0+\ii y_1+\jj y_2+\kk y_3\), so that
\(p=q+\ell r\) is the eight-dimensional argument. Put
\(e_0=1,e_1=\ii,e_2=\jj,e_3=\kk\), and define the left operator
\begin{equation}
 \mathcal D_8=\sum_{\mu=0}^3e_\mu\partial_{x_\mu}
                    +\sum_{\mu=0}^3(\ell e_\mu)\partial_{y_\mu}.
 \label{eq:hc-D8}
\end{equation}
We call a field \emph{left monogenic} if \(\mathcal D_8\mathcal O=0\).
This fixes the operator convention explicitly; general octonionic
Cauchy--Riemann systems are discussed in \cite{kauhanenorelma}.

\begin{proposition}[Explicit octonionic extension]
\label{prop:hc-octonion-extension}
For each \(m\ge0\), the field
\begin{equation}
 \mathcal O_m(t,q,r)
 =\left(G_m(t,q)-\tfrac12G_{m+3}(t,q)\overline r,
                        G_{m+2}(t,q)\right)
 \label{eq:hc-octonion-lift}
\end{equation}
is globally real analytic and satisfies
\begin{equation}
 \mathcal D_8\mathcal O_m=0,\qquad
 \partial_{x_0}\mathcal O_m=\mathcal O_{m+1},\qquad
 \partial_t\mathcal O_m=-\mathcal O_{m+2},\qquad
 \mathcal O_m(t,q,0)=(G_m,G_{m+2}).
 \label{eq:hc-octonion-identities}
\end{equation}
\end{proposition}
\begin{proof}
Applying \eqref{eq:hc-octonion-product} separately to each term of
\eqref{eq:hc-D8} gives the component formula
\begin{equation}
 \mathcal D_8(A,B)=
 \left(D_qA-\sum_{\mu=0}^3B_{y_\mu}\overline e_\mu,
       \overline D_qB+\sum_{\mu=0}^3A_{y_\mu}e_\mu\right).
 \label{eq:hc-D8-components}
\end{equation}
For \eqref{eq:hc-octonion-lift}, \(B\) is independent of \(r\).
Moreover \(D_q(G_{m+3}\overline r)=(D_qG_{m+3})\overline r=0\),
since the component algebra is associative and \(r\) is held fixed.
The first component of \eqref{eq:hc-D8-components} is therefore zero.
The second is
\[
 2G_{m+3}-\tfrac12\sum_{\mu=0}^3
                         G_{m+3}\overline e_\mu e_\mu
 =2G_{m+3}-2G_{m+3}=0.
\]
The derivative identities follow directly from
\eqref{eq:hc-jet-operators}, and the explicit formula proves global
real analyticity.
\end{proof}

The correction is necessary: the extension \((G_m,G_{m+2})\) constant
in \(r\) has defect \((0,2G_{m+3})\). More generally, two left
Fueter-regular fields \(A(q),B(q)\) admit the extension
\((A-\tfrac12B_{x_0}\overline r,B)\), by the same calculation.
Thus the construction imposes no new condition on the original pair.
For \(m=0\), its value away from \(r=0\) involves \(G_3\), hence
\(h_4\); it is not a closed evolution of only four jets.

On the physical trace \((q,r)=(x,0)\),
\[
 \mathcal O_0=(h_0+\jj h_1)+\ell(h_2+\jj h_3)
 \in\operatorname{span}_{\R}\{1,\jj,\ell,\ell\jj\}.
\]
This span is an associative quaternion subalgebra, as can be checked
from \eqref{eq:hc-octonion-product}. Its total norm is
\begin{equation}
 |\mathcal O_0|^2=E+h_2^2+h_3^2.
 \label{eq:hc-total-norm}
\end{equation}
At a nondegenerate double zero, \(E=0\) but \(h_2\ne0\); the total
octonionic norm is then positive. A collision criterion must retain
the distinguished first pair.

\subsection{Oriented associators preserve the signed observations}

The \emph{associator} measures the failure of associativity:
\[
 [A,B,C]=(AB)C-A(BC),\qquad \mathbf u=\ell\jj.
\]
For embedded quaternions \(a,b\in\HH\), the pair multiplication gives
\begin{equation}
 [a,b,\ell]=\ell(\overline{ab}-\overline{ba})
             =-\ell(ab-ba).
 \label{eq:hc-commutator-associator}
\end{equation}
The last equality holds because a quaternion commutator has zero real
part. The untwisted physical jets lie in the commutative plane
\(\R+\jj\R\), so their associator with \(\ell\) vanishes.
Fixed right multiplication by \(\ii\) and \(\kk\) produces the needed
signed quantities while preserving quaternion left regularity.

\begin{lemma}[Dot products and determinants as associators]
\label{lem:hc-twisted-associators}
For \(a=A+\jj B\), \(b=C+\jj D\) with real \(A,B,C,D\),
\begin{equation}
 [a\ii,b\ii,\ell]=-2\mathbf u(AD-BC),
 \qquad
 [a\ii,b\kk,\ell]=2\mathbf u(AC+BD).
 \label{eq:hc-twisted-associators}
\end{equation}
\end{lemma}
\begin{proof}
One has \(a\ii=A\ii-B\kk\), \(b\ii=C\ii-D\kk\), and
\(b\kk=D\ii+C\kk\). Their quaternion commutators are, respectively,
\(2\jj(AD-BC)\) and \(-2\jj(AC+BD)\). Substitute these into
\eqref{eq:hc-commutator-associator}.
\end{proof}

In particular define, on the physical slice,
\begin{align}
 \mathcal A_E&=[F\ii,F\kk,\ell]=2\mathbf u E,
       \label{eq:hc-associator-energy}\\
 \mathcal A_j&=[G_j\ii,G_{j+1}\ii,\ell]
                    =2\mathbf u\mathcal L_j\quad(j=0,1),
       \label{eq:hc-associator-signs}\\
 \mathcal A_t&=[F\ii,F_t\kk,\ell]=\mathbf u E_t.
       \label{eq:hc-associator-time}
\end{align}
These are immediate from Lemma~\ref{lem:hc-twisted-associators} and
\eqref{eq:hc-energy-evolution}. In contrast to
\eqref{eq:hc-total-norm}, \(\mathcal A_E=0\) is equivalent to \(F=0\).
Equation~\eqref{eq:hc-stationary-energy} becomes
\[
 h_1=0\ \Longrightarrow\ \mathcal A_t=\mathcal A_0,
 \qquad
 h_2=0\ \Longrightarrow\ \mathcal A_t=\mathcal A_1.
\]

There is no order on \(\OO\) being invoked. Specify the oriented ray
\(\mathcal C_{\mathbf u}=\{s\mathbf u:s\ge0\}\).
The two stationary sign conditions are exactly
\begin{equation}
 \mathcal A_t\in\mathcal C_{\mathbf u}
          \quad\hbox{on}\quad\{h_1h_2=0\}.
 \label{eq:hc-oriented-sign}
\end{equation}
Likewise the locally bounded loss condition is exactly
\begin{equation}
 \mathcal A_t+K\mathcal A_E\in\mathcal C_{\mathbf u}
          \quad\hbox{on}\quad\{h_1h_2=0\},
 \qquad K\ge0,
 \label{eq:hc-oriented-loss}
\end{equation}
with \(K\) locally bounded through every candidate collision. Indeed
the coefficient on this ray is \(E_t+2KE\). Thus the associators
preserve both the collision detector and the required orientation.
Their norms alone would replace the signed quantities by absolute values.

\subsection{Why a source-specific estimate remains necessary}

The constructions above apply to every jointly entire backward-heat
family real on the physical slice. The following control shows that
even strict positivity and analytic decay of its cosine-transform source
do not supply the missing sign.

\begin{proposition}[A positive entire source with a positive-time collision]
\label{prop:hc-positive-source-control}
There exist strictly positive, even, entire functions \(\Phi_*\),
decaying superexponentially on the real axis, whose heat cosine
transforms are jointly entire and have a nondegenerate real double zero
at a positive time. Their quaternionic and octonionic lifts satisfy all
the preceding identities.
\end{proposition}
\begin{proof}
Set \(\rho(v)=2e^{-v^4}/\Gamma(1/4)\), so that
\(\int_{\R}\rho(v)\,dv=1\), and
\(\rho_\epsilon(v)=\epsilon^{-1}\rho(v/\epsilon)\) for real
\(\epsilon>0\). Use
\[
 \Phi_\epsilon(u)=
 2\{\rho_\epsilon(u-1)+\rho_\epsilon(u+1)\}
   +\rho_\epsilon(u-2)+\rho_\epsilon(u+2).
\]
This source is strictly positive, even, and entire, with quartic decay
on the real axis. Its half-line heat cosine transform is
\begin{align*}
 \widetilde H_\epsilon(t,z)
 &=2\int_{\R}\rho(v)e^{t(1+\epsilon v)^2}
                         \cos\bigl(z(1+\epsilon v)\bigr)\,dv\\
 &\quad+\int_{\R}\rho(v)e^{t(2+\epsilon v)^2}
                         \cos\bigl(z(2+\epsilon v)\bigr)\,dv.
\end{align*}
Quartic decay supplies an integrable majorant on every compact set of
complex \((t,z,\epsilon)\). Hence this formula extends jointly
entirely through \(\epsilon=0\), obeys backward heat, and there becomes
\[
 \widetilde H_0(t,z)=2e^t\cos z+e^{4t}\cos(2z).
\]
At \(T_0=(\log2)/3\), \(x_0=\pi\), it satisfies
\(\widetilde H_0=(\widetilde H_0)_x=0\) and
\((\widetilde H_0)_{xx}=-6e^{T_0}\ne0\).
By the heat equation the Jacobian of
\((\widetilde H_\epsilon,(\widetilde H_\epsilon)_x)\)
with respect to \((t,x)\), at this point and \(\epsilon=0\), has
determinant \(-((\widetilde H_0)_{xx})^2\ne0\).
The real analytic implicit-function
theorem therefore gives a nearby nondegenerate real double zero at a
positive time for every sufficiently small \(\epsilon>0\).
\end{proof}

The failure is visible in the signed observation. At such a collision
\((T,a)\), put \(c=H_{xx}(T,a)\ne0\). The implicit-function theorem
gives a real stationary branch \(H_x(t,\chi(t))=0\) through it.
For \(\delta=T-t\downarrow0\), backward heat gives
\[
 H(t,\chi(t))=c\delta+O(\delta^2),\qquad
 E(t,\chi(t))=c^2\delta^2+O(\delta^3),\qquad
 \mathcal L_0(t,\chi(t))=-c^2\delta+O(\delta^2).
\]
Thus \eqref{eq:hc-oriented-loss} would require
\(K\ge\delta^{-1}+O(1)\), contradicting local boundedness at the
collision. The lift and associator identities alone cannot provide that
bound. For the theta family, its particular arithmetic structure must
enter a further signed estimate.

For the genuine theta family, Theorem~\ref{thm:outer} supplies the
additional sign information in its declared outer region:
$\mathcal A_0$ and $\mathcal A_1$ lie strictly on
$\mathcal C_{\mathbf u}$ there. At genuine stationary points the same
is true of $\mathcal A_t$. This is the geometric consequence of the
common arithmetic amplitude, not a consequence of monogenicity alone.


\section{Centered residuals and the common theta carrier}\label{sec:residuals}

The outer theorem uses oscillatory cancellation to reach a nonzero zeta
comparison amplitude. To study the remaining region, we return to the
exact theta source.
The aim is to retain the same four derivatives and their signed
correlations before estimating any error.

\subsection{A contour identity for the physical derivatives}
\begin{proposition}[Complete-source contour shift]\label{prop:complete-shift}
For real $t\ge0$, $x\ge0$, $0<a<\pi/8$, and every integer $k\ge0$,
\begin{equation}
 h_k(t,x)=e^{-ax}\Re\int_0^\infty
 (i(u+ia))^k e^{t(u+ia)^2}\Phi(u+ia)e^{ixu}\dd u.
 \label{eq:complete-shift}
\end{equation}
\end{proposition}
\begin{proof}
The theta series is normally convergent in the strip
$|\Im z|<\pi/8$ because $\Re(e^{4z})>0$ there. Analytic continuation
of the real reflection identity makes $\Phi$ even throughout the
strip; Schwarz symmetry then makes $\Phi(iy)$ real. On the left edge
of the contour shift, every factor of
$(iz)^ke^{tz^2}\Phi(z)e^{ixz}$ is real at $z=iy$, while $\dd z=i\dd y$.
Its contribution has zero real part. On the right edge,
$\cos(4y)\ge\cos(4a)>0$ gives uniform double-exponential decay.
Cauchy's theorem proves the formula. The evenness argument applies
to the complete source; it must precede any separation into individual
theta summands.
\end{proof}

The all-index and integration bounds in Appendix~\ref{app:payments}
make truncation of \eqref{eq:complete-shift} rigorous. They control
absolute errors at the gamma exponential scale. Obtaining a sign still
requires information about the retained, correlated jets.

\subsection{A residual adapted to genuine stationarity}
Freeze real $a,b$ with $0<a<\pi/8$ and $b\ge0$, and put $\mu=a^2+b^2$.
For every integer $j\ge0$, define
\begin{equation}
 C_j=e^{-ax}\Re\int_0^\infty (u^2-b^2)(i(u+ia))^j
             e^{t(u+ia)^2}\Phi(u+ia)e^{ixu}\dd u.
 \label{eq:C-integral}
\end{equation}
The required stationary forms use $j=0,1$; higher residuals will enter
exact heat-transport identities. The polynomial identity
$u^2-b^2=-(i(u+ia))^2-2a\,i(u+ia)-\mu$ gives
\begin{equation}
 C_j=-h_{j+2}-2ah_{j+1}-\mu h_j,
 \qquad h_{j+1}=0\ \Longrightarrow\
 \Lag_j=\mu h_j^2+h_jC_j.
 \label{eq:C-stationary}
\end{equation}
Thus the contour drift disappears at exactly the original stationary
constraint. No derivative of an exponentially normalized observation
has replaced that constraint.

\begin{proposition}[A sufficient residual cone]\label{prop:residual-cone}
If $\eta>0$ and
\begin{equation}
 \Delta:=\eta^2E-C_0^2-C_1^2\ge0,
 \label{eq:cone-slack}
\end{equation}
then the respective stationary sets satisfy
\begin{equation}
 \Lag_0\ge(\mu-\eta)E\quad(h_1=0),\qquad
 \Lag_1\ge-\frac{\eta^2}{4\mu}E\quad(h_2=0).
 \label{eq:cone-losses}
\end{equation}
A sufficient common negative-loss coefficient is $K=\eta^2/(4\mu)$.
It is locally bounded wherever the prescribed parameters are locally
bounded and $\mu$ stays positive.
\end{proposition}
\begin{proof}
At $h_1=0$, $E=h_0^2$ and $|C_0|\le\eta\sqrt E$, so the first
inequality follows from \eqref{eq:C-stationary}. At $h_2=0$, complete
the square in $\mu h_1^2-\eta|h_1|\sqrt E$ to obtain the second.
Finally $\eta-\mu\le\eta^2/(4\mu)$ because
$(\eta-2\mu)^2\ge0$. No division by $E$ occurs.
\end{proof}

For use with exact algebra or interval certificates, the same argument
has the following identities before stationarity is imposed:
\begin{align}
 2\eta[\Lag_0+(\eta-\mu)E]
 &= (C_0+\eta h_0)^2+C_1^2+\Delta\notag\\
 &\quad+2\eta h_1[(1-\mu+\eta/2)h_1+2ah_0],
 \label{eq:cone-SOS0}\\
 4\mu\Lag_1+\eta^2E
 &= (2\mu h_1+C_1)^2+C_0^2+\Delta
             +4\mu h_2(h_2+2ah_1).
 \label{eq:cone-SOS1}
\end{align}
Each final term is a multiple of the genuine stationary derivative.
The cone is sufficient, but a one-sided bound for $h_jC_j$ can be
substantially weaker and therefore more useful.

For example, suppose paid normalized approximations give
$Z_j=e^{\pi x/8}h_j=p+\delta$ and
$e^{\pi x/8}C_j=q+\epsilon'$, where
$|\delta|\le\varepsilon_0$, $|\epsilon'|\le\zeta_0$.
At $h_{j+1}=0$, expansion of \eqref{eq:C-stationary}, retaining
$\mu\delta^2\ge0$, gives
\begin{equation}
 e^{\pi x/4}\Lag_j\ge\mu p^2+pq
 -|2\mu p+q|\varepsilon_0-|p|\zeta_0-\varepsilon_0\zeta_0.
 \label{eq:paid-signed-residual}
\end{equation}
Every truncation, quadrature, and parameter error must be included in
the two bounds. The retained product $pq$ preserves its sign.

\subsection{One carrier for all four derivatives}
There is an exact arithmetic reduction behind the centered residual.
Put $d_n=\tfrac12\log n$. The source summands obey
\begin{equation}
 \phi_n(z)=n^{-1/2}\phi_1(z+d_n).
 \label{eq:theta-translation}
\end{equation}
For $v\ge d_n$, define
\begin{equation}
 A_n(t,x;v)=n^{-1/2}
       \exp\{td_n^2-2td_n(v+ia)-ixd_n\}.
 \label{eq:An}
\end{equation}
Changing variables $v=u+d_n$ in the complete shifted source yields
\begin{equation}
 h_j=e^{-ax}\Re\int_0^\infty
 e^{t(v+ia)^2}\phi_1(v+ia)e^{ixv}
 \sum_{n\le e^{2v}}A_n(t,x;v)[i(v+ia-d_n)]^j\dd v.
 \label{eq:shared-carrier}
\end{equation}
In particular $j=0,1,2,3$ uses only the first four logarithmic moments
of the same arithmetic weights. The staircase cutoff $n\le e^{2v}$
is part of the identity. It is not an independently chosen truncation.

To justify the rearrangement, let $0\le t\le\tau$. Since
$0\le d_n\le v$, we have $|A_n|\le n^{-1/2}$ and
$\sum_{n\le e^{2v}}n^{-1/2}\le2e^v$. Therefore the absolute integrand
in \eqref{eq:shared-carrier}, before $e^{-ax}$, is bounded by
\begin{equation}
 2(2\pi^2+3\pi)(v+a)^j
           e^{\tau v^2+10v-\pi\cos(4a)e^{4v}}.
 \label{eq:shared-majorant}
\end{equation}
This is integrable. It also gives a full integration tail, stated in
Appendix~\ref{app:payments}; no unretained index is omitted from that
bound.

The common oscillatory exponential in the original $n$-th channel has
phase $(x+2ta)u-\pi n^2\sin(4a)e^{4u}$. Its critical location is
\begin{equation}
 b_n=\frac14\log\frac{x+2ta}{4\pi n^2\sin(4a)}
     =b_1-d_n.
 \label{eq:channel-centers}
\end{equation}
Only nonnegative locations belong to the integration half-line.
This observation supplies a centering choice; it does not assert
first-channel dominance or a full stationary-phase approximation.
Taking $b=b_1$, the multiplier defining $C_j$ becomes
\begin{equation}
 (v-d_n)^2-b^2
 =\underbrace{-b\log n+\tfrac14\log^2n}_{\delta_n}
       +2(b-d_n)(v-b)+(v-b)^2.
 \label{eq:logarithmic-residual}
\end{equation}
Thus the remaining signed residual involves logarithmic weights of the
same arithmetic sum, rather than unrelated errors for each derivative.
At $v=b$ those weights have the transparent form
\begin{equation}
 A_n=n^{-1/2-tb-i(x/2+ta)}e^{t\log^2n/4}.
 \label{eq:saddle-weights}
\end{equation}
The factors $A_n$ are complex and their products retain arithmetic
phases. A positive-weight covariance inequality does not apply to them.

For perspective, if $b>0$ the frozen-parameter balanced quaternion
$W=h_0+\jj(h_1+ah_0)/b$ obeys
\begin{equation}
 W_t=(b^2-a^2-2ab\jj)W+R,\qquad
 R=C_0+\jj(C_1-aC_0)/b.
 \label{eq:balanced-residual-evolution}
\end{equation}
The skew term rotates $W$ without changing its norm. Equation
\eqref{eq:balanced-residual-evolution} separates the expected oscillatory
motion from the residual, whereas \eqref{eq:C-stationary} keeps the
precise scalar conditions needed for a sign. If $a,b$ vary with time,
the additional term is
$\jj(a_t\Re W/b-b_t\Im_{\jj}W/b)$; all cell computations below freeze
them before differentiating.


\subsection{Staircase derivatives and exact entry terms}
Put $T=x/2+ta$ and, for integers $r\ge0$,
\[
 B_r(v)=\sum_{n\le e^{2v}}(\log n)^r A_n(t,x;v),
 \qquad
 D_j(v)=\sum_{n\le e^{2v}}A_n(t,x;v)[i(v+ia-d_n)]^j.
\]
For $r=0$ the logarithmic weight is one, including at $n=1$.
On $v>0$, distributional differentiation gives
\begin{equation}
 B_r'=-tB_{r+1}+
 \sum_{n\ge2}(\log n)^r A_n(t,x;d_n)\delta(v-d_n).
 \label{eq:staircase-recurrence}
\end{equation}
Between entry points,
\begin{equation}
 D_j'=-2t(v+ia)D_j-2itD_{j+1}+ijD_{j-1},
 \label{eq:staircase-jet-recurrence}
\end{equation}
where the last term is omitted for $j=0$. The jump at $d_n$ is
$A_n(t,x;d_n)(-a)^j$. If
$K(v)=e^{t(v+ia)^2}\phi_1(v+ia)e^{ixv}$, its full contribution is
\begin{equation}
 K(d_n)A_n(t,x;d_n)(-a)^j
   =e^{-ta^2}\phi_n(ia)(-a)^j.
 \label{eq:staircase-cancellation}
\end{equation}
Indeed $\partial_v A_n=-2td_nA_n$ and
$d_n=v+ia+i[i(v+ia-d_n)]$ give the smooth recurrence; direct
substitution and \eqref{eq:theta-translation} give the jump formula.
At $v=0$ the index $n=1$ is already present and belongs to the lower
endpoint of any integration by parts. Every entering-index phase in
\eqref{eq:staircase-cancellation} cancels in $x$. These terms cannot
be assigned an artificial oscillatory gain.

The oscillatory block estimate in Section~\ref{sec:outer-proof} can also
be applied to this exact carrier. On a block $M<n\le R(v)$, where
$R(v)=\min(2M,e^{2v})$, the logarithmic-moment weight is
\[
 w_r(y)=(\log y)^r e^{t\log^2y/4}y^{-1/2-tv}.
\]
For $M\ge e^{2r}$ it is decreasing, since
$y(\log w_r)'\le r/\log y-1/2\le0$; for $r=0$ this holds for
all $M\ge1$. Abel summation therefore gives, when $|T|\ge1$,
\[
 \left|\sum_{M<n\le R(v)}w_r(n)n^{-iT}\right|
 \ll w_r(M)\bigl(|T|^{1/6}M^{1/2}+|T|^{-1/6}M\bigr).
\]
For other blocks use the full total variation of $w_r$, and for
$|T|<1$ use the trivial bound. The actual shortened endpoint remains
$R(v)$; these complex block estimates do not imply a signed product
bound.

\subsection{The complete paired observations}
Define the staircase measure space and its coordinates by
\begin{equation}
 \begin{split}
 \Omega&=\{(v,n):n\ge1,\ v\ge d_n\},\qquad
 d\nu=dv\times\text{counting measure},\\
 y_\alpha&=v-d_n,\qquad q_\alpha=-a+iy_\alpha,\\
 W_\alpha&=e^{t(v+ia)^2}\phi_1(v+ia)e^{ixv}A_n(t,x;v),
 \qquad \alpha=(v,n).
 \end{split}
 \label{eq:pair-domain}
\end{equation}
Then $h_k=e^{-ax}\Re\int_\Omega W_\alpha q_\alpha^k\,d\nu(\alpha)$.
The two factors in the following integrals have independent coordinates
and independent indices.
\begin{proposition}[Paired Laguerre and residual forms]
For $j=0,1$,
\begin{align}
 \Lag_j=\frac{e^{-2ax}}4\Re\iint_{\Omega^2}\bigl[&
 (y_\alpha+y_\beta)^2 W_\alpha\overline{W_\beta}
       (q_\alpha\overline{q_\beta})^j\notag\\
 &+(y_\alpha-y_\beta)^2 W_\alpha W_\beta
       (q_\alpha q_\beta)^j\bigr]\,d\nu(\alpha)d\nu(\beta).
 \label{eq:paired-laguerre}
\end{align}
Writing
\begin{equation}
 J_j(\alpha,\beta)=W_\alpha\overline{W_\beta}
       (q_\alpha\overline{q_\beta})^j
       +W_\alpha W_\beta(q_\alpha q_\beta)^j,
 \label{eq:paired-J}
\end{equation}
the centered product is
\begin{equation}
 h_jC_j=\frac{e^{-2ax}}4\Re\iint_{\Omega^2}
 (y_\alpha^2+y_\beta^2-2b^2)J_j(\alpha,\beta)\,
 d\nu(\alpha)d\nu(\beta).
 \label{eq:paired-residual}
\end{equation}
Before imposing stationarity, the exact difference is
\begin{equation}
 \Lag_j-\mu h_j^2-h_jC_j
       =h_{j+1}(h_{j+1}+2ah_j).
 \label{eq:stationary-ideal}
\end{equation}
\end{proposition}
\begin{proof}
Use $\Re A\Re B=\tfrac12\Re(A\overline B+AB)$ and symmetrize
in $\alpha,\beta$. The Hermitian difference in the Laguerre form is
\[
 q_\alpha\overline{q_\beta}
 -\tfrac12(q_\alpha^2+\overline{q_\beta}^{2})
 =-\tfrac12(q_\alpha-\overline{q_\beta})^2
 =\tfrac12(y_\alpha+y_\beta)^2;
\]
the bilinear difference is
$-\tfrac12(q_\alpha-q_\beta)^2
=\tfrac12(y_\alpha-y_\beta)^2$.
Conjugate-swapping the Hermitian term preserves its real integral.
The same polarization applied to the real multiplier $y^2-b^2$
gives \eqref{eq:paired-residual}, and \eqref{eq:C-stationary} gives
\eqref{eq:stationary-ideal}. The majorant
\eqref{eq:shared-majorant} has every polynomial moment integrable.
Together with $(y\pm z)^2\le2(y^2+z^2)$, this proves absolute
convergence of each displayed integral and justifies Fubini.
\end{proof}

The squared frequencies in \eqref{eq:paired-laguerre} do not give the
phase of either amplitude product a sign. As a finite algebra control,
two real signed atoms at frequencies $1,2$, with weights $1,-1/2$
and $a=0$, give $(h_0,h_1,h_2)=(1/2,0,1)$ and $\Lag_0=-1/2$.
The diagonal contributions sum to $2$ and the ordered cross terms to
$-5/2$. This control lies outside the theta class and positive-contour
hypotheses; it explains why retaining only diagonal pairs loses the
observation.

\subsection{A complete omitted-pair payment}
For a measurable retained set $A\subset\Omega$, denote its jets by
$h_k^A$ and set $\Lag_j^A=(h_{j+1}^A)^2-h_j^Ah_{j+2}^A$. Define
absolute masses before the exterior exponential normalization:
\[
 M_k^A=\int_A|W_\alpha||q_\alpha|^k\,d\nu,
 \qquad T_k^A=\int_{\Omega\setminus A}|W_\alpha||q_\alpha|^k\,d\nu.
\]
Then
\begin{equation}
 |\Lag_j-\Lag_j^A|\le e^{-2ax}
 \bigl(M_j^AT_{j+2}^A+M_{j+2}^AT_j^A+T_j^AT_{j+2}^A\bigr).
 \label{eq:pair-tail-payment}
\end{equation}
Indeed, including the factor $1/4$ in \eqref{eq:paired-laguerre},
the absolute sum of its two kernels is at most
\[
 \tfrac12\bigl(|W_\alpha||q_\alpha|^{j+2}
              |W_\beta||q_\beta|^j
       +|W_\alpha||q_\alpha|^j
              |W_\beta||q_\beta|^{j+2}\bigr),
\]
using $(y+z)^2+(y-z)^2=2(y^2+z^2)$ and $y^2\le|q|^2$.
Integrating over $A\times A^c$, $A^c\times A$, and $A^c\times A^c$
gives the claim, including all omitted mixed and omitted--omitted pairs.
No bound for $T_{j+1}^A$ is required.

For $A=\{(v,n)\in\Omega:v\le V,\ n\le N\}$, use the complete
index payment and the shared-coordinate integration payment in
Appendix~\ref{app:payments}, removing their $e^{-ax}$ normalization
before using them for $T_k^A$. Overlap may be overpaid. Retained
quadrature and parameter variation require their own errors. The
absolute payment \eqref{eq:pair-tail-payment} is not a relative bound
through an unknown zero of $E$.

\subsection{Signed staircase cutoffs on the closed outer domain}
Fix $a=\pi/16$ and retain exactly
\begin{equation}
 \Omega_N=\{(v,n)\in\Omega:0\le v\le V_N\},\qquad
 V_N=\tfrac12\log N,
 \label{eq:signed-staircase-domain}
\end{equation}
where the intrinsic restriction $n\le e^{2v}$ remains part of
$\Omega$. There is no additional rectangular approximation to that
staircase. Write $h_k^{(N)}$ for the restricted common-source moments
and $\Lag_j^{(N)}=(h_{j+1}^{(N)})^2-h_j^{(N)}h_{j+2}^{(N)}$.
All cross terms in these products are retained.

To avoid confusing source constants with the centered residual $C_0$,
use the notation
\begin{equation}
 \begin{gathered}
 \beta_0=\pi\cos(4a)=\frac\pi{\sqrt2},\qquad
 \mathsf C_0=2(2\pi^2+3\pi),\qquad
 c=\beta_0/2,\qquad \mathsf D_0=\frac{\mathsf C_0}{2\beta_0},\\
 \mathsf M_k=\mathsf C_0\int_0^\infty(v+a)^k
   \exp\{v^2/2+10v-\beta_0e^{4v}\}\dd v,\qquad 0\le k\le3,\\
 B_j=(\mathsf M_j+\mathsf M_{j+2})\mathsf D_0+\mathsf D_0^2,
 \qquad j=0,1.
 \end{gathered}
 \label{eq:signed-cutoff-constants}
\end{equation}
These finite constants depend only on the fixed complete source and
$a$, not on $t,x,N$.

\begin{corollary}[Signed staircase cutoff estimates]
\label{cor:signed-staircase-cutoffs}
For all integers $N\ge8$, $j=0,1$, and $(t,x)$ in
\eqref{eq:outer-domain},
\begin{equation}
 \Lag_j^{(N)}\ge-A_j^*(t,x)E,
 \qquad
 \Lag_j^{(N)}\ge-e^{-2ax}B_j e^{-cN^2},
 \label{eq:signed-cutoff-conclusions}
\end{equation}
where
\begin{equation}
 A_j^*(t,x)=
 \frac{e^{-2ax}\mathsf M_j\mathsf M_{j+2}\zeta(9/8)^2}
      {\Gamma(t,x)^2},\qquad
 \Gamma(t,x)=|M_t((1+ix)/2)|.
 \label{eq:signed-cutoff-relative-coefficient}
\end{equation}
The coefficient is prescribed independently of $N$ and of the measured
physical energy; it is locally bounded but is not asserted uniformly
bounded as $x\to\infty$. The second estimate has zero energy loss and
height-uniform Gaussian tail constants. The observations and stationary
sets remain those of the full function.
\end{corollary}
\begin{proof}
For $0\le t\le1/2$, \eqref{eq:shared-majorant} bounds every retained
absolute $k$-th moment by $\mathsf M_k$. Summing the absolute values of
the Hermitian and bilinear kernels in \eqref{eq:paired-laguerre}, using
$(y+z)^2+(y-z)^2=2(y^2+z^2)$ and its factor $1/4$, gives
\begin{equation}
 |\Lag_j^{(N)}|\le
 e^{-2ax}\mathsf M_j\mathsf M_{j+2}.
 \label{eq:signed-cutoff-absolute}
\end{equation}
The energy envelope \eqref{eq:closed-sign-envelopes} and
\eqref{eq:signed-cutoff-relative-coefficient} therefore prove the first
inequality in \eqref{eq:signed-cutoff-conclusions}, with a spare factor
two. Its denominator is the explicitly nonzero analytic carrier, not
$E$; this proves the asserted independence and local boundedness.

For the second inequality, the entire omitted source region is
$v>V_N$. For $v\ge1$,
\[
 \log(v+a)\le v,\qquad
 v^2/2+13v\le(\beta_0/2)e^{4v}.
\]
The latter follows because $e^{-4v}(v^2/2+13v)$ is decreasing there
and its value at one is less than $\beta_0/2$. Since $V_N\ge1$ for
$N\ge8$, every omitted moment of order $k\le3$ satisfies
\begin{equation}
 T_k^{(N)}\le
 \frac{\mathsf C_0}{2\beta_0N^2}e^{-cN^2}
 \le\mathsf D_0e^{-cN^2}.
 \label{eq:signed-cutoff-tail-moments}
\end{equation}
Indeed substitute $q=e^{4v}$ in the integral of $e^{-ce^{4v}}$ and
use $q^{-1}\le N^{-2}$. Apply the complete payment
\eqref{eq:pair-tail-payment}, including both mixed sectors and the
omitted--omitted sector. With the constants in
\eqref{eq:signed-cutoff-constants} it gives
\begin{equation}
 |\Lag_j-\Lag_j^{(N)}|
 \le e^{-2ax}B_j e^{-cN^2}.
 \label{eq:signed-cutoff-complete-tail}
\end{equation}
The independently proved full sign $\Lag_j>0$ establishes the second
inequality in \eqref{eq:signed-cutoff-conclusions}.
\end{proof}

The absolute and tail estimates in the proof hold for every $x\ge0$
and $0\le t\le1/2$; the signed conclusions use the independent outer
positivity theorem and are asserted only in its domain. They do not
assert positivity of every truncated form. Also $h_{j+1}=0$ does not
imply $h_{j+1}^{(N)}=0$; no such substitution is made. At $E=0$ on the corresponding full stationary set, an unpaid inequality $\Lag_j^{(N)}\ge-AE$ would require
$\Lag_j^{(N)}\ge0$ for every cutoff, which does not follow merely from
$\Lag_j=0$. The Gaussian additive allowance avoids imposing that extra
cutoff condition.

One exact alternative retaining the complete amplitude is
\[
 \mathcal J_j^{(N)}=-h_jh_{j+2}^{(N)}
 \quad\hbox{on }h_{j+1}=0,\qquad j=0,1.
\]
It obeys
\begin{equation}
 \mathcal J_j^{(N)}-\Lag_j=h_j(h_{j+2}-h_{j+2}^{(N)}),\qquad
 |\mathcal J_j^{(N)}-\Lag_j|
       \le\sqrt E\,e^{-ax}T_{j+2}^{(N)}.
 \label{eq:full-amplitude-cutoff-error}
\end{equation}
Thus $\mathcal J_j^{(N)}$ vanishes at $E=0$ without imposing positivity
of a raw truncated quadratic form. This is an exact error identity,
not a new signed estimate in the remaining cusp.

\section{The bounded region: complete certificates through height 112}
\label{sec:fc-certificate}

The bounded calculation supplies a genuine norm gap, so it excludes
collisions directly.  Its role is to join the central analytic estimates
to the large-height theorem, while leaving the intervening and
small-$t\log|x|$ regions explicitly visible.

\begin{theorem}[Certified closed bounded region]
\label{fc:height104}
For the complete theta source,
\begin{equation}
 |F(t,x)|=\sqrt{h_0(t,x)^2+h_1(t,x)^2}>10^{-16}
 \quad(0\le t\le1/2,\ |x|\le104).
 \label{fc:height104-gap}
\end{equation}
On this rectangle, $h_1=0$ implies $\mathcal L_0>0$, and $h_2=0$
implies $\mathcal L_1>0$.  The norm bound combines an analytic central
estimate with the validated interval records described below.
\end{theorem}

\subsection{A central estimate with no theta truncation}

\begin{proposition}[Central norm gap]
\label{fc:central}
For $0\le t\le1$ and $|x|\le8$,
$|F(t,x)|>1/120000$.  Both stationary signs in
Theorem~\ref{fc:height104} hold on this larger time interval.
\end{proposition}
\begin{proof}
Put $w_t(u)=e^{tu^2}\Phi(u)$ and
$M_k(t)=\int_0^\infty u^kw_t(u)\,du$.
The complete source obeys
$0<\Phi(u)<4\pi^2e^{9u-\pi e^{4u}}$ for $u\ge0$:
after dropping each summand's negative term, the ratio of successive
$n^4e^{-\pi e^{4u}n^2}$ terms is at most $16e^{-3\pi}<1/2$.
The following elementary lower cores and upper tails hold uniformly
for $0\le t\le1$:
\begin{equation}
 \begin{array}{c|c@{\qquad}c|c}
 \text{core }[0,d]&\Phi(u)>c_d&\text{tail }u\ge b&w_t(u)<A_be^{-q_b(u-b)}\\ \hline
 d=1/10&6/25&b=7/25&(1/20)e^{-28(u-b)}\\
 d=3/20&3/25&b=3/10&(1/49)e^{-32(u-b)}\\
 d=1/5&1/20&&
 \end{array}
 \label{fc:source-envelopes}
\end{equation}
For completeness, the lower bounds follow from
$\phi_1(u)=e^ug(\pi e^{4u})$ with
$g(q)=q(2q-3)e^{-q}$ and $g'(q)=-(2q-1)(q-3)e^{-q}$.
On the three cores the respective arguments are less than
$33/7,6,36/5$.  The bounds
$e<11/4$, $e^{2/7}>9/7$, and $e^{1/5}<5/4$ give the three stated
rational lower bounds for $g$.
For the tails, $\psi_t(u)=9u+tu^2-\pi e^{4u}$ is strictly concave.
At $b=7/25$ one has $\psi_t(b)<-69/10$, $\psi_t'(b)<-28$;
at $b=3/10$ one has $\psi_t(b)<-38/5$, $\psi_t'(b)<-32$.
Together with $4\pi^2<40$, $e^{69/10}>800$, and $e^{38/5}>1960$,
these prove the full upper tails in \eqref{fc:source-envelopes}.
All comparisons use rational Taylor bounds for the indicated
exponentials; in particular they do not truncate the theta source.

Integrating $u^2(1-25u^2/2)$ over the first positive core and paying
its entire negative tail with $b=7/25$ gives
\begin{align}
 M_2-\frac{25}{2}M_4
 &>\frac{37}{500000}
   -\frac58\frac{3497259}{33614000000}
   =\frac{2413193}{268912000000}>\frac1{120000}.
 \label{fc:moment24}
\end{align}
Indeed the polynomial is nonnegative up to $b$ and is bounded below
on the tail by $-(25/2)u^2(u^2-b^2)$.
Likewise, with $f(u)=u^4(1-32u^2/3)$ and $b=3/10$, the three
successive cores and the entire negative tail give
\begin{align}
 M_4-\frac{32}{3}M_6
 &>\frac{111799}{33600000000}
   -\frac{32}{147}\frac{85837761}{6710886400000}
   =\frac{2092283587}{3853516800000000}>\frac1{2000000}.
 \label{fc:moment46}
\end{align}
The positive term is
$(6/25)\int_0^{1/10}f+(3/25)\int_{1/10}^{3/20}f
 +(1/20)\int_{3/20}^{1/5}f$;
the tail integral is
$\int_b^\infty u^4(u^2-b^2)e^{-32(u-b)}\,du$.
Thus both displayed differences are exact rational integrations.

The global inequalities $\sin y\ge y-y^3/6$ for $y\ge0$ and
$\cos y\ge1-y^2/2$ now imply, for $0<x\le8$,
\[
 -h_1\ge x(M_2-x^2M_4/6)>x/120000,
 \qquad h_3\ge x(M_4-x^2M_6/6)>x/2000000.
\]
For $|x|\le1$, integrate the first core using
$\cos(xu)\ge199/200$, discard the nonnegative middle interval
$[1/10,7/25]$, and pay the entire remaining tail.  This gives
\[
 H(t,x)>\frac{597}{25000}-\frac1{560}
 =\frac{7733}{350000}>\frac1{50}.
\]
The value bound covers $|x|\le1$, and the slope bound covers
$1\le|x|\le8$, proving the norm gap by evenness.
The first stationary set is exactly $x=0$, where
$\mathcal L_0=H(t,0)M_2>0$.
For $0<|x|\le8$,
$\mathcal L_1=h_2^2-h_1h_3>x^2/240000000000$;
at $x=0$ one has $h_2=-M_2<0$, so the second stationary set is absent.
\end{proof}

\subsection{What the interval certificates enclose}

The records use Arb/Acb interval arithmetic to enclose the complete
physical jets on closed space--time cells.  Each enclosure includes
the returned finite-quadrature radius, every omitted theta index,
the complete integration tail, rounding, and spatial and temporal
Taylor remainders.  The older compensated calculation also pays the
complete Gaussian transforms and their truncation remainders.
Finite compensated channels are transported using the raw heat
equation followed by their explicit time-dependent compensation;
they are not assigned a homogeneous heat equation.

For a heat-transported sequence $Y_r$, a cell with half-widths
$d_t,d_x$ uses
\[
 Y_r(t_c+\tau,x_c+\eta)=
 \sum_{\ell=0}^{q-1}\sum_{m=0}^{p-1}
 \frac{(-\tau)^\ell\eta^m}{\ell!m!}Y_{r+2\ell+m}(t_c,x_c)
 +\mathcal R_r.
\]
If $M_k$ bounds the complete $k$th derivative on the cell, Taylor's
theorem first in time and then in space gives
\begin{equation}
 |\mathcal R_r|\le
 \sum_{\ell=0}^{q-1}\frac{d_t^\ell}{\ell!}
 \frac{d_x^p}{p!}M_{r+2\ell+p}
 +\frac{d_t^q}{q!}M_{r+2q}.
 \label{fc:taylor-payment}
\end{equation}
The higher bands sum the retained theta coefficients before applying
the common Taylor monomials.  This preserves cancellation between
different indices while adding every positive remainder payment.

Stationary decisions concern genuine derivatives.  If
$h_k=p_k+r_k$ with $|r_k|\le\varepsilon_k$, then at $h_{j+1}=0$
one has
\begin{equation}
 \mathcal L_j\ge p_{j+1}^2-p_jp_{j+2}
 -\bigl(\varepsilon_{j+1}^2+|p_j|\varepsilon_{j+2}
       +|p_{j+2}|\varepsilon_j+\varepsilon_j\varepsilon_{j+2}\bigr).
 \label{fc:stationary-payment}
\end{equation}
Here $p_{j+1}^2\le\varepsilon_{j+1}^2$ follows from genuine
stationarity; the square payment has a negative sign.
The products retain all cross terms between retained channels.
A cell is discharged either by exclusion of the stationary derivative
or by a strictly positive paid signed lower bound.
Monotone endpoint exclusions, where used, employ the next physical
derivative and uniform endpoint signs.  No zero of an ordinary
derivative of an exponentially rescaled function is substituted.

Independently, if $I_0,I_1$ enclose the genuine $h_0,h_1$ throughout
a cell, then
\begin{equation}
 E\ge \operatorname{dist}(0,I_0)^2+\operatorname{dist}(0,I_1)^2.
 \label{fc:interval-norm}
\end{equation}
Exact rational comparisons and downward square-root rounding turn
these energy floors into the reported norm floors.  No division by
an unproved norm is used.

\subsection{Coverage, the centered-residual pilot, and reproducibility}

Every numerical band covers the full closed time interval $[0,1/2]$
using eight slabs $[r/16,(r+1)/16]$.  The records give the following
conservative norm bounds; the central row is analytic.
\begin{center}
\begin{tabular}{c r c c}
Positive-height band&Cells&Precision&Proved lower bound for $|F|$\\ \hline
$[0,8]$&---&exact estimates&$1/120000$\\
$[7,41]$&272&256 bits&$9/5000000$\\
$[40,100]$&510&256 bits&$10^{-16}$\\
$[100,104]$&32&192 bits&$4\cdot10^{-16}$
\end{tabular}
\end{center}
The 510 cells comprise 160, 170, and 180 cells in $[40,60]$,
$[60,80]$, and $[80,100]$.  Their exact norm floor is
\[
 \frac{22271745300542300149807}
 {85070591730234615865843651857942052864}>10^{-16}.
\]
Exact rational cover checks include the spatial subdivisions,
all shared boundaries, and both time endpoints.  The spatial overlaps
$[7,8]$, $[40,41]$, and the shared boundary $x=100$, together with
evenness, prove \eqref{fc:height104-gap}.

The last band also certifies the centered-residual mechanism.  On each
of its 32 cells, freeze
\[
 a=\frac\pi{16},\qquad
 b=\frac14\log\frac{x_c+2t_ca}{4\pi\sin(4a)},\qquad
 \mu=a^2+b^2,
 \qquad C_j=-h_{j+2}-2ah_{j+1}-\mu h_j.
\]
The complete even source is shifted before theta indices are separated.
With eight retained indices, integration cutoff $3/2$, spatial order
12 and time order 8, the computation encloses all center jets through
order 28 and complete derivative majorants through order 29.
Centered coefficients are combined before Taylor transport, with
the parameters frozen on the cell.  The cell-constant normalization
$e^{\pi x_c/8}$ rescales each physical jet equally.
For normalized enclosures $P_0,P_1,Q_0,Q_1$ of $h_0,h_1,C_0,C_1$,
the actual decisions are
\[
 E_{\rm low}=\operatorname{dist}(0,P_0)^2+
 \operatorname{dist}(0,P_1)^2>0,\qquad
 E_{\rm low}/9-\sup|Q_0|^2-\sup|Q_1|^2>0.
\]
Thus $C_0^2+C_1^2\le E/9$ and $\mu\ge2/5$ throughout every cell.
The exact stationary identity
$\mathcal L_j=\mu h_j^2+h_jC_j$ at $h_{j+1}=0$ yields
\[
 \mathcal L_0\ge E/15\quad(h_1=0),\qquad
 \mathcal L_1\ge-5E/72\quad(h_2=0).
\]
The second strict stationary sign in
Theorem~\ref{fc:height104} is a separate validated signed decision;
it is stronger than the bounded-loss conclusion supplied by this cone.
At each level the band has 24 derivative-exclusion cells and eight
strict signed cells, with no claim that a signed cell alone proves
existence or uniqueness of a stationary branch.

Appendix~\ref{app:replay} identifies the retained drivers, records,
exact auditors, and earlier independent internal replays. In particular,
the higher-height auditor checks 6682 assertions and the centered-band
auditor checks 547. Those checks concern arithmetic, decisions, complete
coverage, and source bindings. The complete-source contour bounds used
for the latest band are proved in Appendix~\ref{app:payments}.
The computations establish the stated closed bounded region; they do
not extend it to arbitrary heights.


% Prepared for Edward Baker, 14 October 2026.
% LLM assistance: gpt-6-astra (Codex), reasoning effort ultra.
% Replacement for the former Section 7, integrating Notes 21--22.

\subsection{Extension through height 112 and all earlier-time comparisons}
\label{subsec:height112}

\begin{theorem}[Extended closed bounded certificate]
\label{fc:height112}
For the complete theta source,
\[
 \sqrt{h_0(t,x)^2+h_1(t,x)^2}>10^{-17}
       \qquad(0\le t\le1/2,\ |x|\le112).
\]
Both stationary implications $h_{j+1}=0\Rightarrow\Lag_j>0$,
$j=0,1$, hold on this rectangle.  With the fixed centering
$a=\pi/16$, $b=1/4$, $\mu=a^2+b^2$, the following stronger delayed
statements hold on the positive-height extension $104\le x\le112$:
at every current stationary point and every $0\le s\le t$,
\begin{equation}
 h_j(t,x)C_j(s,x)>0,\qquad
 h_j(t,x)\{C_j(t,x)-C_j(s,x)/2\}>0.
 \label{eq:height112-causal}
\end{equation}
There is exactly one $h_1=0$ point in $(107,108)$ and exactly one
$h_2=0$ point in $(109,110)$ at each time in $[0,1/2]$; there are
no other stationary points of these respective observations in
$[104,112]$.
\end{theorem}
\begin{proof}[Proof from the complete-source interval certificates]
The two band records in the numerical archive use the complete source
and error bounds already specified above:
\begin{quote}\small\raggedright
\path{numerics/FIXED_CUSP_BAND_104_108_20261014.json}\\
\path{numerics/FIXED_CUSP_BAND_108_112_20261014.json}.
\end{quote}
Each
has four unit spatial cells and eight closed time slabs
$[k/16,(k+1)/16]$, $0\le k\le7$.  Their fixed numerical parameters
are 192-bit arithmetic in python-flint 0.9.0, $a=\pi/16$, retained
index $N=8$, integration endpoint $U=3/2$, spatial order 12, and time
order 8.  Center physical jets through order 28 and complete moment
payments through order 29 include quadrature, full-index, integration,
rounding, spatial, and temporal errors.

The exact audit proves the whole-cell physical norm bound and the
normalized lower bound $e^{\pi x/4}E>3000$ on both new bands.  The stated
physical floor follows from the stronger cell-specific norm enclosures
and a rational exponential estimate using $\pi<22/7$; the displayed
coarse normalized bound alone is not used to obtain that floor.  The 128 stationary decisions are
\begin{center}
\begin{tabular}{ccl}
Band&Level&Complete closed-cell decisions\\\hline
$[104,108]$&$j=0$&24 exclusions; 8 strict signs in $(107,108)$\\
$[104,108]$&$j=1$&32 exclusions\\
$[108,112]$&$j=0$&32 exclusions\\
$[108,112]$&$j=1$&24 exclusions; 8 strict signs in $(109,110)$.
\end{tabular}
\end{center}
In each candidate spatial cell, paid endpoint signs of $h_{j+1}$ are
opposite and the full-cell $h_{j+2}$ enclosure has one strict sign.
This proves existence and uniqueness throughout the closed time slab.
The other cells exclude stationarity, so the implications cover branch
creation as well as regular branches.  The shared boundaries at 104 and
108 and the preceding certificate leave no spatial gap.  Physical parity
reflects the norm and stationary signs to negative heights.

For the delayed products the auditor reconstructs the residual from the
physical jets with $b=1/4$, rather than the saddle centering used for
the initial stationary test.  For current slab $k$, the earlier slabs
$p=0,\ldots,k$ cover all $0\le s\le t$, conservatively including
some $s>t$ in slab $k$.  The current and earlier time displacements
$d,e\in[-1/32,1/32]$ are independent; the spatial displacement
$\eta\in[-1/2,1/2]$ is the same in both observations.  If their paid
Taylor polynomials are $T_k(\eta,d)$ and $T_p(\eta,e)$, form the
coefficients of $T_k(\eta,d)-T_p(\eta,e)/2$ before ranging the common
$\eta$.  Explicitly, for coefficients $c_{k,r,m}$, the polynomial is
\[
 \sum_{m=0}^{11}\left[c_{k,0,m}-\frac12c_{p,0,m}
 +\sum_{r=1}^{7}\left(c_{k,r,m}d^r-\frac12c_{p,r,m}e^r\right)
 \right]\eta^m.
\]
Add the full current residual payment plus half the earlier payment;
only at the current time intersect with
$C_j=-h_{j+2}-\mu h_j$.  No earlier stationarity is imposed.
Taking the union over all earlier slabs and multiplying the complete
integrated observations preserves all source cross terms.

In the common constant normalization $Z_j=e^{\pi x_c/8}h_j$, with
the same factor for each residual, the final lower bounds for all
earlier times are
\begin{center}
\begin{tabular}{ccc}
Current branch&$Z_j\widehat C_j(s,x)$&
$Z_j(\widehat C_j(t,x)-\widehat C_j(s,x)/2)$\\\hline
$h_1=0$, $x_c=107.5$&$>15511$&$>7559$\\
$h_2=0$, $x_c=109.5$&$>6910$&$>3371$.
\end{tabular}
\end{center}
These enclosures prove \eqref{eq:height112-causal}.  They use complete
earlier fields, so they do not apply a multiplier bound restricted to
$\delta\le1/4$ when $t-s$ is larger.
\end{proof}

The earlier uniform residual-cone enclosure remains undecided on 24
cells in the first band and 16 in the second, all outside the candidate
stationary cells.  Those original verdicts are retained.  The successful
norm and stationary decisions do not turn the failed cone enclosures
into cone proofs or negative physical signs.  Likewise the delayed
products in \eqref{eq:height112-causal} are asserted for the positive-height
protocol actually computed; fixed drift in $C_j$ prevents an automatic
unchanged-drift reflection of that numerical claim.

The generating driver is
\path{numerics/certify_fixed_cusp_band_extension.py}, run with
\texttt{--x-left 104} and \texttt{--x-left 108}, each with an
\texttt{--output} path.  The driver
\path{numerics/audit_fixed_cusp_band_extension.py} takes both records
through repeated \texttt{--record} arguments and an \texttt{--output}
path.  Its retained
\path{numerics/FIXED_CUSP_BAND_EXTENSION_AUDIT_20261014.json}
contains 1321 passing checks, all 128 stationary observations, and all
16 genuine-branch comparisons.  Fresh 192-bit replays reproduced both
generator records after omitting only \texttt{elapsed\_seconds}; the
exact audit was reproduced in full.  Source bindings and these historical
replay claims are documented in
\path{numerics/PHASE_CONTROLLED_CUSP_CONTINUATION_CHECKS_20261014.json}
and Review 19.  They are internal validations, not formal verification
of the analytic source lemmas or the interval library.

\section{One-sided stationary estimates and their obstructions}
\label{sec:remaining}

The outer theorem and the finite cover give complementary information.
The unresolved positive-time domain is contained in the transition strip
$112<|x|<X_*$ and the high-height inner region
\begin{equation}
 0<t\le\tfrac12,\qquad |x|\ge X_*,\qquad
 t\log(|x|/(4\pi))<\frac{61}{24}.
 \label{eq:remaining-domain}
\end{equation}
The closed-boundary assertion is asymptotic: it does not supply a usable
numerical value of $X_*$. The separate explicit threshold-five bound is
not a numerical implementation of this sharper theorem. Even the
conditional completion down to $5/2$ leaves a smaller unbounded cusp.
The remaining analytic problem is a signed estimate in this domain,
with coefficients locally bounded through a candidate positive-time
collision. The following obstructions delimit that problem.

\subsection{Absolute dominance and pointwise positivity are too strong}

At the exact common-carrier coordinate $v$, put
$M=\lfloor e^{2v}\rfloor$. For $\lceil M/2\rceil\le n\le M$,
\[
 td_n^2-2td_nv=t[(d_n-v)^2-v^2]\ge-tv^2.
\]
Consequently
\begin{equation}
 \sum_{n\le e^{2v}}|A_n|
 \ge(M/2-1)M^{-1/2}e^{-tv^2}.
 \label{eq:absolute-channel-obstruction}
\end{equation}
With fixed $a$, $v=b_1$, and $t=\vartheta/\log x$ for fixed
$0\le\vartheta<4$, this grows as a positive constant times
$x^{1/4-\vartheta/16}$, whereas $|A_1|=1$. Absolute first-channel
dominance therefore fails in this scaling. Cancellation within the sum
and after integration remains available; indeed it is essential to the
improved outer theorem.

Pointwise positivity is also unavailable. For fixed $v=b$, write
\[
 D_j=\sum_{n\le e^{2b}}A_n(t,x;b)[-a+i(b-d_n)]^j,\qquad p_j=\Re(UD_j),
\]
where $U$ is the unit phase of
$e^{t(b+ia)^2}\phi_1(b+ia)e^{ixb}$. Put
\[
 \mathfrak H_j=|D_{j+1}|^2-\Re(D_j\overline{D_{j+2}}),\qquad
 \mathfrak W_j=D_{j+1}^2-D_jD_{j+2}.
\]
Then
\begin{equation}
 2(p_{j+1}^2-p_jp_{j+2})
 =\mathfrak H_j+\Re(U^2\mathfrak W_j).
 \label{eq:phase-projection}
\end{equation}
The stronger phase-independent condition
$\mathfrak H_j\ge|\mathfrak W_j|$ fails for the actual common-carrier
sums. With $a=\pi/16$ and $b=b_1$ from
\eqref{eq:channel-centers}, retained Arb/Acb enclosures give
\begin{center}
\begin{tabular}{@{}lll@{}}
\toprule
$x$ & $t$ & $\mathfrak H_0-|\mathfrak W_0|$ (approximate value)\\
\midrule
$211$ & $1/\log211$ & $-0.0277198$\\
$4898.7$ & $2/\log4898.7$ & $-1.26105$\\
\bottomrule
\end{tabular}
\end{center}
At the second point the actual phase also gives
$p_1^2-p_0p_2\approx-0.01099197<0$. The stored enclosures are strictly
negative; Appendix~\ref{app:replay} identifies their replay. These are
pointwise integrand forms, not the integrated Laguerre forms, and no
physical stationarity is asserted at these points. The paired identities
in Section~\ref{sec:residuals} retain correlations between distinct
coordinates as well as distinct indices. A proof requiring every
integrand to be positive would discard that possibility.
The staircase boundary identity \eqref{eq:staircase-cancellation} must
likewise be retained: entering terms have an exact cancellation of their
arithmetic $x$-phase, not a freely discardable endpoint error.

\subsection{A sufficient split that preserves favorable curvature}

Throughout this section let $0<a<\pi/8$, prescribe a real $b$, and put
$\mu=a^2+b^2$. At each point the centered observations are
\begin{equation}
 C_k=-h_{k+2}-2ah_{k+1}-\mu h_k,\qquad k\ge0.
 \label{eq:all-centered-observations}
\end{equation}
Whenever these observations are differentiated or compared at different
times below, $a,b$ are held fixed. The exact identity is
\[
 \Lag_j=\mu h_j^2+h_jC_j+h_{j+1}(h_{j+1}+2ah_j).
\]
Thus the remaining target can be stated as
\begin{equation}
 \boxed{h_jC_j\ge-\mu h_j^2-K_jE
 \quad\hbox{on }h_{j+1}=0,\qquad j=0,1,}
 \label{eq:final-target}
\end{equation}
where $K_j\ge0$ is prescribed without dividing by an unknown energy and
is locally bounded through every candidate positive-time collision.
This is precisely the stationary loss condition in
Section~\ref{sec:fc-collision}.

\begin{lemma}[Favorable curvature and a residual square]
\label{lem:favorable-split}
Prescribe an exact split $C_j=G_j+R_j$ and parameters
$0<\kappa\le\mu$, $\rho\ge0$. Define
\begin{equation}
 \Delta_j=4\kappa h_jG_j+\rho^2E-R_j^2.
 \label{eq:split-slack}
\end{equation}
Without imposing stationarity,
\begin{align}
 4\kappa\Lag_j+\rho^2E
 &=(R_j+2\kappa h_j)^2+\Delta_j
       +4\kappa(\mu-\kappa)h_j^2 \notag\\
 &\quad+4\kappa h_{j+1}(h_{j+1}+2ah_j).
 \label{eq:split-stationary-ideal}
\end{align}
In particular, $\Delta_j\ge0$ on $h_{j+1}=0$ implies
$\Lag_j\ge-\rho^2E/(4\kappa)$. This loss is locally bounded if
$\rho$ is locally bounded and $\kappa$ is locally bounded below by a
positive constant.
\end{lemma}
\begin{proof}
Expand the square and use $G_j+R_j=C_j$. All terms are polynomial in
the physical jets, so neither $E>0$ nor division by a jet is required.
On stationarity the last term vanishes and the other displayed squares
and $4\kappa(\mu-\kappa)h_j^2$ are nonnegative.
\end{proof}

The split must expose an independently estimable source contribution.
For the exact domain and fields of \eqref{eq:pair-domain}, choose a
prescribed real multiplier $m_j(\alpha)$ with at most polynomial growth
in $y_\alpha$, and set
\begin{align}
 G_j&=e^{-ax}\Re\int_\Omega m_j(\alpha)W_\alpha q_\alpha^j\dd\nu(\alpha),
       \notag\\
 R_j&=e^{-ax}\Re\int_\Omega r_j(\alpha)W_\alpha q_\alpha^j\dd\nu(\alpha),
 \qquad r_j(\alpha)=y_\alpha^2-b^2-m_j(\alpha).
 \label{eq:prescribed-source-split}
\end{align}
A proper source filter takes
$m_j=(y_\alpha^2-b^2)\mathbf1_A$, with the retained set $A$ intersected
with the original staircase. The filter and its parameters are fixed
before evaluating the physical observations. For real $m_j,r_j$, the
polarization \eqref{eq:paired-J} gives the exact scalar obligation
\begin{align}
 \Delta_j=\frac{e^{-2ax}}2\Re\iint_{\Omega^2}
 \big[&\{2\kappa(m_j(\alpha)+m_j(\beta))
                  -r_j(\alpha)r_j(\beta)\}J_j(\alpha,\beta)\notag\\
       &+\rho^2(J_0(\alpha,\beta)+J_1(\alpha,\beta))\big]
       \dd\nu(\alpha)\dd\nu(\beta).
 \label{eq:paired-split-slack}
\end{align}
Indeed $h_jG_j$ is $e^{-2ax}/4$ times the real part of the integral
of $(m_j(\alpha)+m_j(\beta))J_j$, whereas $R_j^2$ has coefficient
$e^{-2ax}/2$ and integrand $r_j(\alpha)r_j(\beta)J_j$.
The complete source majorants justify these identities and their tail
payments; they do not prove the sign in \eqref{eq:paired-split-slack}.
The formula requires real multipliers. A complex filter requires its
own conjugated polarization or an independent physical-field identity.

Even a quotient-free split can be vacuous: choosing
$m_j=y_\alpha^2-b^2+2\kappa$ makes $R_j=-2\kappa h_j$ and cancels
the first square in \eqref{eq:split-stationary-ideal}. With $\kappa=\mu$,
the resulting slack condition only restates the target and supplies no
theta estimate.

\subsection{Collision controls and the limitation of the square criterion}

For the backward-heat polynomial $P(t,y)=y^2-d^2-2t$, put
$D=d^2+2t>0$. Its Laguerre forms are $2y^2+2D>0$ and $4$.
At $y=0$, $h_0=-D$, $h_1=0$, $E=D^2$, $C_0=-2+\mu D$, and
$C_1=-4a$. The whole-residual cone would require
\begin{equation}
 \eta\ge\frac{\sqrt{(2-\mu D)^2+16a^2}}{D}.
 \label{eq:close-pair-cone}
\end{equation}
For bounded $a,\mu$, this grows without bound as the pair closes,
although the required negative loss is zero. The prescribed split
$G_0=-2$, $R_0=\mu D$ instead gives
\[
 \Delta_0=8\kappa D+(\rho^2-\mu^2)D^2.
\]
It accepts $\rho=0$ whenever $0<D\le8\kappa/\mu^2$, preserving the
large favorable curvature. This control is not a theta counterexample;
growth as $t\downarrow0$ is also compatible with local boundedness at
each positive time.

The same algebra rejects adverse collision predecessors. For
$P=y^2-2(t-T)$ at $y=0$, $t=T-s<T$, the split
$G_0=-2$, $R_0=-2\mu s$ yields
\[
 \Delta_0=-16\kappa s+4(\rho^2-\mu^2)s^2,\qquad
 \Lag_0=-4s,\qquad E=4s^2.
\]
For $P=y^3-6(t-T)y$ at $y=0$, $t=T-s<T$, the second stationary
set $h_2=0$ has $h_1=6s$, $h_3=6$. Taking $G_1=-6$,
$R_1=-6\mu s$ gives
\[
 \Delta_1=-144\kappa s+36(\rho^2-\mu^2)s^2,\qquad
 \Lag_1=-36s,\qquad E=36s^2.
\]
Both slacks are negative for sufficiently small $s>0$ when $\rho$
is bounded and $\kappa$ stays positively bounded below. The second
stationary set cannot be omitted.

There is also a precise restriction at the candidate collision itself:
if $h_0=h_1=0$, then $\Delta_0=-R_0^2$. The square criterion therefore
forces $R_0=0$ exactly; if also $h_2=0$, it forces $R_1=0$.
A small absolute bound on an omitted source does not imply this
vanishing. Thus Lemma~\ref{lem:favorable-split} is a sufficient
criterion with an extra cancellation obligation, not a solution of the
remaining theta problem.

Nor can its sign be justified by arbitrary-source positivity. On
$h_1=0$ the arbitrary-jet form $\Lag_0+KE=-h_0h_2+Kh_0^2$ is
indefinite for every finite $K$; on $h_2=0$ the form is
$-h_1h_3+K(h_0^2+h_1^2)$ and is likewise indefinite. Terms in the
stationary ideal vanish there. Any useful positivity argument must
therefore use relations of the actual theta source.

\subsection{Direct product absorption}

The residual-square condition can be replaced by a one-sided product
condition which imposes no residual vanishing at $E=0$.
\begin{lemma}[Fixed-fraction absorption]\label{lem:direct-absorption}
For a prescribed split $C_j=G_j+R_j$, let $0\le\lambda\le1$ and
$K_j\ge0$. On the genuine stationary set $h_{j+1}=0$, suppose
\begin{equation}
 h_jG_j\ge0,\qquad h_j(R_j+\lambda G_j)\ge-K_jE.
 \label{eq:direct-absorption}
\end{equation}
Then $\Lag_j\ge-K_jE$.
\end{lemma}
\begin{proof}
On this set,
\[
 \Lag_j=\mu h_j^2+(1-\lambda)h_jG_j
                      +h_j(R_j+\lambda G_j)\ge-K_jE.
\]
\end{proof}
At a first-pair zero the displayed products vanish automatically, but
the estimate must still hold with a locally bounded prescribed loss in
a neighborhood of that point. A favorable retained product alone does
not suffice. The finite results below use $\lambda=1/2$ and
$K_0=K_1=0$, without fitting a loss to an observed energy floor.


\section{Proper source filters and complete theta correlations}
\label{sec:source-filters}

We first test the preceding mechanisms on populated stationary branches
inside the existing height-$104$ cover. We then prove why fixed sharp
filters cannot support uniformly small separate payments at arbitrary
heights, and derive a complete-source alternative. The finite tests do
not extend the bounded cover or settle \eqref{eq:remaining-domain}.

\subsection{Two certified stationary branches and a signed comparison}

Fix
\begin{equation}
 a=\pi/16,\qquad b=1/4,\qquad \mu=a^2+b^2,\qquad
 \lambda=1/2,\qquad K_0=K_1=0.
 \label{eq:filter-pilot-parameters}
\end{equation}
Use the constant normalization $Z_k=e^{\pi x_c/8}h_k$, with
$x_c=101.5$ on $[101,102]$ and $x_c=103.5$ on $[103,104]$.
Hats denote the same normalization of centered observations.
The complete-source certificate and paid endpoint Taylor transport give,
for every $t\in[0,1/2]$,
\begin{align}
 Z_1(t,101)>37,\qquad& Z_1(t,102)<-20,\qquad Z_2<0
                       \quad(101\le x\le102),\notag\\
 Z_2(t,103)<-26,\qquad& Z_2(t,104)>14,\qquad Z_3>0
                       \quad(103\le x\le104).
 \label{eq:pilot-stationary-endpoints}
\end{align}
The intermediate value theorem and strict monotonicity therefore prove
exactly one $h_1=0$ point in $(101,102)$ and exactly one $h_2=0$ point
in $(103,104)$ at each time. The implicit function theorem gives smooth
branches. These are stationary points of the physical derivatives,
not derivatives of a variable normalization.

A first signed comparison uses the full residual on the first branch.
Exact rational extraction from the eight stored closed time slabs gives
\begin{equation}
 \begin{gathered}
 Z_0>0,\qquad\widehat C_0>0,\qquad
 Z_0\widehat C_0>6000,\qquad\widehat C_0-\mu Z_0>7/5,\\
 \mu Z_0^2+Z_0\widehat C_0>9400,\qquad
 \mu^2Z_0^2-\widehat C_0^2-\widehat C_1^2<-78.
 \end{gathered}
 \label{eq:full-residual-signed-pilot}
\end{equation}
Thus the signed condition holds with $K_0=0$, whereas every prescribed
whole-residual cone budget $\eta\le\mu$ fails on this branch.
The previously certified saddle-centered $\eta=1/3$ cone still passes;
the comparison concerns the different fixed centering
\eqref{eq:filter-pilot-parameters}, not a failure of every norm cone.
Taking $G=C$, $R=0$ here is a full-residual comparison, not a proper
source split. The next experiment supplies such a split.

\subsection{A proper filter absorbs an adverse complement on 64 cells}

Before evaluating partial integrals, prescribe the four filters
\[
 (N,U)\in\{(1,1),(1,3/2),(2,1),(2,3/2)\}.
\]
For $z=u+ia$, define
\begin{align}
 g_k(t,x)&=e^{-ax}\Re\sum_{n=1}^N\int_0^U
             (iz)^k e^{tz^2+ixu}\phi_n(z)\dd u,\notag\\
 G_j&=-g_{j+2}-2ag_{j+1}-\mu g_j,\qquad R_j=C_j-G_j.
 \label{eq:rectangular-filter}
\end{align}
This is precisely the real source multiplier
$(u^2-b^2)\mathbf1_{n\le N,\,0\le u\le U}$. Its shared-carrier
domain is $n\le N$, $d_n\le v\le d_n+U$, so all entry boundaries
are retained. Each of the eight time slabs
$[r/16,(r+1)/16]$, $0\le r\le7$, is tested on both branches:
four filters, two branches, and eight slabs give $64$ closed-cell tests.
Even-index slabs were designated development and odd-index slabs holdout
before evaluation; no filter, parameter, or loss was changed between them.
Their interiors are disjoint and their common boundary times are covered.

Both values of $U$ give the following conservative bounds on every
corresponding slab. Products are normalized by $e^{\pi x_c/4}$.
\begin{center}
\begin{tabular}{@{}ccrrr@{}}
\toprule
$N$ & $j$ & $Z_j\widehat G_j$ & $Z_j\widehat R_j$
    & $Z_j(\widehat R_j+\widehat G_j/2)$\\
\midrule
$1$ & $0$ & $>0$ & $<0$ & $<-1{,}200{,}000$\\
$1$ & $1$ & $>0$ & $<0$ & $<-250{,}000$\\
$2$ & $0$ & $>7800$ & $(-212,-120)$ & $>3800$\\
$2$ & $1$ & $>3800$ & $(-46,-25)$ & $>1900$\\
\bottomrule
\end{tabular}
\end{center}
Here $j=0$ means the $h_1=0$ branch and $j=1$ the $h_2=0$ branch.
The two-index filters therefore pay a genuinely adverse complementary
product with the prescribed half of a favorable retained product.
Lemma~\ref{lem:direct-absorption} applies with the fixed zero loss.
For one index the same prescribed absorption test actually fails on
nonempty stationary sets: the negative upper bounds are not an interval
precision failure. The large one-index terms nearly cancel in the
complete residual.

For completeness, the certification encloses retained center derivatives
through order $28$ with $192$-bit Arb quadrature, then uses space order
$12$ and time order $8$ for transport. Fixed endpoints give
$\partial_xg_k=g_{k+1}$ and $\partial_tg_k=-g_{k+2}$ exactly.
The complete-source moment payments remain valid for the three masks
$\mathbf1_A$, $1-\mathbf1_A$, and $1-\mathbf1_A/2$, all of which
lie in $[0,1]$. Centered payments are bounded by
$P_{j+2}+2aP_{j+1}+\mu P_j$ for each spatial or temporal remainder.
The full $C$ arrays already include quadrature error, all omitted theta
indices, and the infinite integration tail. The complement is the exact
difference $C-G$, and $C-G/2$ is formed before transport to preserve
cancellation. On current stationarity its enclosure is intersected with
$-h_{j+2}-\mu h_j-G_j/2$; the complement is treated similarly.
The product always uses the full physical jet and an integrated
observation, retaining all cross terms. No individual theta index is
assigned the vanishing vertical contour integral of the complete source.
The complete $64$-case receipt, including failures, is identified in
Appendix~\ref{app:replay}.

\subsection{The endpoint terms of a fixed sharp filter}

Now fix $t\in(0,1/2]$, $0<a<\pi/8$, $b\ge0$, $N\ge1$, and
$U>0$, independently of $x$. Write
\[
 \Phi_N(z)=\sum_{n=1}^N\phi_n(z),\qquad q=-a+iu=iz,\qquad
 f_{N,j}(u)=(u^2-b^2)q^j e^{tz^2}\Phi_N(z).
\]
Then the retained residual in \eqref{eq:rectangular-filter} is
\begin{equation}
 G_j(x)=e^{-ax}\Re\int_0^U f_{N,j}(u)e^{ixu}\dd u,\qquad R_j=C_j-G_j.
 \label{eq:sharp-filter-integral}
\end{equation}
These fixed filtered fields still solve backward heat. The following
effect comes from source integration by parts, not a heat-equation defect.

\begin{lemma}[Complete endpoint expansion]\label{lem:sharp-endpoints}
For $x>0$,
\begin{align}
 e^{ax}G_j(x)&=\frac{B_j(x)}x+\mathcal E_j(x),\qquad
 B_j(x)=\Im\{f_{N,j}(U)e^{ixU}-f_{N,j}(0)\},\notag\\
 |\mathcal E_j(x)|&\le\frac{M_{N,j}}{x^2},\qquad
 M_{N,j}=|f'_{N,j}(U)|+|f'_{N,j}(0)|
                      +\int_0^U|f''_{N,j}(u)|\dd u.
 \label{eq:sharp-endpoint-expansion}
\end{align}
In particular the lower coefficient is
\begin{equation}
 -\Im f_{N,j}(0)=b^2(-a)^j e^{-ta^2}\Im\Phi_N(ia).
 \label{eq:raw-index-endpoint}
\end{equation}
\end{lemma}
\begin{proof}
Twice integrating $\int_0^Ufe^{ixu}\dd u$ by parts gives
\[
 \frac{[fe^{ixu}]_0^U}{ix}
 -\frac{[f'e^{ixu}]_0^U}{(ix)^2}
 +\frac{1}{(ix)^2}\int_0^Uf''e^{ixu}\dd u.
\]
Taking real parts uses $\Re(-iZ)=\Im Z$ and gives the stated bound.
Substitution at $u=0$ proves \eqref{eq:raw-index-endpoint}.
\end{proof}

For the complete theta source, evenness and Schwarz symmetry make
$\Phi(ia)$ real. More generally its centered integrand obeys
$f_j(-u)=\overline{f_j(u)}$. Its even derivatives at zero are real
and its odd derivatives purely imaginary, so every lower-endpoint
coefficient in the real integration-by-parts expansion vanishes.
A finite raw-index sum need not have this symmetry.
If instead a common cutoff $v\le V$ is used, retain
$U_n=V-d_n>0$ for every included index; the leading term becomes
\[
 \Im\left\{\sum_{\substack{n\le N\\d_n<V}}
 f_{n,j}(U_n)e^{ixU_n}
 -\sum_{\substack{n\le N\\d_n<V}}f_{n,j}(0)\right\}.
\]
Here $f_{n,j}$ is defined with the individual term $\phi_n$ in place
of $\Phi_N$. The remainder is the sum of the bounds in
\eqref{eq:sharp-endpoint-expansion}. An index with $d_n=V$ has an
empty interval and contributes zero, not a single unmatched endpoint.
This is the same entry cancellation as
\eqref{eq:staircase-cancellation}.

The upper endpoint is provably nonzero for the pilot family, without a
genericity assumption. More generally suppose
$a=\pi/16$, $U\ge1$, and $U\ne b$. Put $D_4=e^{4U}>50$ and
$\beta=\pi\cos(4a)D_4>100$. For $n\ge2$,
\[
 \left|\frac{\phi_n(U+ia)}{\phi_1(U+ia)}\right|
 \le n^2\frac{2\pi n^2D_4+3}{2\pi D_4-3}
             e^{-\beta(n^2-1)}
 \le2n^4e^{-100(n^2-1)}.
\]
Here $\phi_1(U+ia)\ne0$ since
$|2\pi e^{4(U+ia)}-3|\ge2\pi D_4-3>0$.
Using $n^4\le e^{n^2-1}$ and
$n^2-1\ge3+5(n-2)$ gives
\begin{equation}
 \sum_{n=2}^\infty2n^4e^{-100(n^2-1)}
 \le\frac{2e^{-297}}{1-e^{-495}}<\tfrac12.
 \label{eq:nonzero-sharp-endpoint}
\end{equation}
Thus $\Phi_N(U+ia)\ne0$ for every finite $N$, and
$f_{N,0}(U)\ne0$. This includes $U=1,3/2$, $b=1/4$.

\begin{theorem}[Adverse sharp-filter products on genuine stationary sets]
\label{thm:sharp-filter-obstruction}
Fix the parameters of \eqref{eq:sharp-filter-integral} and assume
$f_{N,0}(U)\ne0$. For each $j=0,1$ there is a sequence
$x_k\to\infty$ of genuine stationary points $h_{j+1}(t,x_k)=0$ such
that $h_jG_j>0$, $h_jR_j<0$, and
\begin{equation}
 \frac{-h_jR_j}{E}(t,x_k)
   =\exp\{(\pi/8-a)x_k+o(x_k)\}.
 \label{eq:sharp-filter-adverse-rate}
\end{equation}
For every fixed $0\le\lambda<1$, along the same sequence eventually
\begin{equation}
 h_j(R_j+\lambda G_j)<0,\qquad
 \frac{-h_j(R_j+\lambda G_j)}E
   =\exp\{(\pi/8-a)x_k+o(x_k)\}.
 \label{eq:sharp-filter-absorption-rate}
\end{equation}
There are also stationary sequences with $h_jR_j>0$ and the same
exponential rate for its absolute ratio. All these points eventually
belong to the proved outer region, where $E>0$ and both complete
Laguerre forms are positive.
\end{theorem}
\begin{proof}
At fixed $t>0$, the outer four-jet estimates apply for all sufficiently
large $x$. With $A(x)=2|Q_t(x)|$, $\omega=\log(x/(4\pi))/4$, and
a continuous phase $\varphi=\arg Q_t$, their fixed-time consequences are
\begin{equation}
 \begin{gathered}
 h_k=A\omega^k\{\cos(\varphi+k\pi/2)+O(1/\log x)\},
       \qquad 0\le k\le3,\\
 \varphi'(x)=\omega+O(1)>0,\qquad
 \log A(x)=-\pi x/8+O_t((\log x)^2).
 \end{gathered}
 \label{eq:fixed-time-phase-windows}
\end{equation}
For $j=0$, take phase windows centered at $\varphi=k\pi$ with
half-width $\pi/6$. For all sufficiently large windows, $h_1$ has
opposite endpoint signs and $h_2$ has a fixed nonzero sign throughout.
There is exactly one $h_1=0$ point in each window, and the sign of
$h_0$ alternates. There $|h_0|\asymp A$ and $E=h_0^2\asymp A^2$.

For $j=1$, center the windows at $\varphi=\pi/2+k\pi$. Now $h_2$
has opposite endpoint signs and $h_3$ has a fixed nonzero sign.
There is exactly one $h_2=0$ point in each window and $h_1$ alternates
sign. At these points $|h_1|\asymp A\omega$, while $h_2=0$ in
\eqref{eq:fixed-time-phase-windows} gives $h_0=O(A/\log x)$;
therefore $E\asymp A^2\omega^2$.

Since $q(U)\ne0$, the hypothesis makes each $B_j$ a nonconstant
sinusoid plus a constant. It has infinitely many fixed-width intervals
on which its sign is fixed and $|B_j|\ge c_j>0$. The spacing of the
stationary points just constructed is $O(1/\log x)$. Each such interval
at sufficiently large height therefore contains stationary points of
both signs of $h_j$. By \eqref{eq:sharp-endpoint-expansion}, $G_j$
has the sign of $B_j$ there and $|G_j|\asymp e^{-ax}/x$.
In contrast, $C_j$ is a fixed linear combination of complete physical
jets, so \eqref{eq:fixed-time-phase-windows} makes it exponentially
smaller than $G_j$. Hence $R_j=-G_j(1+o(1))$ on these intervals.

Choose the stationary points with $h_jG_j>0$ or with $h_jG_j<0$.
The energy estimates above give \eqref{eq:sharp-filter-adverse-rate}
and its opposite-sign version: all powers of $x,\log x$ and the
$O_t((\log x)^2)$ amplitude correction contribute $o(x)$ to the
logarithm. Finally
$h_j(R_j+\lambda G_j)=h_jC_j-(1-\lambda)h_jG_j$, and the same
comparison proves \eqref{eq:sharp-filter-absorption-rate}.
\end{proof}

Thus no constant, polynomial, or other budget $K_j(x)=\exp(o(x))$
can make the second inequality in \eqref{eq:direct-absorption} hold
globally for this fixed sharp filter and a fixed $\lambda<1$.
This conclusion does not exclude losses locally bounded at every finite
positive-time point: exponential growth in $x$ is still locally bounded.
It is a fixed-positive-time obstruction in an already favorable region,
not an obstruction to an adaptive filter in the unresolved moving cusp.
Neither $\lambda=1$ nor a fraction tending exponentially fast to one is
covered. In the complete product $h_jC_j$, the artificial endpoint
terms cancel exactly.

\subsection{A compact filter has no finite rational jet closure}

\begin{proposition}\label{prop:no-rational-filter-closure}
For any fixed $t,a,b,N,U$ as in
\eqref{eq:sharp-filter-integral}, no finite family of real rational
functions $r_0,\ldots,r_m$ can satisfy
\begin{equation}
 R_j(x)=\sum_{k=0}^m r_k(x)h_k(t,x)
 \label{eq:forbidden-rational-closure}
\end{equation}
on a nonempty open real interval free of their poles, for $j=0$ or $1$.
\end{proposition}
\begin{proof}
Since $C_j$ is a constant-coefficient combination of physical jets,
clearing denominators would give
\[
 p(x)G_j(x)=\sum_{k=0}^m p_k(x)h_k(t,x),\qquad p\ne0,
\]
with polynomials. Both sides extend to entire functions, so this
identity holds for all real $x$.
Complete evenness and Schwarz symmetry give
\begin{equation}
 e^{ax}h_k(t,x)=\frac12\int_{\R}
 e^{t(u+ia)^2}\Phi(u+ia)[i(u+ia)]^k e^{ixu}\dd u.
 \label{eq:whole-line-filter-closure}
\end{equation}
The Fourier density is real analytic on all of $\R$, and it and all its
derivatives decay rapidly: holomorphy follows from the open theta strip,
positive-end decay from the theta series, and negative-end decay from
the complete evenness. Polynomial multiplication in $x$ corresponds
to finite derivatives in $u$, preserving real analyticity.

By contrast, $e^{ax}G_j$ is the Fourier transform of the compactly
supported $L^1$ density equal to $f_{N,j}(u)/2$ on $(0,U)$,
$\overline{f_{N,j}(-u)}/2$ on $(-U,0)$, and zero elsewhere.
Multiplication by $p(x)$ gives a compactly supported distribution.
Fourier injectivity for tempered distributions equates this compact
distribution to a real analytic function. That function vanishes on
$(U,\infty)$, hence everywhere. Therefore $pG_j=0$ and $G_j=0$.

This is impossible. The finite sum $\Phi_N$ is a nonzero holomorphic
function, since every $\phi_n$ is positive on the positive real axis.
The other factors in $f_{N,j}$ have only isolated zeros, so its compact
Hermitian extension is not zero. Fourier injectivity implies $G_j$ is
not identically zero. The injectivity used here follows, for example,
by Gaussian regularization and the approximate-identity limit; it
requires no information about zeros of $H$.
\end{proof}

The proposition excludes rational-coefficient finite-jet identities for
these fixed filters, including constant and polynomial coefficients.
It does not exclude general transcendental coefficients, merely locally
bounded coefficients, or moving filters. Nor can complete symmetry be
restored by naively adding $\Phi_N(-z)$: as real $u\to\infty$,
\[
 \Phi_N(-u)=-3\pi\Big(\sum_{n=1}^Nn^2\Big)e^{-5u}+O_N(e^{-9u}),
\]
so multiplication by $e^{tu^2}$ destroys integrability at every
positive time. Complete theta evenness is an infinite-index identity.
A vertical contour correction alone does not restore this raw-index
symmetry. Smooth compact cutoffs still require matched endpoint
derivatives and a new relative signed estimate.

\subsection{Complete-source Gaussian transport and a 14-cell pilot}

There is an exact alternative without sharp source endpoints. For
fixed $\delta\ge0$, multiply the complete source on the shifted
contour by $e^{-\delta z^2}$, $z=u+ia$. With the same frozen $a,b$
at both times, its retained observation and complement are
\begin{equation}
 G_j(t,x)=C_j(t-\delta,x),\qquad
 R_j(t,x)=C_j(t,x)-C_j(t-\delta,x).
 \label{eq:complex-gaussian-transport}
\end{equation}
This follows directly from $e^{tz^2}e^{-\delta z^2}=e^{(t-\delta)z^2}$
and the complete contour identity. The multiplier preserves the
symmetry on the imaginary segment and introduces no raw-index or upper
cutoff boundary. It is generally complex on the horizontal contour:
\[
 e^{-\delta(u+ia)^2}
     =e^{\delta a^2-\delta u^2}e^{-2i\delta au}.
\]
Accordingly \eqref{eq:paired-split-slack}, which assumes real
multipliers, is not applied unchanged. Formula
\eqref{eq:complex-gaussian-transport} supplies the physical products
directly. Numerical truncation errors still require full payment.

For the parameters \eqref{eq:filter-pilot-parameters}, fix
$\delta=1/16$ before testing. The two current-time branches in
\eqref{eq:pilot-stationary-endpoints} and the seven closed time slabs
$[r/16,(r+1)/16]$, $1\le r\le7$, give $14$ predeclared cases.
All satisfy $h_jG_j>0$, $h_jR_j>0$, and
$h_j(R_j+G_j/2)>0$. Conservative normalized margins are
\begin{center}
\begin{tabular}{@{}crrr@{}}
\toprule
$j$ & $Z_j\widehat G_j$ & $Z_j\widehat R_j$
    & $Z_j(\widehat R_j+\widehat G_j/2)$\\
\midrule
$0$ & $>8300$ & $>54$ & $>4200$\\
$1$ & $>4200$ & $>7$ & $>2100$\\
\bottomrule
\end{tabular}
\end{center}
Thus this second finite pilot has favorable retained and complementary
products, with the same prescribed zero loss. Stationarity is imposed
only at $(t,x)$, never at $(t-\delta,x)$.

The proof is exact rational extraction from the stored full-source
certificate and requires no new quadrature. Previous and current slabs
have the same spatial center and time displacements. Center coefficient
differences are formed before transport. The complete spatial and
temporal payments are $P_{\rm prior}$ for $G$,
$P_{\rm current}+P_{\rm prior}$ for $R$, and
$P_{\rm current}+P_{\rm prior}/2$ for $R+G/2$.
All full center arrays already contain the complete index and integration
tails. Intersecting with the current stationary identity
$C_j=-h_{j+2}-\mu h_j$ preserves that constraint, and products of full
integrated observations retain every cross term. The receipt in
Appendix~\ref{app:replay} records all $14$ cases. It proves the stated
finite comparison on $1/16\le t\le1/2$, not persistence of the
correlation near an unknown collision or suitability of a fixed delay
as $t\downarrow0$.

The distinct real Gaussian multiplier
$(u^2-b^2)e^{-\delta u^2}$ also retains the complete source. Denote
its retained residual by $G^{\rm real}_{\delta,j}$. It satisfies
\begin{equation}
 G^{\rm real}_{\delta,j}(t,x)
     =e^{a^2\delta}C_j(t-\delta,x+2a\delta).
 \label{eq:real-gaussian-transport}
\end{equation}
Indeed
$tz^2-\delta u^2=(t-\delta)z^2+2i\delta au-\delta a^2$.
After including the physical prefactor $e^{-ax}$, the corresponding
filtered jet is $e^{a^2\delta}h_k(t-\delta,x+2a\delta)$.
Applying $-\partial_x^2-2a\partial_x-\mu$ proves
\eqref{eq:real-gaussian-transport}.
Both identities converge for all real times; $0\le\delta\le t$
is imposed when using a theorem restricted to nonnegative time.
The real mask obeys the hypotheses of the paired multiplier formula,
but no numerical certificate for this additional mask is asserted.
These are standard Fourier/heat identities specialized to the complete
theta source, not consequences of an unproved positivity property.

Their exact complementary remainders expose the next analytic obligation.
Since the frozen observations satisfy $\partial_tC_j=-C_{j+2}$,
\begin{align}
 R_j(t,x)&=-\int_0^\delta C_{j+2}(t-s,x)\dd s,
       \label{eq:complex-gaussian-remainder}\\
 C_j(t,x)-G^{\rm real}_{\delta,j}(t,x)
   &=-\int_0^\delta e^{a^2s}
       [C_{j+2}+2aC_{j+1}+a^2C_j](t-s,x+2as)\dd s.
       \label{eq:real-gaussian-remainder}
\end{align}
The earlier or displaced point need not be stationary.

\subsection{Complex delay polarization and moving-parameter identities}
\label{subsec:complex-delay-polarization}

The Gaussian comparison uses a complex multiplier on the shifted contour.
Its paired formula must therefore conjugate the multiplier in the second
Hermitian factor.  To state this without identifying source coordinates,
use $\Omega_0=[0,\infty)\times\mathbb N$.  For $\alpha=(u,n)$ set
\[
 z_\alpha=u+ia,\quad q_\alpha=iz_\alpha,\quad
 D_\alpha=u^2-b^2,\quad
 W_\alpha=e^{tz_\alpha^2+ixu}\phi_n(z_\alpha),
\]
and put
$\mathcal F_j(k)=e^{-ax}\Re\int_{\Omega_0}W_\alpha q_\alpha^j
k_\alpha\,d\nu(\alpha)$.  The complete-source contour identity gives
$h_j=\mathcal F_j(1)$, $C_j=\mathcal F_j(D)$.
For $G_j=C_j(t-\delta,x)$, the multipliers for $P_j=h_jG_j$ and
$Q_j=h_j(C_j-G_j/2)$ are respectively
\[
 k^P_\alpha=D_\alpha e^{-\delta z_\alpha^2},\qquad
 k^Q_\alpha=D_\alpha(1-e^{-\delta z_\alpha^2}/2).
\]
The complete source is shifted before its arithmetic terms are separated.
In particular no individual raw theta term is assumed even.

\begin{lemma}[Complex polarization and full pair-tail payment]
\label{lem:complex-delay-polarization}
For the above multipliers, with independent $\alpha,\beta$,
\begin{align}
 h_j\mathcal F_j(k)=\frac{e^{-2ax}}4\Re\iint_{\Omega_0^2}
 \big[&(k_\alpha+\overline{k_\beta})W_\alpha\overline{W_\beta}
                       (q_\alpha\overline{q_\beta})^j\notag\\
     &+(k_\alpha+k_\beta)W_\alpha W_\beta(q_\alpha q_\beta)^j
 \big]d\nu(\alpha)d\nu(\beta).
 \label{eq:complex-delay-polarization}
\end{align}
For $a=\pi/16$, $b=1/4$, $0\le t\le1/2$ and
$0\le\delta\le\min(t,1/4)$, put $M=e^{a^2/4}$.  If $A$ is
a retained set, let $M_{A,k}=\int_A|W||q|^k$ and
$T_k=\int_{A^c}|W||q|^k$.  The absolute multiplier masses satisfy
\[
 B_A^P\le M(M_{A,j+2}+b^2M_{A,j}),\qquad
 B_A^Q\le(1+M/2)(M_{A,j+2}+b^2M_{A,j}),
\]
and the same bounds hold on $A^c$ with $M_{A,k}$ replaced by $T_k$.
Truncating both factors of either product has error at most
\begin{equation}
 e^{-2ax}\{M_{A,j}B_{A^c}+T_jB_A+T_jB_{A^c}\}.
 \label{eq:complex-delay-tail}
\end{equation}
\end{lemma}
\begin{proof}
Use $\Re U\Re V=\frac12\Re(U\overline V+UV)$ and then average
under interchange of the two source coordinates.  This gives the two
sectors and the conjugated multiplier in
\eqref{eq:complex-delay-polarization}.  Absolute convergence follows,
with $c_a=\cos(4a)$ and $\rho_0=16e^{-3\pi c_a}<1$, from
\[
 \sum_n|W_{(u,n)}||q_{(u,n)}|^k
 \le\frac{2\pi^2+3\pi}{1-\rho_0}(u+a)^k
           e^{u^2/2+9u-\pi c_a e^{4u}},
\]
which has every polynomial moment.  The inequality
$|e^{-\delta(u+ia)^2}|\le M$ and
$|u^2-b^2|\le|q|^2+b^2$ prove the multiplier bounds.  Expanding a
product into its retained and omitted factors proves
\eqref{eq:complex-delay-tail}, including both mixed sectors and the
omitted--omitted sector.  The complete index and integration bounds of
Appendix~\ref{app:payments} supply the $T_k$ on their declared
rectangles.  Quadrature and parameter-transport errors remain additional.
\end{proof}

The real horizontal Gaussian $e^{-\delta u^2}$ is a different mask;
its time-and-position transport identity above must not replace the
complex mask in this lemma.  At a fixed observation point the latter
has the exact time identity
\[
 G_j=C_j(t-\delta,x),\qquad
 C_j(t,x)-G_j=-\int_0^\delta C_{j+2}(t-s,x)\,ds.
\]
Earlier stationarity is never imposed.  If $\delta$ depends on $t,x$,
its derivatives also enter transport.  With superscript $-$ denoting
evaluation at $(t-\delta(t,x),x)$ and fixed $a,b$,
\begin{align}
 G_{j,t}&=-(1-\delta_t)C_{j+2}^-,&
 G_{j,x}&=C_{j+1}^-+\delta_x C_{j+2}^-,\notag\\
 (\partial_t+\partial_x^2)G_j
 &=(\delta_t+\delta_{xx})C_{j+2}^-
       +2\delta_xC_{j+3}^-+\delta_x^2C_{j+4}^-.
 \label{eq:general-moving-delay}
\end{align}
These identities follow by the chain rule and $\partial_tC_k=-C_{k+2}$.
All needed derivatives are justified by the complete-source moment
majorant, not by differentiating a four-jet asymptotic.
On a regular current branch $x=\xi(t)$,
\[
 \xi'=h_{j+3}/h_{j+2},\qquad
 \dot h_j=-h_{j+2},\qquad
 \dot G_j=\xi'C_{j+1}^--(1-\delta_t-\xi'\delta_x)C_{j+2}^-.
\]
Thus $\dot P_j=-h_{j+2}G_j+h_j\dot G_j$ and
\[
 \dot Q_j=-h_{j+2}(C_j-G_j/2)
       +h_j[-C_{j+2}+\xi'C_{j+1}-\dot G_j/2].
\]
The denominator-free tangent operator
$\mathcal T_j=h_{j+2}\partial_t+h_{j+3}\partial_x$ satisfies
\[
 \mathcal T_jh_j=-h_{j+2}^2,\quad
 \mathcal T_jG_j=-[h_{j+2}(1-\delta_t)-h_{j+3}\delta_x]C_{j+2}^-
                              +h_{j+3}C_{j+1}^-.
\]
It may vanish at a singular point.  None of these transport identities
alone proves continuation or positivity through branch creation.

\subsection{Why every fixed proportional delay fails its exact retained sign}
\label{subsec:proportional-delay-obstruction}

\begin{theorem}[Proportional-delay phase obstruction]
\label{thm:proportional-delay-obstruction}
Fix $t_0\in(0,1/2]$, $0<\rho<1$, and finite real centering
parameters $a,b$.  Put $\delta=\rho t_0$, $\tau=(1-\rho)t_0$.
For each $j=0,1$, the genuine current set $h_{j+1}(t_0,x)=0$
contains sequences $x\to+\infty$ on which
$h_j(t_0,x)C_j(\tau,x)<0$, and other sequences with the opposite
strict sign.  Nevertheless, for every fixed $\lambda\in[0,1]$,
\[
 h_j(t_0,x)\{C_j(t_0,x)-(1-\lambda)C_j(\tau,x)\}>0
\]
at every sufficiently high current stationary point.  Both full
Laguerre forms are also positive there.
\end{theorem}
\begin{proof}
Use the analytic carrier $Q_v$ in \eqref{eq:analytic-model} at each
fixed positive time $v=t_0,\tau$, and put
$A_v=2|Q_v|$, $\phi_v=\arg Q_v$, $L=\log(x/(4\pi))$,
$\omega=L/4$.  The four-jet theorem applies at both times for all
sufficiently large $x$, giving
\[
 h_k(v,x)=A_v\omega^k\{\cos(\phi_v+k\pi/2)+O(L^{-1})\},
 \quad 0\le k\le3,\qquad \phi_{t_0}'=\omega+O(1)>0.
\]
The exact carrier quotient and
$\alpha(s)=L/2+i\pi/4+O(x^{-1})$, $s=(1+ix)/2$, give
\begin{equation}
 D:=\phi_{t_0}-\phi_\tau=\frac{\delta\pi L}{16}+o(1),\qquad
 r_\delta:=\frac{A_\tau}{A_{t_0}}
  =\exp\{-\delta L^2/16+\delta\pi^2/64+o(1)\}.
 \label{eq:fixed-delay-phase-amplitude}
\end{equation}
Indeed the two zeta arguments have real parts tending to infinity, so
their Dirichlet series and their analytic logarithms tend to 1 and 0,
respectively.  Continuous carrier phases may be chosen consistently.

At current stationarity let $\psi=\phi_{t_0}+j\pi/2$; then
$\sin\psi=O(L^{-1})$.  The lower-order drift and mass terms in the
residual contribute only $O(L^{-1})$ after normalization.  Hence
\[
 \frac{h_j(t_0,x)C_j(\tau,x)}
      {A_{t_0}A_\tau\omega^{2j+2}}=\cos D+O(L^{-1}),\qquad
 \frac{h_j(t_0,x)C_j(t_0,x)}
      {A_{t_0}^2\omega^{2j+2}}=1+O(L^{-1}).
\]
The normalized absorption expression tends to 1 because
$r_\delta\to0$.

To obtain actual stationary points, take each sufficiently high phase
window $|\psi-k\pi|\le\pi/6$.  The signs of $h_{j+1}$ at its
endpoints are opposite, and $h_{j+2}$ has one strict sign throughout,
so the window contains exactly one current stationary point.  On the
logarithmic bands
\[
 \left|L-\frac{16(2n+1)}\delta\right|\le\frac8{3\delta},
\]
the phase $D$ is within $\pi/6+o(1)$ of $(2n+1)\pi$ and
$\cos D\le-1/2$ eventually.  Their lengths in $x$ grow exponentially,
whereas the phase windows have width of order $1/L$, so each sufficiently
high band contains complete windows of both levels.  Bands centered at
$L=32n/\delta$ give the positive sign.  All these points lie in the
already proved outer-sign region.
\end{proof}

The obstruction also excludes a height-independent rule $\delta(t)$
whenever $0<\delta(t_0)<t_0$ at one positive time.  It concerns the
exact retained-sign requirement, not the bounded-loss program: on these
proved nonzero fixed-time branches,
\[
 P_j/E=r_\delta\omega^2\{\cos D+O(L^{-1})\}\longrightarrow0.
\]
The sign alternates while its energy-relative magnitude vanishes.
No first counterexample height is evaluated, and this fixed-time argument
does not settle a moving inner cusp.

For the particular half-time rule $\delta=t/2$, a complete-source pilot
in $(101,102)$ on $h_1=0$ and $(103,104)$ on $h_2=0$ certifies
$P_j,Q_j>0$ in all 16 branch/slab cases over $0\le t\le1/2$.
In the constant branch normalizations its lower bounds for $(P_j,Q_j)$
are $(8108,3999)$ and $(4100,2025)$.  The complement $h_j(C_j-G_j)$
is undecided on the first two time slabs of each branch, and equals zero
at $t=0$; these four undecided intervals are not negative signs.
The exact earlier-time map on slab $k$ is
$d^-=d/2+\epsilon_k$, with $\epsilon_k=-1/64$ for even $k$ and
$1/64$ for odd $k$.  Substitution in the common Taylor polynomial,
before interval ranging and with all original payments retained, proves
the pilot.  Its record is
\path{numerics/PROPORTIONAL_GAUSSIAN_SOURCE_FILTER_PILOT_20261014.json}.
This bounded success is consistent with the asymptotic obstruction.

\subsection{A height-dependent delay with a uniform outer comparison}
\label{subsec:phase-controlled-delay}

Keep $a=\pi/16$, $b=1/4$ fixed, and prescribe
\begin{equation}
 \ell(x)=\log(1+x^2),\quad B=1+t\ell(x),\quad
 \delta(t,x)=\frac{t}{2B},\quad\tau=t-\delta,
 \quad G_j=C_j(\tau,x).
 \label{eq:smooth-delay-rule}
\end{equation}
For $t\ge0$, $0\le\delta\le t/2$ and $t/2\le\tau\le t$.

\begin{theorem}[Phase-controlled outer comparison]
\label{thm:phase-controlled-delay}
There is a finite unevaluated $X$ such that, uniformly for
$|x|\ge X$, $0<t\le1/2$, and
$\theta=t\log(|x|/(4\pi))\ge13/4$, the rule
\eqref{eq:smooth-delay-rule} satisfies at current stationarity
\[
 P_j\ge\frac18A_tA_\tau\omega^{2j+2}>0,\qquad
 Q_j\ge\frac12A_t^2\omega^{2j+2}>0,\qquad j=0,1,
\]
where $A_v=2|Q_v|$ is the analytic carrier amplitude at time $v$ and
$\omega=\frac14\log(|x|/(4\pi))$.
\end{theorem}
\begin{proof}
Take $x>4\pi$ and write $L=\log(x/(4\pi))$, $d=\delta L$.
Since $\ell(x)>2L$,
\[
 0<d<\frac{\theta}{2+4\theta}<\frac14,\qquad
 \tau L>f(\theta):=\theta-\frac{\theta}{2+4\theta}.
\]
The increasing function $f$ satisfies $f(13/4)=91/30=3+1/30$.
Thus the uniform four-jet asymptotic is available at both times with
a fixed margin, even when those times decrease with height.  Write
$s=(1+ix)/2$, $w_v=s+v\alpha(s)/2$.  The full relative phase is
\begin{equation}
 D:=\phi_t-\phi_\tau
 =\Im\left\{\frac\delta4\alpha(s)^2+
       \int_\tau^t\frac{\alpha(s)}2
                \frac{\zeta'}{\zeta}(w_v)\,dv\right\}.
 \label{eq:smooth-delay-full-phase}
\end{equation}
The analytic prime-power logarithm is valid throughout this path:
$\Re w_v\ge5/4$ uniformly once $X$ is large.  In that half-plane,
\[
 \left|\frac{\zeta'}\zeta(w)\right|
 \le\sum_{n\ge2}(\log n)n^{-5/4}<18.
\]
For the last inequality bound the $n=2,3$ terms by 1 each and the
decreasing remainder by
$\int_3^\infty(\log u)u^{-5/4}du<16$.
Since $\alpha=L/2+i\pi/4+O(x^{-1})$ and $d<1/4$, the gamma and
arithmetic parts of \eqref{eq:smooth-delay-full-phase} are bounded by
$\pi/64+o(1)$ and $18/16+o(1)$, respectively.  Increase $X$ until
\[
 |D|<5/4,\qquad \cos D\ge1-D^2/2\ge7/32.
\]
At current stationarity, $\psi=\phi_t+j\pi/2$ has
$\sin\psi=O(L^{-1})$, so
\[
 \frac{P_j}{A_tA_\tau\omega^{2j+2}}
 =\cos\psi\cos(\psi-D)+O(L^{-1})=\cos D+O(L^{-1}).
\]
This proves the bound $1/8$ after a further uniform increase of $X$.

For absorption, the current jets give
$h_jC_j/(A_t^2\omega^{2j+2})=1+O(L^{-1})$.
The delayed/current amplitude ratio is exactly
\[
 r_\delta=\frac{A_\tau}{A_t}
  =e^{-\delta\Re(\alpha^2)/4}
       \left|\frac{\zeta(w_\tau)}{\zeta(w_t)}\right|.
\]
The second factor is at most $\zeta(5/4)^2\le25$, by the
absolutely convergent series for zeta and its reciprocal.  Uniformly
for $\theta\ge13/4$, the fact that $\ell/L\to2$ implies $d\ge1/5$
eventually (it suffices that $\ell/L\le21/10$).  Therefore
\[
 r_\delta\le25e^{-L/80+O(L^{-1})},\qquad
 \frac{Q_j}{A_t^2\omega^{2j+2}}
 =1+O(L^{-1})-\frac{r_\delta}{2}\{\cos D+O(L^{-1})\}
 \ge\frac12
\]
eventually.  For $x=-y<0$, the delay is unchanged and
$h_k(t,-y)=(-1)^kh_k(t,y)$.  The drift in the delayed residual
changes sign under reflection but remains in the same lower-order
$O(L^{-1})$ payment.  The leading margins are consequently unchanged.
\end{proof}

This theorem proves a delayed comparison within the known outer-sign
region.  It does not improve the full-sign threshold.  Using only the
historical condition $\tau L>3$ in the same proof gives, for every
fixed $\varepsilon>0$, the comparison on
$\theta\ge(11+\sqrt{217})/8+\varepsilon$, with an unevaluated
$\varepsilon$-dependent height.  No claim is made at the limiting
value, nor is that historical threshold asserted optimal for the filter.

For completeness, the exact moving derivatives in
\eqref{eq:general-moving-delay} are supplied by
\[
 \delta_t=\frac1{2B^2},\quad
 \delta_x=-\frac{t^2\ell'}{2B^2},\quad
 \delta_{xx}=-\frac{t^2\ell''}{2B^2}
                  +\frac{t^3(\ell')^2}{B^3},
\]
where $\ell'=2x/(1+x^2)$ and
$\ell''=2(1-x^2)/(1+x^2)^2$.  The pointwise theorem requires no
regular stationary branch; its transport formulas do not themselves
handle branch creation.  The finite all-earlier-time result
\eqref{eq:height112-causal} includes this rule and the half-time rule
on the two new positive-height branches.

The exact checkers
\path{numerics/check_proportional_delay_identities.py} and
\path{numerics/check_phase_controlled_delay_identities.py}
verify 20 and 21 finite algebraic statements, respectively.  The affine
transport audit performs 1200 rational point comparisons.  Their
receipts and Reviews 18--19 distinguish these checks from the analytic
phase proof and the complete-source interval calculations.  No check
evaluates the asymptotic height or proves an unbounded inner-cusp estimate.

\subsection{The remaining cross-time estimate}

One optional sufficient target is to prescribe a delay rule $0<\delta(t,x)\le t$,
a fraction $0\le\lambda<1$, and losses locally bounded through a
candidate collision, and prove on the \emph{current} set $h_{j+1}(t,x)=0$
that
\begin{align}
 h_j(t,x)C_j(t-\delta,x)&\ge0,\notag\\
 h_j(t,x)\{C_j(t,x)-(1-\lambda)C_j(t-\delta,x)\}
        &\ge-K_j(t,x)E(t,x),\qquad j=0,1.
 \label{eq:cross-time-target}
\end{align}
Lemma~\ref{lem:direct-absorption} would then supply the required
stationary loss. This is a theta-specific correlation statement: neither
backward heat, complete-source symmetry, nor the transport identities
prove its sign. The finite comparisons establish it only on their declared
cells; Theorem~\ref{thm:phase-controlled-delay} gives the separate
outer comparison. Any adaptive parameters remain prescribed independently of the
unknown energy; differentiating an adaptive filter additionally requires
all chain-rule and moving-boundary terms. The paired source formulas and
the geometric formulation identify what must be controlled, while
\eqref{eq:cross-time-target} remains an open all-height analytic estimate.
The combined signed observation can be bounded directly without
requiring both separate products to have favorable signs;
Section~\ref{pg:research-program} develops that less restrictive target.

\subsection{Polynomial cutoff losses as a quantitative target}
\label{subsec:polynomial-cutoff-transfer}

The useful role of logarithmic slack is to permit deterioration with
source resolution while retaining a collision-excluding estimate.  The
following transfer is independent of any unproved theta sign.

\begin{lemma}[Polynomial cutoff transfer]
\label{lem:polynomial-cutoff-transfer}
On a declared stationary domain, suppose independently established
$A,B\ge0$, $c>0$, a fixed finite $p\ge0$, and an integer $N_0\ge1$
give, for every integer $N\ge N_0$,
\begin{equation}
 \Lag_j\ge-A N^p E-B e^{-cN^2}.
 \label{eq:polynomial-cutoff-hypothesis}
\end{equation}
Assume $A,B,N_0,c^{-1}$ are locally bounded through candidate events.
For every fixed $M>0$ there is an independently prescribed, locally
bounded $C$ such that
\begin{equation}
 \Lag_j\ge-C E[1+\log_+(M/E)]^{p/2},
 \label{eq:polynomial-cutoff-conclusion}
\end{equation}
with value zero on the right at $E=0$.
\end{lemma}
\begin{proof}
For $E>0$ set $\ell_E=1+\log_+(M/E)$,
$d_B=\log_+(B/M)$ if $B>0$ and $d_B=0$ otherwise.  Take
\[
 N=\left\lceil\max\left\{N_0,
       \sqrt{c^{-1}\log_+(B/E)}\right\}\right\rceil,
\]
where the logarithm is defined as zero for $B=0$.  Then
$B e^{-cN^2}\le E$ and
$\log_+(B/E)\le d_B+\log_+(M/E)\le(1+d_B)\ell_E$.
Consequently $N\le D\sqrt{\ell_E}$ for
\[
 D=N_0+1+c^{-1/2}\sqrt{1+d_B},\qquad C=AD^p+1.
\]
This proves \eqref{eq:polynomial-cutoff-conclusion} and its local
uniformity.  At $E=0$, let $N\to\infty$ in the hypothesis.
\end{proof}

The same conclusion follows if a retained observation
$\Lag_j^{(N)}$ obeys the corresponding inequality and its difference
from $\Lag_j$ has a separately paid Gaussian tail.  This is an
all-$N$ family with fixed coefficient rules: proving one cutoff or
choosing coefficients from an unknown energy floor does not suffice.
The final optimization may depend on $E$ because all members of the
family have already been proved with the same prescribed constants.

For the common-carrier source, use the exact staircase cutoff
\[
 \Omega_N=\{(v,n):0\le v\le\tfrac12\log N,\quad n\le e^{2v}\}.
\]
The complete-source majorant in \eqref{eq:shared-majorant} implies
Gaussian omitted moments $O(e^{-cN^2})$ at bounded time, after reducing
$c>0$ to absorb polynomial and logarithmic factors.  The paired error
\eqref{eq:pair-tail-payment} pays both mixed sectors and the
omitted--omitted sector.  Thus a sufficient new arithmetic target on an
unresolved domain is
\begin{equation}
 \Lag_j^{(N)}\ge-A_j(t,x)N^pE-B_j(t,x)e^{-cN^2}
 \quad\text{on the full set }h_{j+1}=0,
 \quad N\ge N_0,\quad j=0,1.
 \label{eq:inner-signed-cutoff-target}
\end{equation}
All retained jet cross terms must remain, and both $E$ and stationarity
refer to the full physical jets.  The coefficients must be derived
independently of $E$ and locally bounded through candidate events.
A formulation with retained energy and a paid stationary defect is also
possible: the tail implies $E_N\le2E+O(e^{-2cN^2})$ and
$|h_{j+1}^{(N)}|=O(e^{-cN^2})$ at full stationarity.
None of these tail estimates proves the signed target.

\subsection{A cumulative signed moment and its stronger sufficient bound}
\label{subsec:cumulative-signed-moment}

An alternative exposes the signed correlation directly on the original
real source.  For $j=0,1$, put
\[
 a_j(u)=e^{tu^2}\Phi(u)u^j\cos(xu+j\pi/2),\quad
 W=h_j=\int_0^\infty a_j(u)\,du,\quad
 F_j(r)=\int_0^r a_j(u)\,du.
\]
At full stationarity, $\Lag_j=W\int_0^\infty u^2a_j(u)\,du$.
An integration by parts, with the full tail retained, gives
\begin{equation}
 \Lag_j=\int_0^R2r\{W^2-WF_j(r)\}\,dr
       +W\int_R^\infty(u^2-R^2)a_j(u)\,du.
 \label{eq:cumulative-signed-identity}
\end{equation}
Indeed integrating $u^2\,dF_j(u)$ on $[0,R]$ gives
$R^2F_j(R)-\int_0^R2rF_j(r)dr$, and
$W-F_j(R)=\int_R^\infty a_j$ supplies the displayed tail.

Suppose, as an additional unproved source estimate, that
\begin{equation}
 WF_j(r)\le W^2+A_jE\qquad(r\ge0)
 \label{eq:cumulative-overshoot-target}
\end{equation}
with a prescribed locally bounded $A_j\ge0$.  For $0\le t\le1/2$,
$0\le k\le3$, and $R\ge1$, the complete real-source envelope yields
\begin{equation}
 \int_R^\infty u^ke^{tu^2}\Phi(u)\,du
 \le2\pi e^{-4R}e^{-(\pi/2)e^{4R}}
 \le2\pi e^{-(\pi/2)e^{4R}}.
 \label{eq:cumulative-real-tail}
\end{equation}
To verify this, use $\Phi(u)<4\pi^2e^{9u-\pi e^{4u}}$ and, for
$u\ge1$, $k\log u+tu^2+9u\le u^2/2+12u\le(\pi/2)e^{4u}$.
Integrating the remaining exponential using its derivative gives the
first bound.  Since $|W|\le\sqrt E$ for both levels,
\eqref{eq:cumulative-signed-identity} implies
\[
 \Lag_j\ge-A_jR^2E-2\pi\sqrt E\,e^{-(\pi/2)e^{4R}}.
\]
At $E=0$ the stationary identity gives $\Lag_j=0$ directly. For $E>0$, set $\ell_E=1+\log_+(M/E)$,
$R=1+\tfrac14\log\ell_E$, and $c_0=(\pi/2)e^4$.
Then $R^2\le25\ell_E/16$ and
$\sqrt E e^{-c_0\ell_E}\le e^{-c_0}E/\sqrt M$, both for
$E\le M$ and for $E\ge M$.  The conditional consequence is
\[
 \Lag_j\ge-\left(\frac{25}{16}A_j+
                 \frac{2\pi e^{-c_0}}{\sqrt M}\right)
                  E[1+\log_+(M/E)].
\]
Keeping $R^2=(1+\frac14\log\ell_E)^2$ gives the sharper squared
iterated-log factor.  This is a deduction from
\eqref{eq:cumulative-overshoot-target}, not a proof of that target.
An estimate on the integrated expression in
\eqref{eq:cumulative-signed-identity} would demand less than pointwise
control for every $r$.

Positive source weights alone cannot establish the needed alignment.
For the positive two-frequency heat family
$H_*(t,x)=e^t\cos x+e^{-3T}e^{4t}\cos(2x)$, there is a double
collision at $(T,\pi)$.  At $t=T-s$, $x=\pi$,
$W=-3e^Ts+O(s^2)$, $E=9e^{2T}s^2+O(s^3)$, and
$\Lag_0=-9e^{2T}s+O(s^2)$.  For a source cutoff $1<r<2$,
$F_0(r)=-e^{T-s}$ and $WF_0(r)=3e^{2T}s+O(s^2)$, contradicting
\eqref{eq:cumulative-overshoot-target} for bounded $A_0$.
This atomic-source control is outside the theta class.  It shows why
an arithmetic relation among the actual theta contributions is needed;
evenness, rapid decay, or an absolute source-tail payment does not supply
an energy-relative signed estimate.

% Prepared for Edward Baker, 15 October 2026.
% Substantial LLM assistance: gpt-6-astra, reasoning effort ultra,
% as verified and supplied by the parent session for this turn.
% eq:wpm, eq:RS-input, eq:shared-majorant, fc:predecessor.

\section{Conditional transfers and the remaining research program}
\label{pg:research-program}

The preceding outer theorem does not cover the entire positive-time
collision domain.  This section records three further deductions that
clarify possible next steps.  Their transfer implications are proved;
the arithmetic hypotheses proposed below are not established here.
In particular, this section adds no unconditional sign region.
Write $S_j=\{h_{j+1}=0\}$ for $j=0,1$, and retain
$E=h_0^2+h_1^2$.  All stationary conditions below refer to the full
physical jets.

\subsection{A weighted remainder sufficient for the closed strip}

Retain the branches $s_\pm,w_\pm,Q_\pm$ of
\eqref{eq:alpha-Mt}--\eqref{eq:analytic-model} and the analytic model
$\mathcal P_t=Q_++Q_-$.  For a real center $x$, put
\[
 L=\log(x/(4\pi)),\qquad D_x=\{z:|z-x|\le L^{-1}\},\qquad
 \Gamma(t,x)=|M_t(s_+(x))|.
\]
For $z=u+iv\in D_x$, define
\[
 \ell=\log(u/(4\pi)),\qquad Y=e^{19\ell/100},\qquad
 N_{\rm RS}=\left\lfloor\sqrt{u/(4\pi)+t/16}\right\rfloor.
\]
The notation $N_{\rm RS}$ distinguishes this cutoff from the
staircase source cutoff used below.

\begin{proposition}[Conditional weighted-remainder transfer]
\label{pg:weighted-transfer}
Suppose that, as $x\to\infty$, for both signs one has
\begin{equation}
 \sum_{Y<n\le N_{\rm RS}}e^{t\log^2n/4}n^{-w_\pm(z)}=o(1)
 \label{pg:weighted-target}
\end{equation}
uniformly for real centers satisfying
\[
 0<t\le\tfrac12,\qquad \tfrac52\le tL\le\tfrac83,
 \qquad z\in D_x.
\]
Then, with $Q=Q_+(x)\ne0$, $\omega=L/4$, and $\phi=\arg Q$,
\begin{equation}
 \frac{1}{2|Q|}\left(h_0,\frac{h_1}{\omega},
                 \frac{h_2}{\omega^2},\frac{h_3}{\omega^3}\right)
 = (\cos\phi,-\sin\phi,-\cos\phi,\sin\phi)+o(1)
 \label{pg:weak-four-jet}
\end{equation}
uniformly on this closed strip.  Consequently $E,\Lag_0,\Lag_1$
are strictly positive there at sufficiently large height.  The
hypothesis alone does not imply an $O(L^{-1})$ error in
\eqref{pg:weak-four-jet}.
\end{proposition}

\begin{proof}
The same disc estimates as in the closed-strip argument give
$\Re w_\pm\ge11/10$ at sufficiently large height.  On $n\le Y$,
$e^q-1\le qe^q$ bounds the discrepancy from the ordinary Dirichlet
series by
\[
 \frac t4\sum_{n\ge2}(\log n)^2
 n^{-\{1/2+(81/400)t\ell-O(L^{-1})\}}=O(L^{-1}).
\]
Indeed, at the real lower boundary the exponent tends to
$1/2+(81/400)(5/2)=161/160$, and disc perturbations preserve
the lower bound $1+1/320$.  Also $t=O(L^{-1})$ on this fixed
strip.  The ordinary Dirichlet tail beyond $Y$ is
$O(e^{-19\ell/1000})$.  Adding \eqref{pg:weighted-target} therefore
gives $S_{t,N_{\rm RS}}(w_\pm)-\zeta(w_\pm)=o(1)$ on the disc.

The complete imported value errors in \eqref{eq:RS-input} remain
negligible: $e_A+e_B=O(L^2u^{-3/4})$ and
$e_{C,0}=O(u^{-2/5})$ suffice here.  Carrier magnitudes on $D_x$
are comparable to $\Gamma$.  The upper-half-disc approximation and
Schwarz reflection thus yield
\[
 H(t,z)-\mathcal P_t(z)=o(\Gamma),\qquad z\in D_x.
\]
This difference is analytic even though the intermediate cutoffs
can jump.  Cauchy's estimates, applied only to this analytic
difference, give errors $o(\Gamma L^k)$ for $0\le k\le3$ at $x$.

In $\Re w\ge11/10$, the zeta factors are bounded above and away
from zero, and their fixed-order logarithmic derivatives are bounded.
The explicit carrier differentiation consequently gives
\[
 Q_\pm^{(k)}(x)=(\pm i\omega)^kQ_\pm(x)
                         +O(\Gamma L^{k-1}),\qquad 1\le k\le3.
\]
Together with $Q_-(x)=\overline{Q_+(x)}$ and $|Q|\asymp\Gamma$,
this proves \eqref{pg:weak-four-jet}.  Substitution into the quadratic
forms gives
\[
 h_0^2+(h_1/\omega)^2=4|Q|^2(1+o(1)),\qquad
 \Lag_j=4|Q|^2\omega^{2j+2}(1+o(1)),
\]
so the strict signs follow uniformly in phase.
\end{proof}

Thus a uniform logarithmic saving for the actual weighted remainder
can suffice.  The hypothesis retains its full disc uniformity and
both endpoints; a pointwise real-axis estimate is not a substitute.
An estimate obtained only after combining the two branches may also
suffice if it directly proves the displayed analytic comparison
$H-\mathcal P_t=o(\Gamma)$.  Neither alternative is proved here.

\subsection{A full-amplitude cutoff and its conditional loss transfer}

Recall the exact staircase domain
\[
 \Omega_N=\{(v,n):0\le v\le\tfrac12\log N,\ n\le e^{2v}\},
 \qquad N\ge8,
\]
and its retained jets $h_k^{(N)}$.  Retain $a=\pi/16$, $\mathsf M_k$, $\mathsf D_0$, and $c$ from
\eqref{eq:signed-cutoff-constants}, and put $c_0=c$.
In particular $|h_k|,|h_k^{(N)}|\le e^{-ax}\mathsf M_k$ for
$0\le t\le1/2$ and $x\ge0$.  The complete staircase tail bound is
\begin{equation}
 |h_k-h_k^{(N)}|\le e^{-ax}\mathsf D_0e^{-c_0N^2},\qquad
 \mathsf D_0=\frac{\mathsf C_0}{2\beta_0},\quad c_0=\frac{\beta_0}{2}.
 \label{pg:jet-tail}
\end{equation}
It includes all omitted source indices and coordinates.

On $S_j$ define the full-amplitude observation
\begin{equation}
 J_{j,N}:=\mathcal J_j^{(N)}=-h_jh_{j+2}^{(N)},\qquad j=0,1.
 \label{pg:full-amplitude}
\end{equation}
This notation is distinct from the two-source integrand $J_j(\alpha,\beta)$.
Full stationarity gives the exact identity and estimate
\begin{equation}
 J_{j,N}-\Lag_j=h_j(h_{j+2}-h_{j+2}^{(N)}),\qquad
 |J_{j,N}-\Lag_j|\le D_j\sqrt E\,e^{-c_0N^2},
 \quad D_j=e^{-ax}\mathsf D_0.
 \label{pg:J-tail}
\end{equation}
In particular $J_{j,N}=0$ at $E=0$ for both levels.  No stationarity
condition is imposed on the retained jets.

\begin{proposition}[Conditional full-amplitude cutoff transfer]
\label{pg:J-transfer}
Let $\mathcal D$ be a declared domain in $0<t\le1/2$, $x>0$.
Suppose, on $S_j\cap\mathcal D$, that the source-prescribed inequality
\begin{equation}
 J_{j,N}\ge-A_jN^pE-B_je^{-cN^2}
 \quad\hbox{for every integer }N\ge N_0\ge8
 \label{pg:J-target}
\end{equation}
holds with a fixed finite $p\ge0$, nonnegative $A_j,B_j$, and $c>0$.
Assume $A_j,B_j,c^{-1},N_0$ are locally bounded through candidate
positive-time collisions.  Then, for each fixed $M>0$, a
source-prescribed locally bounded $K_j$ satisfies
\begin{equation}
 \Lag_j\ge-K_jE\,[1+\log_+(M/E)]^{p/2}
 \quad\hbox{on }S_j\cap\mathcal D,
 \label{pg:J-log-loss}
\end{equation}
where the right side is zero at $E=0$ and
$\log_+v=\max\{0,\log v\}$.  If this is proved for both levels
on neighborhoods containing the predecessor branches of every
candidate collision in $\mathcal D$, it excludes those collisions.
\end{proposition}

\begin{proof}
Young's inequality in \eqref{pg:J-tail} gives
\[
 D_j\sqrt E\,e^{-c_0N^2}
 \le\tfrac12E+\tfrac12D_j^2e^{-2c_0N^2}.
\]
Set $\widehat A_j=A_j+1/2$, $\widehat B_j=B_j+D_j^2/2$, and
$\widehat c=\min\{c,2c_0\}$.  Since $N^p\ge1$,
\[
 \Lag_j\ge-\widehat A_jN^pE-\widehat B_je^{-\widehat cN^2}.
\]
For $E>0$ choose
\[
 N=\left\lceil\max\left\{N_0,
          \sqrt{\widehat c^{-1}\log_+(\widehat B_j/E)}
                           \right\}\right\rceil.
\]
Put $\ell_E=1+\log_+(M/E)$ and
$d_j=\log_+(\widehat B_j/M)$.  Then the tail is at most $E$, and
\[
 N\le\Delta_j\sqrt{\ell_E},\qquad
 \Delta_j=N_0+1+\widehat c^{-1/2}\sqrt{1+d_j}.
\]
Thus \eqref{pg:J-log-loss} holds with
$K_j=\widehat A_j\Delta_j^p+1$, which is locally bounded under the
stated hypotheses.  At $E=0$, let $N\to\infty$ in the full-form
family.  All cutoff inequalities precede this optimization; none of
their coefficients is defined using a measured adverse quotient.

For the collision conclusion, the predecessor calculation in
Lemma~\ref{fc:predecessor}, applied to $H$ or $H_x$ according to
the parity of the multiplicity, gives
\[
 E\asymp(T-t)^{2r},\qquad
 -\Lag_j/E=\frac{r}{T-t}+O(1),\qquad r\ge1.
\]
The locally bounded right coefficient in \eqref{pg:J-log-loss}
permits at most a fixed power of $1+|\log(T-t)|$, contradicting
this pole.  This proves exclusion only in the covered neighborhoods.
\end{proof}

The proposition proves a transfer, not \eqref{pg:J-target} for the
theta source in the unbounded inner cusp.  A bound only for one chosen
cutoff, or with coefficients that depend on a sampled energy floor,
does not establish its hypothesis.

\subsection{Localizing the missing estimate near small energy}

The nonzero scale $\Gamma$ defined from \eqref{eq:M0} and
\eqref{eq:alpha-Mt} contains no zeta factor or physical jet denominator.
It can therefore be used even where a zeta-carrier nonvanishing claim
has not been established.  The following observation separates a
simple region from the difficult cancellation regime.

\begin{lemma}[Small-energy localization]
\label{pg:small-energy}
Let $\mathcal D\subset(0,1/2]\times(0,\infty)$ be a declared domain
on which $\Gamma$ is continuous and strictly positive, and fix $\delta>0$.
For either observation
$\mathcal O_{j,N}=J_{j,N}$ or $\mathcal O_{j,N}=\Lag_j^{(N)}$,
suppose a polynomial-cutoff lower bound
\[
 \mathcal O_{j,N}\ge-A_jN^pE-B_je^{-cN^2}
\]
has been proved for all $N\ge N_0\ge8$ on
\begin{equation}
 S_j\cap\mathcal D\cap\{E\le\delta\Gamma^2\}.
 \label{pg:small-energy-domain}
\end{equation}
Here $p\ge0$ is fixed and finite, and the prescribed nonnegative
coefficients are defined on $\mathcal D$ with the local boundedness
conditions of Proposition~\ref{pg:J-transfer}.  The same type of
bound then holds on all of $S_j\cap\mathcal D$ after replacing $A_j$ by
\begin{equation}
 \widetilde A_j=
 \max\left\{A_j,
   \frac{e^{-2ax}\mathsf M_j\mathsf M_{j+2}}
        {\delta\Gamma^2}\right\}.
 \label{pg:extended-coefficient}
\end{equation}
This new coefficient is locally bounded and independent of $N$.
\end{lemma}

\begin{proof}
For $J_{j,N}$ the separate complete and retained jet envelopes give
$|J_{j,N}|\le U_j$, where
$U_j=e^{-2ax}\mathsf M_j\mathsf M_{j+2}$.
For the raw cutoff the complete paired identity gives the same
envelope $|\Lag_j^{(N)}|\le U_j$.  On $E>\delta\Gamma^2$,
\[
 \mathcal O_{j,N}\ge-U_j
       \ge-\frac{U_j}{\delta\Gamma^2}E.
\]
Use $N^p\ge1$ and combine this with the assumed bound on
\eqref{pg:small-energy-domain}.  Local boundedness follows from
continuity and positivity of $\Gamma$ at every finite point.
\end{proof}

For any hypothetical finite collision with $\Gamma>0$ whose sufficiently
nearby predecessor points stay in $\mathcal D$, continuity and $E=0$
at the collision place those points in \eqref{pg:small-energy-domain}.
Thus, wherever the declared domain covers these predecessors, the
restricted small-energy region contains the essential collision obstruction.  The
localization assumes no positive energy floor there and does not prove
the missing lower bound.  Growth of the coefficient as height tends to
infinity or as $t\downarrow0$ is compatible with the local criterion.

\subsection{Arithmetic directions and limits of the present coverage}

The main proposed cusp program is to prove
\eqref{pg:J-target}, or its paid raw-cutoff analogue, using relations
among the actual theta amplitudes.  Exact Poisson or modular relations
are candidates only if they constrain adverse interactions at the
physical energy scale.  Both Hermitian and bilinear sectors, independent
source coordinates, and all staircase boundary terms must remain.
The paired kernel has no universal sign for arbitrary amplitudes;
positive source weights and the heat equation alone are insufficient.

The cumulative identity \eqref{eq:cumulative-signed-identity}
provides a related diagnostic. Grouping consecutive oscillatory lobes
before estimating the integrated signed overshoot may demand less than
\eqref{eq:cumulative-overshoot-target}. The complete tail must still be
paid at the physical energy scale; a merely absolute error proportional
to $|h_j|$ is insufficient near a collision.

The central arithmetic gap is not removed by the published systematic
exponent table \cite[Table 1]{taotrudgianyang25}. At block scale
$\beta=2/5$, its bound is
\[
 \frac{18}{199}+\frac{521}{796}\frac25=\frac{701}{1990},
 \qquad \frac{701}{1990}-\frac7{20}=\frac9{3980}>0.
\]
Here $7/20$ is the exponent allowed by the $\theta=5/2$ target before
its additional saving. This is a limitation of the quoted bounds, not
an obstruction to a logarithmic-phase-specific or signed estimate.

The outer program seeks \eqref{pg:weighted-target} through exact
partitions, transformed logarithmic-phase sums, or cancellation across
neighboring scales.  A small diagnostic computation should specify
its observation and coefficient rule before validation, retain complete
tails, and increase the cutoff and precision when the physical carrier
scale requires it.  Along a sampling curve $t=\theta/L(x)$ the derivative
used to isolate a root of the physical stationary observation is
\[
 \frac d{dx}h_{j+1}(\theta/L(x),x)
   =h_{j+2}+\frac{\theta}{xL(x)^2}h_{j+3}.
\]
An undecided enclosure is not a sign violation; favorable finite cells
do not establish an unbounded-domain estimate.

Finally, the nominal zeta argument has real part $1/2+\theta/4$.
The elementary nonvanishing input for the present carrier works with
a fixed margin when $\theta>2$.  At and below two that justification
no longer applies.  A different argument may cross the boundary, but
dividing by a zeta factor whose nonvanishing has not been proved would
import the missing information.  This is a limitation of the present
method, not an impossibility theorem for all carrier methods.

Even a sequence of outer thresholds $\theta_n\downarrow0$ does not
settle the global problem without controlling its associated starting
heights $X_n$.  At a fixed hypothetical collision $(T,x_*)$, choosing
$n$ so that $\theta_n<T\log(x_*/(4\pi))$ may also give $X_n>x_*$.
A finite-height or other global argument must close that quantifier
gap.  Effective bridging should therefore target an evaluated reachable
analytic cutoff; repeated bounded extensions alone do not address the
unbounded small-time cusp.

\appendix
\section{Complete source and integration error bounds}\label{app:payments}

The following bounds make the numerical and shared-source constructions
independent of an unpaid finite theta truncation. Set
\[
 S_k(q,d)=\sum_{r=0}^k\binom kr\frac{d^{k-r}r!}{q^{r+1}}
         =\int_0^\infty(v+d)^ke^{-qv}\dd v\qquad(q>0).
\]
Let $0\le t\le\tau\le1/2$, $0<a<\pi/8$, and $x\ge0$.
For an index cutoff $N$, put $m=N+1$ and
\[
 c_a=\cos4a,\quad \alpha=\pi m^2,\quad \beta=c_a\alpha,
 \quad q_N=4\beta-9,\quad
 \rho_N=((m+1)/m)^4e^{-\pi c_a(2m+1)}.
\]
If $\rho_N<1$, $q_N>0$, and $8\beta\ge\tau$, the complete omitted
index contribution in \eqref{eq:complete-shift} is at most
\begin{equation}
 E_{N,k}=e^{-ax}\frac{(2\alpha^2+3\alpha)e^{-\beta}}{1-\rho_N}
                         S_k(q_N,a).
 \label{eq:tail-index}
\end{equation}
For a cutoff $U\ge0$, put
\[
 \beta_U=\pi c_ae^{4U},\quad \rho_U=16e^{-3\beta_U},\quad
 q_U=4\beta_U-9-2\tau U,\quad B_U=\tau U^2+9U-\beta_U.
\]
If $\rho_U<1$, $q_U>0$, and $8\beta_U\ge\tau$, the entire source
on $u\ge U$ contributes at most
\begin{equation}
 E_{U,k}=e^{-ax}\frac{(2\pi^2+3\pi)e^{B_U}}{1-\rho_U}
                         S_k(q_U,U+a).
 \label{eq:tail-integration}
\end{equation}
Thus $E_{N,k}+E_{U,k}$ pays the omitted domain of an $N$-term integral
on $[0,U]$. Validated retained quadrature and cell variation are
additional errors.

To prove the bounds, use
\[
 |\phi_n(u+ia)|\le(2\alpha_n^2+3\alpha_n)
 e^{9u-c_a\alpha_ne^{4u}},\quad
 |e^{t(u+ia)^2}|\le e^{\tau u^2},\quad |u+ia|^k\le(u+a)^k.
\]
For $n\ge m$, successive coefficient-majorant ratios are bounded by
$\rho_N$ uniformly for $u\ge0$. The inequality
$e^{4u}\ge1+4u+8u^2$ gives
\[
 \tau u^2+9u-\beta e^{4u}
 \le-\beta-q_Nu-(8\beta-\tau)u^2\le-\beta-q_Nu.
\]
Summing the full geometric index bound and integrating proves
\eqref{eq:tail-index}. For $u\ge U$, the exponent
$\tau u^2+9u-\pi c_ae^{4u}$ is concave and has slope $-q_U$ at
$U$. Its tangent bound and the geometric ratio $\rho_U$ give
\eqref{eq:tail-integration}. Adding the bounds overcounts their
intersection, which is harmless.

For the centered residual $C_j$, replace any factor $S_j(q,d)$ in
these payments by $S_{j+2}(q,d)+b^2S_j(q,d)$, since
$|u^2-b^2|\le(u+a)^2+b^2$. For the shared carrier
\eqref{eq:shared-carrier}, set
\[
 \widehat\beta_V=\pi\cos(4a)e^{4V},\quad
 \widehat q_V=4\widehat\beta_V-10-2\tau V,\quad
 \widehat B_V=\tau V^2+10V-\widehat\beta_V.
\]
If $\widehat q_V>0$ and $8\widehat\beta_V\ge\tau$, the majorant
\eqref{eq:shared-majorant} and the same tangent argument give the full
tail payment
\begin{equation}
 e^{-ax}2(2\pi^2+3\pi)e^{\widehat B_V}S_j(\widehat q_V,V+a).
 \label{eq:shared-tail}
\end{equation}
The same degree-two replacement pays the shared expression for $C_j$,
because $0\le v-d_n\le v$.

For completeness, these absolute errors can be made small at the gamma
exponential scale at arbitrary height. For $x\ge41$, $p>0$, and
$\ell_x=\log x$, take
\[
 a_x=\pi/8-1/x,\quad
 m=\left\lceil\sqrt{\frac{x}{2\pi}(p\ell_x+32)}\right\rceil,
 \quad N=m-1,\quad U=\frac{\ell_x+2\log(\ell_x+8)}4.
\]
All preceding payment hypotheses hold, and for fixed $k$,
\begin{equation}
 e^{\pi x/8}E_{N,k}=O_{p,k}(x^{5/2-p}(\log x)^{3/2}),\qquad
 e^{\pi x/8}E_{U,k}=O_{k,A}(x^{-A})\quad(A>0).
 \label{eq:adaptive-payments}
\end{equation}
Indeed $c_a=\sin(4/x)\ge2/x$, $\beta\ge p\ell_x+32$, and
$m^2\ge2x/\pi$. Hence
$-\log\rho_N\ge4\pi m/x-4/m\ge2\pi m/x$ and
$(1-\rho_N)^{-1}\le1+x/(2\pi m)$.
Together with $e^{-a_xx}=e\,e^{-\pi x/8}$, substitution gives the
first bound. For the second,
$\beta_U\ge2\pi(\ell_x+8)^2$, $U\le\ell_x/2+2$, and
$B_U\le-47\ell_x^2/8-181\ell_x/2-364$.
These are absolute error estimates. They neither establish a norm gap
nor replace a signed estimate of the retained source. On a spatial cell,
freeze the contour and cutoffs before differentiating or transporting
physical jets.


% Prepared for Edward Baker; gpt-6-astra, effort ultra; 14 October 2026.
\section{An explicit height at the condition \texorpdfstring{$tL\ge5$}{tL >= 5}}
\label{app:explicit-height}

The exponent-pair argument in Section~\ref{sec:outer-proof} gives the
stronger asymptotic condition $tL\ge61/24$, with an unevaluated
height. The following elementary absolute-sum calculation supplies an
explicit height for $tL\ge5$. These are distinct quantitative statements.

\begin{proposition}[Explicit outer height]\label{prop:explicit-height}
Let $L=\log(x/(4\pi))$, $\omega=L/4$, and let $Q=Q_+(x)$ be the
amplitude in \eqref{eq:analytic-model}. If
\[
 x\ge X_{\mathrm{exp}}:=4\pi e^{10^9},\qquad
 0<t\le\tfrac12,\qquad tL\ge5,
\]
then every coordinate error in \eqref{eq:main-four-jet} is at most
\begin{equation}
 \frac{108000000}{L}\le0.108<\frac18.
 \label{eq:effective-four-jet-budget}
\end{equation}
In particular, for $P_k=h_k/(2|Q|\omega^k)$,
\[
 P_0^2+P_1^2\ge\tfrac12,\qquad
 P_1^2-P_0P_2\ge\tfrac{31}{64},\qquad
 P_2^2-P_1P_3\ge\tfrac{31}{64}.
\]
The same conclusions hold at negative heights after reflection.
\end{proposition}

\subsection{Explicit weights on the complex disc}
All bounds until the last sign step hold for $L\ge1000$,
$0<t\le1/2$, and $tL\ge5$. On $D_x$ from
\eqref{eq:shrinking-disc}, write $z=u+iv$ and $\theta_u=tL_u$. Then
\begin{equation}
 u\ge x/2\ge e^L,\qquad |L_u-L|\le\frac2{xL}<1,
 \qquad\theta_u\ge5-\frac1{xL}\ge\frac{19}{4}.
 \label{eq:effective-disc-geometry}
\end{equation}
For $s=s_\pm(z)$, one has $0\le\Re s\le1$ and
$|s|,|s-1|\ge u/2$. Formula~\eqref{eq:alpha-Mt} therefore yields
\begin{equation}
 \Re\alpha(s)\ge L_u/2-3/u,\qquad
 |\alpha(s)|\le L_u/2+1\le L,
 \qquad |\alpha'(s)|\le1/u+6/u^2\le2/u.
 \label{eq:effective-alpha-disc}
\end{equation}
For the magnitude estimate, use
$|s|\le(u/2)\sqrt{1+4/u^2}$, $|\arg s|\le\pi/2$, and the bound
$3/u$ for the rational terms. It follows that
\begin{equation}
 \sigma_\pm:=\Re w_\pm
 \ge\frac12-\frac1{2L}+\frac{\theta_u}{4}-\frac{3t}{2u}
 \ge\frac32.
 \label{eq:effective-zeta-half-plane}
\end{equation}
The cutoff satisfies
\[
 \tfrac12e^{L_u/2}\le N,\qquad\log N\le L_u/2+1/u.
\]
Indeed the upper correction is at most $\pi t/(8u)$, and flooring
reduces a number at least two by at most a factor of two.
Consequently, for $n\le N$,
\begin{equation}
 e^{t\log^2n/4}n^{-\sigma_\pm}\le n^{-p},\qquad
 p=\frac12+\frac{\theta_u}{8}-\delta,\qquad
 \delta=\frac1{2L}+\frac{7t}{4u}\le\frac1{32}.
 \label{eq:effective-weight}
\end{equation}
Writing $s_0=\theta_u-19/4\ge0$, we obtain
$p\ge17/16+s_0/8$. On the upper half-disc, the same bound holds for
the first branch after multiplication by $n^v$: its leading real part
changes from $(1+v)/2$ to $(1-v)/2$, already included in the estimate.

For $\eta>0$ and an integer $m\ge0$, comparison on $[n-1,n]$ gives
\begin{equation}
 \sum_{n=2}^{\infty}(\log n)^m n^{-1-\eta}
 \le\int_1^\infty(\log y+1)^m y^{-1-\eta}\dd y.
 \label{eq:effective-log-moments}
\end{equation}
At $m=2$, $\eta=1/16$, the right side is
$2/\eta^3+2/\eta^2+1/\eta=8720<9000$. Hence
\begin{align}
 |S_{t,N}(w_\pm)-\zeta(w_\pm)|
 &\le\frac{9000}{4L_u}(19/4+s_0)2^{-s_0/8}+2N^{-1/2}
 \notag\\
 &\le\frac{29250}{L_u}+3e^{-L_u/4}
 \le\frac{31000}{L}.
 \label{eq:effective-sum-discrepancy}
\end{align}
Here $s_0e^{-s_0\log2/8}\le8/(e\log2)<8$, so the time-dependent
numerator is less than $13$. Moreover,
$29250/(L-1)<30000/L$ for $L\ge1000$, while
$e^{(L-1)/4}\ge(L-1)^2/32\ge3L$. Thus the displayed bound includes
both the weighted discrepancy and the complete ordinary zeta tail,
uniformly even as $tL$ grows.

\subsection{All imported approximation errors}
On the upper half-disc, \eqref{eq:RS-ratios} gives $|\gamma|\le2$ and
\[
 |\kappa|\log N
 \le\frac{tv}{2(u-6)}(L_u/2+1/u)<\tfrac12,
 \qquad N^{|\kappa|}\le2.
\]
Since $\zeta(17/16)\le17$, the full positive weight sum in
\eqref{eq:RS-AB} is at most $5\cdot17=85$. Denote that equation's
exponent by
\[
 q_n=\frac{(t^2/16)\log^2(u/(4\pi n^2))+0.626}{u-6.66}.
\]
The estimates $-2/u\le\log(u/(4\pi n^2))\le L_u$,
$(L+1)^2/64+1\le L^2/32$, and $u-6.66\ge u/2$ imply
$0\le q_n\le L^2/(16u)<1$. Using $e^q-1\le2q$ on $[0,1]$ gives
\begin{equation}
 e_A+e_B\le\frac{85L^2}{8u}\le\frac{11L^2}{u}\le\frac1L.
 \label{eq:effective-AB}
\end{equation}
For the final inequality, $u\ge e^L\ge L^4/24\ge11L^3$ suffices.

The positive correction exponent in \eqref{eq:RS-C} is at most
\[
 \frac{1.24(3^v+3^{-v})}{N-0.125}
 +\frac{3|L_u+i\pi/2|+10.44}{u-12}
 \le20e^{-L_u/2}+7L/u\le1.
\]
To verify the last inequality, use
$20e^{-(L-1)/2}\le160/(L-1)^2<1/2$ and
$7Le^{-L}\le168/L^3<1/2$. Its negative main exponent therefore gives
\begin{equation}
 e_{C,0}\le e\exp\{-(1/4+19/64)L_u\}
       \le3e^{-L_u/2}\le\frac1L,
 \qquad e_A+e_B+e_{C,0}\le\frac2L.
 \label{eq:effective-C-and-total}
\end{equation}
All three errors in the imported theorem are included; the constant
$0.626$ in $q_n$ is independent of $t$.

\subsection{Amplitude comparison and all four derivatives}
Let $\Gamma=|M_t(s_+(x))|=|M_t(s_-(x))|$. On $D_x$,
\[
 \left|\frac{d}{dz}\log M_t(s_\pm(z))\right|
 \le\frac{|\alpha|}{2}(1+t|\alpha'|/2)\le L.
\]
Integration along a segment of length at most $1/L$ yields
$|M_t(s_\pm(z))|\le e\Gamma<3\Gamma$. Also
\eqref{eq:effective-zeta-half-plane} and the zeta and reciprocal-zeta
Dirichlet series give
$|\zeta(w)|,|1/\zeta(w)|\le\zeta(3/2)\le3$, hence
$\Gamma\le3|Q|$. Combining \eqref{eq:effective-sum-discrepancy} and
\eqref{eq:effective-C-and-total} with \eqref{eq:RS-input} gives
\begin{align}
 |H(t,z)-\mathcal P_t(z)|
 &\le 2(3\Gamma)\frac{31000}{L}+(3\Gamma)\frac2L
   =\frac{186006\Gamma}{L}
 \notag\\
 &\le\frac{558018|Q|}{L}\le\frac{560000|Q|}{L}.
 \label{eq:effective-disc-error}
\end{align}
Schwarz symmetry extends this upper-half-disc estimate to the full disc.
Cauchy's inequality now gives
\begin{equation}
 \frac{|\partial_x^k(H(t,\cdot)-\mathcal P_t)(x)|}
      {2|Q|\omega^k}
 \le\frac{280000\,k!\,4^k}{L}
 \le\frac{107520000}{L},\qquad 0\le k\le3.
 \label{eq:effective-cauchy-error}
\end{equation}

For the analytic carrier, at $s=(1+ix)/2$ direct differentiation of
\eqref{eq:alpha-Mt} gives
\begin{equation}
 |\alpha-(L/2+i\pi/4)|\le4/x,\quad
 |\alpha'|\le2/x,\quad |\alpha''|\le3/x^2,\quad
 |\alpha'''|\le9/x^3.
 \label{eq:effective-alpha-center}
\end{equation}
The first estimate follows from $s=(ix/2)(1-i/x)$ and
$|\operatorname{Log}(1-i/x)|\le2/x$. For the others one uses
\[
 \alpha'=-\frac1{2s^2}-\frac1{(s-1)^2}+\frac1{2s},\qquad
 \alpha''=\frac1{s^3}+\frac2{(s-1)^3}-\frac1{2s^2},\qquad
 \alpha'''=-\frac3{s^4}-\frac6{(s-1)^4}+\frac1{s^3}.
\]
Let $J(x)=\log M_t(s(x))$. Then
\begin{align*}
 J'&=(i/2)\alpha(1+t\alpha'/2),\\
 J''&=-\tfrac14\{\alpha'+(t/2)[(\alpha')^2+\alpha\alpha'']\},\\
 J'''&=-\tfrac i8\{\alpha''+(t/2)[3\alpha'\alpha''+\alpha\alpha''']\}.
\end{align*}
Equations~\eqref{eq:effective-alpha-disc} and
\eqref{eq:effective-alpha-center} imply
$|J'-i\omega|\le1$, $|J''|,|J'''|\le1$ and
$|w'|,|w''|,|w'''|\le1$. A majorant for the first bound is
$\pi/8+2/x+L/(4x)<1$; the remaining bounds follow by inserting the
stated powers of $x^{-1}$ and $x\ge e^L$.

The prime-power series and \eqref{eq:effective-log-moments} at
$\eta=1/2$ give
\[
 |(\log\zeta)'(w)|\le6,\qquad
 |(\log\zeta)''(w)|\le26,\qquad
 |(\log\zeta)'''(w)|\le158.
\]
Thus for a local logarithm $g=\log Q_+$, the chain rule yields
\[
 |g'-i\omega|\le7,\qquad |g''|\le1+26+6=33,
 \qquad |g'''|\le1+158+3\cdot26+6=243.
\]
The resulting carrier errors are
\begin{equation}
\begin{array}{c|c}
 k&\displaystyle\frac{|Q_+^{(k)}/Q-(i\omega)^k|}{\omega^k}\\[2pt]
 \hline
 1&28/L\\
 2&56/L+1312/L^2\\
 3&84/L+3936/L^2+81856/L^3
\end{array}
\label{eq:effective-carrier-table}
\end{equation}
Every entry in the right-hand column is at most $100/L$ for $L\ge1000$.
At order three, for instance, expanding
$(i\omega+d)^3+3(i\omega+d)g''+g'''$ with $|d|\le7$ bounds the
unnormalized discrepancy by $21\omega^2+246\omega+1279$.
The conjugate branch contributes the same normalized bound; the factor
two is already included in $2|Q|\omega^k$. The carrier error at order
zero is exactly zero.

Adding \eqref{eq:effective-cauchy-error} and
\eqref{eq:effective-carrier-table}, every coordinate error is at most
$107520100/L<108000000/L$. For $L\ge10^9$ this is less than $1/8$,
and \eqref{eq:explicit-perturbation} proves
Proposition~\ref{prop:explicit-height}.

The value $X_{\mathrm{exp}}$ is an explicit consequence of the primary
approximation and elementary inequalities. Its exponential need not be
evaluated numerically. The accompanying exact checker verifies 52 rational
budgets and sufficient polynomial inequalities for every $L\ge1000$ by
shifting $L=1000+y$ and checking nonnegative coefficients. It audits the
arithmetic in this proof, rather than replacing its analytic arguments.
The bound supplies no practical finite-strip coverage between 112 and
$X_{\mathrm{exp}}$, and it does not evaluate the height in the stronger
closed-threshold theorem.

\section{Reproducibility and scope of the checks}\label{app:replay}
\begingroup\raggedright

All paths in this appendix are relative to
\path{papers/investigations-2/active/21_quaternion_theta_geometry/}
in the accompanying repository. The analytic proofs above are distinct
from the finite interval decisions. The latter use complete error
bounds, outward rounding, and closed-cell coverage; hashes identify
the evaluated sources and records, but do not substitute for their
mathematical content.

The bounded results can be regenerated using the drivers
\path{numerics/certify_compensated_stationary_time_extension.py},
\path{numerics/certify_correlated_source_height_extension.py}, and
\path{numerics/certify_phase_centered_residual.py}.
The correlated-height guide
\path{numerics/CORRELATED_SOURCE_HEIGHT_README_20261014.md}
gives the parameters for its three bands. The full-time certificate is
\path{numerics/COMPENSATED_STATIONARY_TIME_EXTENSION_20261014.json};
the final centered-band record is
\path{numerics/THETA_CENTERED_RESIDUAL_100_104_20261014.json}.
The independent internal replays are documented in Reviews 6, 9, and 11
in \path{reviews/}. They reproduced the numerical records after removal
of runtime metadata.

For compact exact audits, run the following from \path{numerics/}:
\begin{quote}\small\ttfamily
python3 audit\_correlated\_source\_height\_records.py\par
python3 audit\_theta\_centered\_residual\_record.py\par
python3 check\_octonion\_jet\_identities.py\par
python3 check\_theta\_centered\_residual\_identities.py
\end{quote}
The centered-band auditor additionally requires
\texttt{--record THETA\_CENTERED\_RESIDUAL\_100\_104\_20261014.json}.
Each requires the option \texttt{--output} followed by a desired JSON
path. The first two inspect the retained records and bound source
versions; the latter two check 248 and nine exact algebraic identities,
respectively. The octonionic identities, centered residual, and outer
analytic argument have separate internal assessments in Reviews 10--12.

The diagnostic
\path{numerics/check_outer_cusp_and_phase_obstructions.py}, also run
with \texttt{--output}, requires python-flint and imports
\path{numerics/check_shifted_contour_payments.py}. Its retained record
\path{numerics/OUTER_CUSP_AND_PHASE_OBSTRUCTION_CHECKS_20261014.json}
contains the pointwise counterexamples in Section~\ref{sec:remaining},
twelve staircase boundary checks, exact four-jet perturbation algebra,
and selected complete-source point enclosures. It does not certify a
new rectangle or prove the uniform asymptotic theorem.

The subsequent paired algebra, explicit outer constants, and signed-split
controls are audited by
\path{numerics/check_signed_carrier_pair_identities.py},
\path{numerics/check_explicit_outer_constants.py}, and
\path{numerics/check_signed_curvature_split.py}.
They require only the Python standard library and an \texttt{--output}
path, and check 34, 52, and 14 exact statements respectively. Their
records are
\path{numerics/SIGNED_CARRIER_PAIR_IDENTITY_CHECKS_20261014.json},
\path{numerics/EXPLICIT_OUTER_CONSTANT_CHECKS_20261014.json}, and
\path{numerics/SIGNED_CURVATURE_SPLIT_IDENTITY_CHECKS_20261014.json}.
These checks support algebra and explicit inequalities; the infinite
source and asymptotic arguments require the proofs in the text.

The full-residual branch comparison is recorded in
\path{numerics/STATIONARY_FAVORABLE_CURVATURE_PILOT_20261014.json}.
The proper source-filter and Gaussian pilots are produced by
\path{numerics/certify_predetermined_source_filter_pilot.py} and
\path{numerics/check_gaussian_source_filter_pilot.py}.
Each accepts \texttt{--record}, \texttt{--original-driver},
\texttt{--contour-source}, and \texttt{--output}; the inputs are the
centered-band record, its driver, and
\path{numerics/check_shifted_contour_payments.py}.
The proper-filter run uses python-flint 0.9.0 at 192 bits and retains
all 64 classifications, including 32 failed absorption conditions.
The Gaussian extraction uses exact rational arithmetic and retains all
14 cases. Their records are
\path{numerics/PREDETERMINED_SOURCE_FILTER_PILOT_20261014.json} and
\path{numerics/GAUSSIAN_SOURCE_FILTER_PILOT_20261014.json}.
Reviews 14--16 document the analytic audits and independent replays,
including the eight full-residual branch slabs. Review 17 audits the
earlier manuscript integration; its fresh exact checks are distinct
from the preceding full interval replays.

The height-112 extension and its all-earlier-time comparisons use
\path{numerics/certify_fixed_cusp_band_extension.py} and the two records
\path{numerics/FIXED_CUSP_BAND_104_108_20261014.json} and
\path{numerics/FIXED_CUSP_BAND_108_112_20261014.json}.
The exact auditor \path{numerics/audit_fixed_cusp_band_extension.py}
accepts both paths through repeated \texttt{--record} options and an
\texttt{--output} path. Its 1321 checks retain all 64 closed cells,
128 stationary observations, and 16 genuine branch/time cases.
The source quadratures were independently replayed in Review 19; the
present integration reruns the exact auditor, not those quadratures.

The newly integrated arithmetic and collision material is audited by
\path{numerics/check_proportional_delay_identities.py} (20 checks),
\path{numerics/check_phase_controlled_delay_identities.py} (21),
\path{numerics/check_collision_behavior.py} (1072),
\path{numerics/check_hybrid_cutoff_lemma.py} (84), and
\path{numerics/check_remaining_strip_lemma.py} (97).
The proportional checker takes a positional output path; the phase and
strip checkers accept \texttt{--output}. The collision driver writes its
named receipt beside its source; the integration replay calls its pure
check routines without replacing that historical receipt.
Together with the finite-band audit, the integration executes 2615 exact
checks. The new receipt is
\path{numerics/MANUSCRIPT_PROGRAM_INTEGRATION_CHECKS_20261015.json}.
These checks verify finite algebra, rational margins, and recorded
interval decisions. They do not prove the imported arithmetic theorems,
the unproved middle-block or weighted-remainder hypotheses, or an
all-height theta sign estimate. Review 24 records the integrated proof
and layout audit, including the conditional deductions from Note 33.

The main remaining obligation is the one-sided inner-region bound
\eqref{eq:final-target}, with all cross terms retained and losses locally
bounded through candidate positive-time collisions. The complete-source
Gaussian identities make a specific cross-time correlation available for
that task. Improving the explicit height and connecting the finite cover
remain separate quantitative work. Neither a fixed-time asymptotic nor
a ratio divided by an unknown physical energy supplies the missing bound.

\par\endgroup

\section*{Acknowledgment}
This manuscript was drafted for Edward Baker with substantial assistance
from OpenAI language models. The present integration used Codex with the
recorded model \texttt{gpt-6-astra} and reasoning effort \texttt{ultra};
earlier research and numerical certificates also used
\texttt{gpt-6.1-sol}. LLM assistance included derivations, exposition,
code, and internal checks. The accompanying reviews distinguish
analytic arguments, validated interval calculations, exact algebra
checks, and their remaining limitations. These internal checks are not
external specialist refereeing or proof-assistant formalization.

\begin{thebibliography}{99}
\raggedright
\bibitem{debruijn}
N. G. de Bruijn,
\emph{The roots of trigonometric integrals},
Duke Mathematical Journal \textbf{17} (1950), 197--226.
\url{https://doi.org/10.1215/S0012-7094-50-01720-0}.

\bibitem{newman}
C. M. Newman,
\emph{Fourier transforms with only real zeros},
Proceedings of the American Mathematical Society \textbf{61} (1976),
245--251.
\url{https://doi.org/10.1090/S0002-9939-1976-0434982-5}.

\bibitem{polymath}
D. H. J. Polymath,
\emph{Effective approximation of heat flow evolution of the Riemann
$\xi$ function, and a new upper bound for the de Bruijn--Newman constant},
arXiv:1904.12438v2 (2019).
\url{https://arxiv.org/abs/1904.12438}.

\bibitem{robert}
O. Robert,
\emph{On van der Corput's $k$-th derivative test for exponential sums},
Indagationes Mathematicae \textbf{27} (2016), 559--589.
\url{https://doi.org/10.1016/j.indag.2015.11.009}.
Author manuscript:
\url{https://perso.univ-st-etienne.fr/rool6510/robert-2015-indag.pdf}.

\bibitem{sudbery}
A. Sudbery,
\emph{Quaternionic analysis},
Mathematical Proceedings of the Cambridge Philosophical Society
\textbf{85} (1979), 199--225.
\url{https://doi.org/10.1017/S0305004100055638}.
% Primary references checked 14 October 2026.
% Existing manuscript already has the sudbery key.
\bibitem{baez}
J. C. Baez,
\emph{The octonions},
Bulletin of the American Mathematical Society \textbf{39} (2002), 145--205.
\url{https://arxiv.org/abs/math/0105155}.
The Cayley--Dickson convention is in Section~2.2;
\url{https://math.ucr.edu/home/baez/octonions/node5.html}.

\bibitem{kauhanenorelma}
J. Kauhanen and H. Orelma,
\emph{Cauchy--Riemann operators in octonionic analysis},
arXiv:1701.08698 (2017).
\url{https://arxiv.org/abs/1701.08698}.

\bibitem{bourgain14}
J. Bourgain,
\emph{Decoupling, exponential sums and the Riemann zeta function},
arXiv:1408.5794v2, especially Theorem 6 and its subinterval discussion.
\url{https://arxiv.org/abs/1408.5794v2}.

\bibitem{trudgianyang23}
T. S. Trudgian and A. Yang,
\emph{Toward optimal exponent pairs},
arXiv:2306.05599v3, especially Section 1.1, Lemma 1.1, and equation (1.10).
\url{https://arxiv.org/abs/2306.05599v3}.

\bibitem{ruiz}
J. M. Ruiz,
\emph{Basic properties of real analytic and semianalytic germs},
Publications mathematiques et informatique de Rennes (1986), no. 4,
29--44, especially (7.2), p. 40.
\url{https://www.numdam.org/item/PSMIR_1986___4_29_0.pdf}.

\bibitem{taotrudgianyang25}
T. Tao, T. Trudgian, and A. Yang,
\emph{New exponent pairs, zero density estimates, and zero additive
energy estimates: a systematic approach}, arXiv:2501.16779v1 (2025),
Table 1. \url{https://arxiv.org/abs/2501.16779v1}.

\end{thebibliography}
\end{document}
